%%%%%%%%%%%%%%%%%%%%%%%%%%%%%%%%%%%%%%%%%%%%%%%%%%%%%%%%%%%%%%%%%%%%%%
% Lớp 11 | Các số đặc trưng đo xu thế trung tâm cho mẫu số liệu | Số trung bình và mốt của mẫu số liệu ghép nhóm
% Dạng 1D5?1-3: Số trung bình
% Số câu: 129  (sinh tự động bởi tach_cau.py; từ nay thầy sửa trực tiếp file này)
%%%%%%%%%%%%%%%%%%%%%%%%%%%%%%%%%%%%%%%%%%%%%%%%%%%%%%%%%%%%%%%%%%%%%%

%%%%%%%%%%%%%%%% TRẮC NGHIỆM 4 PHƯƠNG ÁN %%%%%%%%%%%%%%%%

% nguồn: Ok_Toan11_TracNghiem_Lan2
\begin{ex}%[1D5N1-3]%[ma:3bc5d3a4]
	Tìm cân nặng trung bình (đơn vị: kg) của học sinh lớp 11D cho  trong bảng sau
	\begin{center}
		\begin{tabular}{|l|c|c|c|c|c|c|c|}
			\hline Cân nặng& {$[40{,}5;45{,}5)$} & {$ [45{,}5;50{,}5)$} & {$[50{,}5;55{,}5) $} &  {$[55{,}5;60{,}5)$} & {$[60{,}5;65{,}5)$} & {$[65{,}5;70{,}5)$}\\
			\hline Số học sinh & $10$ &  $7$ & $16$ & $4$& $2$& $3$ \\
			\hline 
		\end{tabular}
	\end{center}
	 \choice
		{ \True $51{,}81$} 
		 { $57{,}6$}
		 { $58{,}7$}
		{  $52{,}6$}
	 \loigiai{
		Trong mỗi khoảng cân nặng, giá trị đại diện là trung bình cộng của giá trị hai đầu mút nên ta có  bảng sau (kết quả làm tròn đến hàng phần trăm).
		\begin{center}
			\begin{tabular}{|l|c|c|c|c|c|c|c|}
				 \hline Cân nặng& {$43$} & {$ 48$} & {$53$} & {$58$} & {$63$} & {$68$}\\
				\hline Số học sinh  & $10$ & $7$ & $16$ & $4$& $2$& $3$ \\
				 \hline
			\end{tabular}
		\end{center}
		Tổng số học sinh là $n=42$. Cân nặng trung bình của học sinh lớp  11D là
	$$\overline{x}=\dfrac{10\cdot 43+7\cdot 48+16\cdot 53+4\cdot  58+2\cdot 63+3\cdot 68}{42}\approx 51{,}81\text{ kg}$$}
\end{ex}

% nguồn: Ok_Toan11_TracNghiem_Lan2
\begin{ex}%[1D5N1-3]%[ma:b3bf0a81]
	Cân nặng (đơn vị: kg) của  $28$ học sinh lớp $11$ được thống kê trong bảng tần số ghép nhóm sau
	\begin{center}
		\begin{tabular}{|l|c|c|c|c|c|}
			\hline Cân  nặng & {$[45;49)$} & {$ [49;53)$} & {$[53;57) $} & {$[57;61)$} & {$[61;65)$} \\
			\hline Số học sinh & $4$ & $5$ & $7$ & $7$& $5$  \\
			\hline 
		\end{tabular}
	\end{center}
	Cân nặng trung bình của $28$ học sinh trên  là (kết quả làm tròn đến hàng phần trăm).
	 \choice
		{ \True $55{,}57$  (kg)}
		{ $54{,}57$ (kg)} 
		{ $53{,}57$  (kg)}
		{  $56{,}57$ (kg)}
	\loigiai{ 
	Cân nặng trung bình của $28$ học sinh trên là $$\overline{x}=\dfrac{47\cdot  4+51\cdot 5+55\cdot7+59\cdot7+63\cdot5}{28}\approx 55{,}57 \mathrm{kg}.$$}
\end{ex}

% nguồn: Ok_Toan11_TracNghiem_Lan2
\begin{ex}%[1D5N1-3]%[ma:d7d1e686]
	Điểm thi môn Toán của hai lớp 11A và 11B được cho trong bảng sau, biết mỗi lớp có  $50$ học sinh
	\begin{center}
		\begin{tabular}{|l|c|c|c|c|c|}
			\hline Điểm thi &  {$[0;2)$} & {$ [2;4)$} & {$[4;6) $} & {$[6;8)$} & {$[8;10)$} \\
			\hline Lớp 11A & $2$ & $4$ & $12$ & $28$ & $4$ \\ 
			\hline Lớp 11B & $4$ &  $10$ & $17$ & $14$ & $5$ \\
			\hline 
		\end{tabular}
	\end{center}
	Chọn khẳng  định đúng trong các khẳng định sau
	 \choice
		{ \True Điểm thi trung  bình của lớp 11A cao hơn lớp 11B}
		{ Điểm thi trung bình của lớp 11B cao hơn lớp  11A}
		{ Điểm thi  trung bình của lớp 11A bằng lớp 11B}
		{ Điểm thi trung bình  của lớp 11A gấp $2$ lần lớp 11B}
	\loigiai{ 
		Ta có bảng thống kê  điểm thi môn Toán theo giá trị đại diện
		\begin{center}
			\begin{tabular}{|l|c|c|c|c|c|}
				\hline Điểm thi & {$[0;2)$} & {$ [2;4)$} & {$[4;6) $} & {$[6;8)$} & {$[8;10)$} \\ 
				\hline Điểm thi đại  diện & $1$ & $3$ & $5$ & $7$ & $9$ \\
				\hline Lớp  11A & $2$ & $4$ & $12$ & $28$ & $4$ \\
				\hline Lớp 11B & $4$ & $10$ & $17$ & $14$ & $5$ \\ 
				\hline 
			\end{tabular}
		\end{center}
		Điểm thi trung bình của lớp 11A là $(1\cdot 2+3\cdot 4+5\cdot 12+7\cdot  28+9\cdot 4):50=6{,}12$ cm.\\
		Điểm  thi trung bình của lớp 11B là $(1\cdot 4+3\cdot 10+5\cdot 17+7\cdot 14+9\cdot 5):50=5{,}24$ cm.\\
	Vậy điểm thi trung bình của lớp 11A cao hơn  lớp 11B.}
\end{ex}

% nguồn: Ok_Toan11_TracNghiem_Lan2
\begin{ex}%[1D5N1-3]%[ma:fd23f19e]
	Tìm hiểu thời gian xem ti vi trong tuần  trước (đơn vị: giờ) của một số học sinh thu được kết quả sau
	\begin{center}
		\begin{tabular}{|l|c|c|c|c|c|}
			\hline Thời gian (giờ) & {$[0;5)$} & {$ [5;10)$} & {$[10;15) $} & {$[15;20)$} & {$[20;25)$}\\ 
			\hline Số học sinh& $8$ & $16$ &  $4$ & $2$& $2$\\
			\hline 
		\end{tabular}
	\end{center}
	Tính thời gian xem ti vi trung bình trong  tuần trước của các bạn học sinh này.
	\choice 
		{ $8{,}2375$ giờ} 
		 { \True $8{,}4375$ giờ}
		{ $8{,}754$ giờ} 
		{ $8{,}457$ giờ} 
	 \loigiai{
		Thời gian xem  ti vi trung bình trong tuần trước của các bạn học sinh này là
	$$\overline{x}=\dfrac{8\cdot 2{,} 5+16\cdot 7{,}5+4\cdot  12{,}5+2\cdot 17{,}5+2\cdot 22{,}5}{8+16+4+2+2}=8{,}4375\text{ giờ}.$$}
\end{ex}

% nguồn: Ok_Toan11_TracNghiem_Lan2
\begin{ex}%[1D5N1-3]%[ma:c81ca021]
	Điều tra về số tiền mua đồ dùng học tập trong một tháng của $40$ học sinh, ta có mẫu số  liệu như sau (đơn vị: nghìn đồng)
	\begin{center}
		\begin{tabular}{|l|c|c|c|c|c|c|c|c|}
			\hline Giá trị  $(x)$ & {$[10;15) $} & {$[15;20)$} & {$[20;25)$}& {$[25;30)$} & {$ [30;35)$} & {$[35;40) $} & Cộng\\
			\hline  Tần số $n$ & $2$ & $5$ & $15$ & $8$& $9$& $1$& $N=40$\\
			\hline 
		\end{tabular}
	\end{center}
	Số trung bình của mẫu số liệu  là
	 \choice
		{  $28$}
		 { $23$}
		{  $32$}
		{ \True $25$} 
	\loigiai{ 
	 Số trung bình cộng của mẫu là $$\overline{x}=\dfrac{12{,}5\cdot 2+17{,}5\cdot5+22{,}5\cdot15+27{,}5\cdot8+32{,}5\cdot9+37{,}5\cdot1}{40}=25.$$}
\end{ex}

% nguồn: Ok_Toan11_TracNghiem_Lan2
\begin{ex}%[1D5N1-3]%[ma:11c61628]
	Thời gian (phút) đi từ nhà đến  trường của bạn A trong $35$ ngày được cho trong bảng sau
	\begin{center}
		\begin{tabular}{|l|c|c|c|c|c|}
			\hline Thời gian (phút)& {$[19;21)$} &  {$ [21;23)$} & {$[23;25) $} & {$[25;27)$} & {$[27;29)$} \\
			 \hline Tần số & $5$ & $8$ & $7$ & $9$& $6$ \\
			\hline 
		\end{tabular}
	\end{center}
	Tìm thời gian trung bình  bạn A thường dành ra để đi từ nhà đến trường là (kết quả làm tròn đến hàng phần chục).
	 \choice
		{ \True $24{,}2$  phút}
		{ $25{,}4$ phút} 
		{ $25{,}8$ phút} 
		 { $24{,}8$ phút}
	 \loigiai{
		Ta có bảng thống kê  thời gian đi từ nhà đến trường của A theo giá trị đại diện
		\begin{center}
			\begin{tabular}{|l|c|c|c|c|c|}
				\hline Thời gian (phút)& {$[19;21)$} & {$ [21;23)$} & {$[23;25) $} & {$[25;27)$} & {$[27;29)$}  \\
				\hline Thời gian đại diện & $20$ & $22$ & $24$ & $26$& $28$ \\ 
				\hline Tần số & $5$ & $8$ &  $7$ & $9$& $6$ \\
				\hline 
			\end{tabular}
		\end{center}
		Thời gian trung bình bạn A thường dành ra để đi từ nhà đến  trường là
	 $$(20\cdot 5+22\cdot 8+24\cdot 7+26\cdot 9+28\cdot 6):35=24{,}2.$$}
\end{ex}

% nguồn: Ok_TracNghiemLop11_Lan1
\begin{ex}%[1D5N1-3]%[ma:167e6d6c]
	Điều  tra về số tiền mua đồ dùng học tập (đơn vị: nghìn đồng) trong một tháng của $40$ học sinh ta có mẫu số liệu như sau:
	\begin{longtable}{|c|c|c|c|c|c|c|c|}
		\hline 
		Số  tiền&$[10;15)$&$[15;20)$&$[20;25)$&$[25;30)$&$[30;35)$&$[35;40)$&Cộng\\
		\hline 
		Số học sinh&$2$&$5$&$15$&$8$&$9$&$1$&$N=40$\\ 
		 \hline
	\end{longtable}
	Số trung bình của  mẫu số liệu là
	\choice 
		 {$22{,}5$}
		 {\True $25$}
		{$25{,}5$} 
		 {$27$}
	\loigiai{ 
		Ta có bảng số liệu với giá trị đại diện  như sau:
		\begin{longtable}{|c|c|c|c|c|c|c|}
			 \hline
			 Số tiền&$[10;15)$&$[15;20)$&$[20;25)$&$[25;30)$&$[30;35)$&$[35;40)$\\
			\hline 
			 Giá trị đại diện&$12{,}5$&$17{,}5$&$22{,}5$&$27{,}5$&$32{,}5$&$37{,}5$\\
			 \hline
			Số học  sinh&$2$&$5$&$15$&$8$&$9$&$1$\\
			 \hline
		\end{longtable}
		Số trung bình  $\overline{x}=\dfrac{12{,}5\cdot 2+17{,}5\cdot 5+22{,}5\cdot 15+27{,}5\cdot 8+32{,}5\cdot 9+37{,}5\cdot 1}{40}=25$.
	 }
\end{ex}

% nguồn: Ok_TracNghiemLop11_Lan1
\begin{ex}%[1D5N1-3]%[ma:f0ec6f28]
	Cân nặng của $28$ học sinh nam lớp $11$ được cho  như sau
	\begin{center}
		\begin{tabular}{llllllllllllll}
			$55{,}4$  & $62{,}6$ & $54{,}2$ & $56{,}8$ & $58{,}8$ & $59{,}4$ & $60{,}7$ & $58$ & $59{,}5$ & $63{,}6$ & $61{,}8$ & $52{,}3$ & $63{,}4$ & $57{,}9$\\
			$49{,}7$ & $45{,}1$  & $56{,}2$ & $63{,}2$ & $46{,}1$ & $49{,}6$ & $59{,}1$ & $55{,}3$ & $55{,}8$ & $45{,}5$ & $46{,}8$ & $54$ & $49{,}2$ & $52{,}6$
		\end{tabular}
	\end{center}
	Số trung bình của mẫu số liệu ghép nhóm trên xấp xỉ bằng 
	 \choice
		 {\True $55{,}6$}
		 {$65{,}5$}
		{$48{,}8$} 
		 {$57{,}7$}
	\loigiai{ 
	\begin{center}
		\begin{tabular}{|c|c|c|c|c|c|}
			 \hline
			Cân nặng &{$[45; 49)$} &{$[49; 53)$} &{$[53; 57)$}  &{$[57; 61)$} &{$[61; 65)$} \\
			\hline 
			Giá trị đại diện &  $47$ & $51$ & $55$ & $59$ & $63$ \\
			 \hline
			Số  học sinh & $4$ & $5$ & $7$ & $7$ & $5$ \\
			 \hline
		\end{tabular}
	\end{center}
	Cân nặng trung bình  của học sinh trong lớp $11$ là\\
	$\overline{x}= \dfrac{47 \cdot 4 + 51 \cdot  5 + 55 \cdot 7 + 59 \cdot 7 + 63 \cdot 5 }{28} \approx 55{,}6$ (kg).
	 }
\end{ex}

% nguồn: Ok_TracNghiemLop11_Lan1
\begin{ex}%[1D5N1-3]%[ma:5d69a704]
	Doanh thu bán hàng trong $20$  ngày được lựa chọn ngẫu nhiên của một của hàng được ghi lại ở bảng sau (đơn vị: triệu đồng):
	\begin{center}
		\begin{tabular}{|c|c|c|c|c|c|}
			\hline 
			Doanh  thu & $\left[5;7\right)$ & $\left[7;9\right)$ &$\left[9;11\right)$ &$\left[11;13\right)$ &$\left[13;15\right)$ \\
			 \hline
			Số ngày & $2$  &$7$ &$7$ &$3$ &$1$ \\
			 \hline
		\end{tabular}
	\end{center}
	Số trung bình của  mẫu số liệu trên thuộc khoảng nào trong các khoảng dưới đây?
	 \choice
		{$\left[7;9\right)$} 
		{\True $\left[9;11\right)$} 
		{$\left[11;13\right)$} 
		{$\left[13;15\right)$} 
	\loigiai{ 
		Bảng tần số ghép nhóm theo giá trị đại  diện là
		\begin{center}
			\begin{tabular}{|c|c|c|c|c|c|}
				\hline 
				 Doanh thu & $\left[5;7\right)$ & $\left[7;9\right)$ &$\left[9;11\right)$ &$\left[11;13\right)$ &$\left[13;15\right)$ \\
				\hline 
				Giá trị đại diện &  $6$ &$8$ &$10$ &$12$ &$14$ \\
				 \hline
				Số học  sinh & $2$ &$7$ &$7$ &$3$ &$1$ \\
				\hline 
			\end{tabular}
		\end{center}
		Số trung bình $\overline{x}=\dfrac{2\cdot 6+7\cdot 8+7\cdot  10+3\cdot 12+1\cdot 14}{20}=9{,}4$.
	 }
\end{ex}

% nguồn: Ok_TracNghiemLop11_Lan1
\begin{ex}%[1D5N1-3]%[ma:e583d005]
	Tìm cân nặng trung bình của học sinh lớp $11$D cho trong bảng sau, làm tròn đến hàng phần trăm. 
	\begin{center}
		\begin{tabular}{|c|c|c|c|c|c|c|}
			 \hline Cân nặng & {$[40,5; 45,5)$} & {$[45,5; 50,5)$} & {$[50,5; 55,5)$} & {$[55,5; 60,5)$} & {$[60,5; 65,5)$} & {$[65,5; 70,5)$} \\
			\hline Số học sinh &  $10$ & $7$ & $16$ & $4$ & $2$ & $3$ \\
			\hline 
		\end{tabular}
	\end{center}
	 \choice
		{$51{,}8$} 
		{\True  $51{,}81$}
		{$52$} 
		{$51{,}809$} 
	\loigiai{ 
		Trong mỗi khoảng cân nặng, giá trị đại diện là trung bình cộng của giá  trị hai đầu mút nên ta có bảng sau
		\begin{center}
			\begin{tabular}{|c|c|c|c|c|c|c|}
				\hline Cân nặng (kg) & $43$ & $48$ & $53$ & $58$ &  $63$ & $68$ \\
				\hline Số học sinh  & $10$ & $7$ & $16$ & $4$ & $2$ & $3$ \\
				\hline 
			\end{tabular}
		\end{center}
		Tổng số học sinh là  $n=42$.\\
		Cân nặng trung  bình của học sinh lớp 11D là
		$$\bar{x}=\dfrac{10\cdot 43+7\cdot  48+16\cdot 53+4\cdot 58+2\cdot 63+3\cdot 68}{42} \approx 51{,}81\,(\text{kg}).$$
	 }
\end{ex}

% nguồn: Ok_TracNghiemLop11_Lan1
\begin{ex}%[1D5N1-3]%[ma:07ea2b9a]
	Một cửa hàng đã ghi lại số tiền bán xăng cho $35$ khách hàng đi xe máy. Mẫu số  liệu gốc có dạng: $x_1$, $x_2$, $\ldots$, $x_{35}$­ trong đó $x_i$ là số tiền bán xăng cho khách hàng thứ $i$. Vì một lí do nào đó, cửa hàng chỉ có mẫu số liệu ghép nhóm dạng sau:
	\begin{center}
		\begin{tabular}{|c|c|c|c|c|}
			 \hline
			\textbf{Số tiền (nghìn đồng)}&  $[0;30)$ & $[30;60)$ & $[60;90)$ & $[90;120)$\\
			 \hline
			\textbf{Số khách hàng}& $3$ & $15$ & $10$  & $7$\\
			\hline 
		\end{tabular}
	\end{center}
	Tính số tiền bán xăng  trung bình của $35$ khách hàng.
	\choice 
		 {$35$}
		{$17{,}5$} 
		{\True  $63$}
		 {$60$}
	\loigiai{ 
		Ta có bảng  giá trị đại diện như sau:
		\begin{center}
			\begin{tabular}{|c|c|c|c|c|}
				\hline 
				\textbf{Số tiền (nghìn đồng)}& $15$ &  $45$ & $75$ & $105$\\
				\hline 
				\textbf{Số khách hàng}& $3$  & $15$ & $10$ & $7$\\
				\hline 
			\end{tabular}
		\end{center}
		Số tiền bán xăng trung bình của $35$ khách hàng là \[\dfrac{3\cdot 15+15\cdot 45+10\cdot  75+7\cdot 105}{35}=63\,\text{(nghìn đồng)}.\]
	 }
\end{ex}

% nguồn: Ok_TracNghiemLop11_Lan1
\begin{ex}%[1D5N1-3]%[ma:18f861ed]
	Điểm thi môn Toán (thang điểm 100, điểm được làm tròn đến 1) của 60  thí sinh được cho trong bảng sau
	\begin{center}
		\begin{tabular}{|l|c|c|c|c|c|}
			\hline Điểm & $[0-9{,}5)$ & $[9{,}5-19{,}5)$ & $[19{,}5-29{,}5)$ & $[29{,}5-39{,}5)$ &  $[39{,}5-49{,}5)$ \\
			\hline Số thí sinh & $1 $& $2$ & $4$ & $6$ & $15$ \\ 
			\hline Điểm &  $[49{,}5-59{,}5)$ & $[59{,}5-69{,}5)$ & $[69{,}5-79{,}5)$ & $[79{,}5-89{,}5)$ & $[89{,}5-99{,}5)$ \\
			\hline Số  thí sinh & $12$ & $10$ & $6$ & $3$ & $1$ \\
			\hline 
		\end{tabular}
	\end{center}
	Tìm tứ phân vị thứ hai của mẫu số liệu. 
	\choice 
		{\True  $Q_2\approx 51{,}17$}
		{$Q_2\approx 51{,}67$} 
		 {$Q_2\approx 49{,}5$}
		{$Q_2\approx  41{,}3$}
	\loigiai{Cỡ  mẫu là $n=60$.\\
		Tứ phân vị thứ nhất $Q_2$ là  $\dfrac{x_{30}+x_{31}}{2}$. Do $x_{30}$, $x_{31}$ đều thuộc nhóm $[49{,}5; 59{,}5)$ nên nhóm này chứa $Q_2$. \\Do đó, $p=6$; $a_6=49{,}5$; $m_6=12$; $m_1+\cdots+m_5=28$; $a_7-a_6=10$ và ta có
		\[Q_2=49,5+\dfrac{\dfrac{60}{2}-28}{12}\cdot  10\approx51{,}17.\]
	 }
\end{ex}

% nguồn: Ok_TracNghiemLop11_Lan1
\begin{ex}%[1D5N1-3]%[ma:5c17e14c]
	Tìm hiểu thời gian xem ti vi trong tuần trước của một số học sinh thu được kết quả  sau
	\begin{center}
		\begin{tabular}{|l|c|c|c|c|c|}
			\hline Thời gian (giờ) & {$[0; 5)$} & {$[5; 10)$} & {$[10;  15)$} & {$[15; 20)$} & {$[20; 25)$} \\
			\hline Số học sinh  & $8$ & $16$ & $4$ & $2$ & $2$ \\
			 \hline
		\end{tabular}
	\end{center}
	Tính thời gian xem ti vi trung bình trong tuần trước của các bạn học sinh  này.
	\choice 
		 {\True $8{,}4375$}
		 {$8{,}125$}
		 {$8{,}75$}
		 {$8{,}28125$}
	 \loigiai{
		Trong mỗi khoảng, giá trị đại diện là trung bình cộng của giá trị hai  đầu mút nên ta có bảng sau
		\begin{center}
			\begin{tabular}{|c|c|c|c|c|c|c|}
				\hline Thời gian (giờ) &$[0; 5)$ & {$[5; 10)$} &  {$[10; 15)$} & $[15;20)$ & $[20;25)$ \\
				\hline Số học sinh & $8$ &  $16$ & $4$ & $2$ & $2$ \\
				\hline Giá trị đại diện &  $2{,}5$ & $7{,}5$ & $12{,}5$ & $17{,}5$ & $22{,}5$ \\
				 \hline
			\end{tabular}
		\end{center}
		Thời gian xem ti vi trung bình  trong tuần trước của các bạn học sinh này là
		$$\bar{x}=\dfrac{8\cdot 2{,}5+16\cdot 7{,}5+4\cdot 12{,}5+2\cdot  17{,}5+2\cdot 22{,}5}{8+16+4+2+2}=8{,}4375.$$
	} 
\end{ex}

% nguồn: Ok_TracNghiemLop11_Lan1
\begin{ex}%[1D5N1-3]%[ma:9aeb1439]
	Doanh thu bán  hàng trong $20$ ngày được lựa chọn ngẫu nhiên của một của hàng được ghi lại ở bảng sau (đơn vị: triệu đồng)
	\begin{center}
		{\tabcolsep =  2mm
		\begin{tabular}{|c|c|c|c|c|c|}
			 \hline
			Doanh thu & $[5; 7)$  & $[7; 9)$ & $[9; 11)$ & $[11; 13)$ & $[13; 15)$ \\
			\hline 
			Số  ngày & $2$ & $7$ & $7$ & $3$ & $1$\\
			\hline 
		\end{tabular}
		 }
	\end{center}
	Số trung bình của mẫu số liệu trên thuộc khoảng nào trong các  khoảng dưới đây?
	\choice 
		{$[7; 9)$} 
		{\True $[9; 11)$} 
		{$[11;  13)$}
		{$[13;  15)$}
	 \loigiai{
	Bảng tần số ghép nhóm theo giá trị đại diện là 
	\begin{center}
		{\tabcolsep  = 2mm
		\begin{tabular}{|c|c|c|c|c|c|}
			\hline 
			Doanh thu & $[5; 7)$ & $[7; 9)$ & $[9; 11)$ & $[11; 13)$ & $[13; 15)$  \\
			\hline 
			Giá trị đại diện & $6$ & $8$ & $10$  & $12$ & $14$ \\
			\hline 
			Số ngày &  $2$ & $7$ & $7$ & $3$ & $1$\\
			\hline 
		\end{tabular}
		} 
	\end{center}
	Số trung bình là $\overline{x}= \dfrac{2 \cdot 6 + 7 \cdot 8 + 7 \cdot 10 + 3  \cdot 12 + 1 \cdot 14}{20}= 9{,}4$.
	} 
\end{ex}

% nguồn: Ok_TracNghiemLop11_Lan1
\begin{ex}%[1D5N1-3]%[ma:146cd69f]
	Tỉ giá giữa đồng tiền mỗi nước với đồng Đô la Mỹ (USD) là một chỉ số quan trọng của kinh tế. Một nhà kinh doanh đã theo dõi trong $25$ ngày liên tiếp tỉ giá giữa đồng EUR (đồng tiền chung của nhiều quốc gia trong Liên minh châu Âu) và đồng USD vào đầu giờ làm việc mỗi ngày. Do sự biến động tỉ giá giữa hai ngày liền nhau thường không nhiều nên nhà kinh doanh đó đã ghép nhóm số liệu thu được  và biểu diễn qua bảng.
	\begin{center}
		\begin{tabular}{|c|c|c|c|c|c|}
			\hline Tỉ  giá EUR/USD & {$[1; 1{,}04)$} & {$[1{,}04; 1{,}08)$} & {$[1{,}08; 1{,}12)$} & {$[1{,}12; 1{,}16)$} & {$[1{,}16; 1{,}2)$} \\
			\hline Tần số & $6$ & $4$ & $5$ & $7$ & $3$ \\ 
			\hline 
		\end{tabular}
	\end{center}
	Tính tỉ giá trung bình của đồng EUR/USD  trong $25$ ngày mà nhà kinh doanh theo dõi.
	\choice 
		{$1{,}1152$} 
		 {$1{,}1052$}
		{\True $1{,}0952$} 
		 {$1{,}0852$}
	\loigiai{Số trung bình của mẫu số liệu cho trong bảng  là
		$$ 
		\overline{x}=\dfrac{6 \cdot 1{,}02+4 \cdot  1{,}06+5 \cdot 1{,}1+7 \cdot 1{,}14+3 \cdot 1{,}18}{25} \approx 1{,}0952 .$$
	 }
\end{ex}

% nguồn: Ok_TracNghiemLop11_Lan1
\begin{ex}%[1D5N1-3]%[ma:538f358b]
	 Kết quả khảo sát cân nặng của $25$ quả cam ở lô hàng $A$ được cho ở bảng sau
	\begin{center}
		{\tabcolsep  = 2mm
		\begin{tabular}{|c|c|c|c|c|c|}
			 \hline
			Cân nặng (g) & $[150; 155)$ & $[155;  160)$ & $[160; 165)$ & $[165; 170)$ & $[170; 175)$ \\
			\hline 
			Số quả cam ở lô hàng A &  $1$ & $3$ & $7$ & $10$ & $4$\\
			 \hline
		\end{tabular}
		 }
	\end{center}
	Cân nặng trung bình của mỗi quả  cam ở lô hàng $A$ xấp xỉ bằng
	 \choice
		 {$162{,}7$}
		{\True $161{,}7$} 
		{$163{,}7$} 
		{$164{,7}$} 
	\loigiai{ 
		Ta có bảng thống kê số  lượng cam theo giá trị đại diện
		\begin{center}
			 {\tabcolsep = 2mm
			\begin{tabular}{|c|c|c|c|c|c|}
				 \hline
				Cân nặng (g) & $[150; 155)$ & $[155; 160)$ & $[160; 165)$ & $[165; 170)$ & $[170; 175)$  \\
				 \hline
				Cân nặng đại diện & $152{,}5$ & $157{,}5$ & $162{,}5$ &  $167{,}5$ & $172{,}5$ \\
				 \hline
				Số quả cam ở lô hàng A & $2$ & $6$ & $12$ &  $4$ & $1$ \\
				\hline 
			\end{tabular}
			 }
		\end{center}
		Cân nặng trung bình của mỗi quả cam ở lô hàng  A là
		$$\overline{x}= \dfrac{2 \cdot 152{,}5  + 6 \cdot 157{,}5 + 12 \cdot 162{,}5 + 4 \cdot 167{,}5 + 172{,}5}{25} \approx 161{,}7 \text{ (g)}.$$
	} 
\end{ex}

% nguồn: Ok_Toan11_TracNghiem_Lan2
\begin{ex}%[1D5H1-3]%[ma:29440747]
	Thống  kê nhiệt độ tại một địa điểm trong $40$ ngày, ta có bảng số liệu sau
	\begin{center}
		\begin{tabular}{|c|c|c|c|c|}
			\hline 
			Nhiệt độ $(^\circ C)$ & $[19;22)$  & $[22;25)$ & $[25;28)$ & $[28;31)$ \\
			\hline 
			Số ngày & $7$ & $15$ & $12$ & $6$  \\
			\hline 
		\end{tabular}
	\end{center}
	Nhiệt độ trung bình  là
	 \choice
		{$25{,}456$} 
		{\True $24{,}775$} 
		{$23{,}021$} 
		{$22{,}036$} 
	 \loigiai{
		Nhiệt độ trung bình  là
		\[\overline{x}=\dfrac{7\cdot 20{,}5 + 15\cdot 23{,}5 +  12\cdot 26{,}5 + 6\cdot 29{,}5}{40} = 24{,}775.\]
	 }
\end{ex}

% nguồn: Ok_Toan11_TracNghiem_Lan2
\begin{ex}%[1D5H1-3]%[ma:e9ab19c1]
	Điểm thi giữa kỳ 1 môn toán của một lớp  học sinh khối 11 thu được mẫu số liệu ghép nhóm sau
	\begin{longtable}{|*{4}{c|}}
		 \hline
		Điểm thi& [$1{,}5$; $4{,}5$) & [$4{,}5$; $7{,}5$) & [$7{,}5$; $10{,}5$) \\ 
		 \hline
		Số học sinh & $7$  & $18$ & $10$ \\
		 \hline
	\end{longtable}
	Trung vị của mẫu  số liệu ghép nhóm trên là
	 \choice
		 {$6$}
		 {$4{,}5$}
		 {$7{,}5$}
		{\True$6{,}25$} 
	 \loigiai{
		Cỡ mẫu là  $n=7+18+10=35$.\\
		Gọi $x_1, x_2, \ldots, x_{35}$ là số điểm của 35 học sinh và giả sử dãy này được sắp xếp theo thứ tự không giảm. Khi đó: trung vị là $x_{18}$  thuộc nhóm $\left[4{,}5; 7{,}5\right)$.\\
		Ta xác định  được $n=35$, $n_m=18$, $C=7$, $u_m=4{,}5$, $u_{m+1}=7{,}5$.\\
		Trung vị của mẫu số liệu ghép nhóm là: $M_e=4{,}5+\dfrac{\dfrac{35}{2}-7}{18} \left(7{,}5-4{,}5\right)=6{,}25$. 
	} 
\end{ex}

% nguồn: Ok_Toan11_TracNghiem_Lan2
\begin{ex}%[1D5H1-3]%[ma:9076f33c]
	Cho mẫu số liệu ghép nhóm về thống kê chiều cao của $35$ cây bạch đàn trong rừng, ta có bảng  số liệu sau
	\begin{center}
		\begin{tabular}{|c|c|c|c|c|}
			 \hline
			Khoảng chiều cao & $[6{,}5;7{,}0)$ &  $[7{,}0;7{,}5)$ & $[7{,}5;8{,}0)$ & $[8{,}0;8{,}5)$ \\
			\hline 
			Giá trị đại diện &  $6{,}75$ & $7{,}25$ & $7{,}75$ & $8{,}25$ \\
			 \hline
			Tần số & $6$ & $15$ & $11$ & $3$ \\ 
			 \hline
		\end{tabular}
	\end{center}
	Tính chiều cao trung bình của $35$ cây  bạch đàn trên.
	\choice 
		{\True $7{,}407$\,m} 
		 {$4{,}707$\,m}
		{$7{,}704$\,m} 
		{$7{,}5$\,m} 
	\loigiai{ 
		Chiều cao trung bình của $35$  cây bạch đàn là
		\[\overline{x}=\dfrac{6{,}75\cdot6 +  7{,}25\cdot15 + 7{,}75\cdot11 + 8{,}25\cdot3}{35}\approx 7{,}407.\]
	 }
\end{ex}

% nguồn: Ok_Toan11_TracNghiem_Lan2
\begin{ex}%[1D5H1-3]%[ma:43d20736]
	Cân nặng của $28$ học sinh  nam lớp $11$ được cho như sau
	\begin{center}
		\begin{tabular}{|c|c|c|c|c|c|c|c|c|c|c|c|c|c|}
			 \hline
			 $55{,}4$&$62{,}6$&$54{,}2$&$56{,}8$&$58{,}8$&$59{,}4$&$60{,}7$&$58$&$59{,}5$&$63{,}6$&$61{,}8$&$52{,}3$&$63{,}4$&$57{,}9$\\
			\hline 
			 $49{,}7$&$45{,}1$&$56{,}2$&$63{,}2$&$46{,}1$&$49{,}6$&$59{,}1$&$55{,}3$&$55{,}8$&$45{,}5$&$46{,}8$&$54$&$49{,}2$&$52{,}6$\\
			 \hline
		\end{tabular}
	\end{center}
	Số trung bình của mẫu số liệu ghép nhóm trên xấp  xỉ bằng
	 \choice
		{\True $55{,}6$} 
		 {$65{,}5$}
		 {$48{,}8$}
		{$57{,}7$} 
	 \loigiai{
		\begin{center}
			\begin{tabular}{|c|c|c|c|c|c|}
				 \hline
				Cân nặng & $[45;49)$ & $[49;53)$ &  $[53;57)$ & $[57;61)$ & $[61;65)$\\
				 \hline
				Giá trị đại diện & $47$ & $51$ &  $55$ & $59$ & $63$\\
				\hline 
				Số học sinh & $4$ & $5$ &  $7$ & $7$ & $5$\\
				\hline 
			\end{tabular}
		\end{center}
		Cân nặng trung bình của  học sinh trong lớp $11$ xấp xỉ
		\[\dfrac{47\cdot4+51\cdot5+55\cdot7+59\cdot7+63\cdot5}{28}\approx  55{,}6\text{ (kg)}.
		\] 
	 }
\end{ex}

% nguồn: Ok_Toan11_TracNghiem_Lan2
\begin{ex}%[1D5H1-3]%[ma:d90c57ca]
	 Kết quả khảo sát cân nặng của $25$ quả cam ở lô hàng $A$ được cho ở bảng sau:
	\begin{center}
		\begin{tabular}{|c|c|c|c|c|c|}
			\hline 
			Cân nặng (g) & $[150;155)$ & $[155;160)$ & $[160;165)$ & $[165;170)$ & $[170;175)$\\ 
			\hline 
			Số quả cam ở lô hàng $A$ & $2$ & $6$ & $12$ & $4$  & $1$\\
			 \hline
		\end{tabular}
	\end{center}
	Cân nặng trung bình của mỗi quả cam  ở lô hàng $A$ xấp xỉ bằng
	\choice 
		 {$162{,}7$}
		 {\True $161{,}7$}
		 {$163{,}7$}
		 {$164{,}7$}
	 \loigiai{
		Ta có bảng thống kê số lượng  cam theo giá trị đại diện
		\begin{center}
			\begin{tabular}{|c|c|c|c|c|c|}
				 \hline
				Cân nặng đại diện (g) &$152{,}5$&$157{,}5$&$162{,}5$& $167{,}5$ & $172{,}5$\\ 
				\hline 
				Số quả  cam ở lô hàng $A$ & $2$ & $6$ & $12$ & $4$ & $1$\\
				 \hline
			\end{tabular}
		\end{center}
		Cân nặng trung bình của mỗi quả cam ở lô hàng $A$  xấp xỉ bằng
		\[\dfrac{2\cdot  152{,}5+6\cdot157{,}5+12\cdot162{,}5+4\cdot167{,}5+1\cdot172{,}5}{25}\approx 161{,}7\text{(g)}.
		\] 
	} 
\end{ex}

% nguồn: Ok_Toan11_TracNghiem_Lan2
\begin{ex}%[1D5H1-3]%[ma:d572fb59]
	Doanh thu bán hàng trong 20 ngày  được lựa chọn ngẫu nhiên của một của hàng được ghi lại ở bảng sau (đơn vị: triệu đồng):
	\begin{longtable}{|*{6}{c|}}
		 \hline
		 Doanh thu & [$5$; $7$) & [$7$; $9$) & [$9$; $11$) & [$11$; $13$) & [$13$; $15$) \\
		\hline 
		Số ngày  & $2$& $7$ & $7$ & $3$ & $1$ \\
		 \hline
	\end{longtable}
	Trung vị của mẫu số liệu trên thuộc khoảng  nào trong các khoảng dưới đây?
	 \choice
		 {$\left[7; 9\right)$}
		{\True $\left[9;  11\right)$}
		{$\left[11; 13\right)$} 
		{$\left[13; 15\right)$} 
	 \loigiai{
		 Gọi $x_1, x_2, \ldots, x_{20}$ là doanh thu bán hàng trong 20 ngày xếp theo thứ tự không giảm.\\
		Khi đó: $x_1, x_2  \in \left[5; 7\right)$, $x_3, \ldots, x_9 \in \left[7; 9\right)$, $x_9, \ldots, x_{16} \in \left[9; 11\right)$,
		$x_{17}$, $ \ldots$, $ x_{19} \in \left[11; 13\right)$,  $x_{20} \in \left[13; 15\right)$.\\
		Do đó trung vị của mẫu số liệu thuộc nhóm $[9;  11)$.
	 }
\end{ex}

% nguồn: Ok_Toan11_TracNghiem_Lan2
\begin{ex}%[1D5H1-3]%[ma:0830018e]
	Kiểm tra cân nặng $25$ quả  dưa của hai nhà vườn $A$ và $B$ cho kết quả như bảng sau
	\begin{center}
		\begin{tabular}{|c|c|c|c|c|c|}
			\hline 
			Cân nặng(kg) & $[1{,}5;2{,}5)$ & $[2{,}5;3{,}5)$ & $[3{,}5;4{,}5)$ &  $[4{,}5;5{,}5)$ & $[5{,}5;6{,}5)$\\
			 \hline
			Số quả của vườn $A$ & $2$ & $6$ & $4$ &  $10$ & $3$\\
			 \hline
			Số quả của vườn $B$ & $3$ & $7$ & $8$  & $5$ & $2$\\
			\hline 
		\end{tabular}
	\end{center}
	Trung bình mỗi quả dưa vườn  $B$ nặng là
	\choice 
		{\True $3{,}84$  kg}
		{$4{,}84$ kg} 
		{$2{,}84$ kg} 
		{$3{,}74$ kg} 
	 \loigiai{
		Trung bình mỗi quả dưa vườn $B$ nặng  số kg
		\[\dfrac{3\cdot 2+7\cdot 3+8\cdot 4+5\cdot 5+2\cdot 6}{25}=3{,}84\text{ (kg)  }.
		 \]
	} 
\end{ex}

% nguồn: Ok_Toan11_TracNghiem_Lan2
\begin{ex}%[1D5H1-3]%[ma:821ce6c1]
	Cho mẫu số liệu ghép nhóm về thống kê điểm số (thang điểm $20$) của $100$ học sinh tham  dự kỳ thi học sinh giỏi toán, ta có bảng số liệu sau
	\begin{longtable}{|*{7}{c|}}
		 \hline
		Điểm & [$8$; $10$) & [$10$; $12$) & [$12$; $14$) & [$14$; $16$) & [$16$; $18$) &  [$18$; $20$) \\
		 \hline
		Số học sinh & $6$ & $21$ & $30$ & $25$ & $14$  & $4$\\
		\hline 
	\end{longtable}
	Tìm trung vị của mẫu số liệu ghép nhóm trên. (Kết quả  làm tròn đến hàng phần trăm).
	 \choice
		 {$12{,}18$}
		 {$12{,}81$}
		{$13{,}35$} 
		{\True $13{,}53$} 
	\loigiai{ 
		Số phần tử của mẫu là $n=100$. Ta có $\dfrac{n}{2}=\dfrac{100}{2}=50$.\\ 
		Vì cỡ mẫu $n=100$ là số chẵn nên trung vị của mẫu số liệu số $x_1 \le x_2 \le x_3 \le \ldots x_{100}$ là  $\dfrac{1}{2} \left(x_{50}+x_{51} \right)\in \left[12; 14\right)$.\\
		Vậy trung vị của mẫu số  liệu ghép nhóm trên là: $M_e=12+\left(\dfrac{50-27}{30} \right)\cdot 2\approx 13{,}53$.
	} 
\end{ex}

% nguồn: Ok_Toan11_TracNghiem_Lan2
\begin{ex}%[1D5H1-3]%[ma:f61c233c]
	Khảo sát thời gian tự học bài ở nhà của một số em  học sinh lớp $11$ thu được mẫu số liệu ghép nhóm như sau
	\begin{center}
		\begin{tabular}{|c|c|c|c|c|c|}
			\hline 
			Thời gian(phút) &  $[0;30)$ & $[30;60)$ & $[60;90)$ & $[90;120)$ & $[120;150)$\\
			 \hline
			Số học sinh & $9$ & $10$ & $9$ & $15$  & $7$\\
			\hline 
		\end{tabular}
	\end{center}
	Thời gian trung bình mỗi em tự học ở  nhà là
	 \choice
		{$75{,}5$  phút}
		{$85{,}6$  phút}
		 {\True $75{,}6$ phút}
		{$94{,}5$ phút} 
	 \loigiai{
		\begin{center}
			\begin{tabular}{|c|c|c|c|c|c|}
				 \hline
				Thời gian(phút) & $[0;30)$ & $[30;60)$ &  $[60;90)$ & $[90;120)$ & $[120;150)$\\
				\hline 
				Giá trị đại diện  & $15$ & $45$ & $75$ & $105$ & $135$\\
				 \hline
				Số học sinh  & $9$ & $10$ & $9$ & $15$ & $7$\\
				 \hline
			\end{tabular}
		\end{center}
		Thời gian trung bình các em tự học ở nhà 
		 \[\dfrac{9\cdot15+10\cdot45+9\cdot75+15\cdot105+7\cdot135}{50}=75{,}6\text{ (phút)}.
		 \]
	 }
\end{ex}

% nguồn: Ok_Toan11_TracNghiem_Lan2
\begin{ex}%[1D5H1-3]%[ma:25b74d86]
	Tiền lãi (nghìn  đồng) trong $30$ ngày được khảo sát ở một quầy bán báo
	\begin{center}
		\begin{tabular}{|c|c|c|c|c|c|c|}
			 \hline
			 Lớp tiền lãi & $[29{,}5;40{,}5)$ & $[40{,}5;51{,}5)$ & $[51{,}5;62{,}5)$ & $[62{,}5;73{,}5)$ & $[73{,}5;84{,}5)$ & $[84{,}5;95{,}5)$ \\
			 \hline
			Tần số & $3$ & $5$ &  $7$ & $6$ & $5$ & $4$ \\
			 \hline
		\end{tabular}
	\end{center}
	Tính số trung bình của mẫu số  liệu.
	 \choice
		{\True  $63{,}23$}
		{$61{,}25$} 
		{$62{,}35$} 
		{$63{,}23$} 
	 \loigiai{
		Số trung bình của mẫu  số liệu là
		\[\overline{x}=\dfrac{3\cdot 35 + 5\cdot 46 + 7\cdot 57 + 6\cdot 68  + 5\cdot 79 + 4\cdot 90}{30} \approx 63{,}23\text{ (nghìn đồng)}.\]
	} 
\end{ex}

% nguồn: Ok_Toan11_TracNghiem_Lan2
\begin{ex}%[1D5H1-3]%[ma:be2e037d]
	Doanh thu bán hàng trong $20$ ngày được lựa chọn  ngẫu nhiên của một cửa hàng được ghi lại ở bảng sau (đơn vị triệu đồng)
	\begin{center}
		\begin{tabular}{|c|c|c|c|c|c|}
			 \hline
			Doanh thu & $[5;7)$ & $[7;9)$ & $[9;11)$ & $[11;13)$ & $[13;15)$ \\ 
			\hline 
			Số ngày & $2$ & $7$ & $7$ &  $3$ & $1$ \\
			\hline 
		\end{tabular}
	\end{center}
	Số trung bình của mẫu số  liệu trên thuộc khoảng nào trong các khoảng dưới đây
	 \choice
		 {$[7;9)$}
		 {\True $[9;11)$}
		{$[11;13)$} 
		{$[13;15)$} 
	 \loigiai{
		Bảng tần số ghép nhóm theo giá trị đại diện  là
		\begin{center}
			\begin{tabular}{|c|c|c|c|c|c|}
				\hline 
				Doanh thu & $[5;7)$ & $[7;9)$ & $[9;11)$  & $[11;13)$ & $[13;15)$ \\
				 \hline
				Giá trị đại diện & $6$ &  $8$ & $10$ & $12$ & $14$ \\
				 \hline
				Số  ngày & $2$ & $7$ & $7$ & $3$ & $1$ \\
				 \hline
			\end{tabular}
		\end{center}
		Số trung bình \[\overline{x}=\dfrac{2\cdot 6 + 7\cdot 8 + 7\cdot 10 +  3\cdot 12 + 1\cdot 14}{20}=9{,}4.\]
	} 
\end{ex}

% nguồn: Ok_Toan11_TracNghiem_Lan2
\begin{ex}%[1D5H1-3]%[ma:103fdfcc]
	Điểm kiểm tra giữa kỳ I của $1$ lớp được cô  giáo chủ nhiệm ghi lại theo bảng sau
	\begin{center}
		\begin{tabular}{|c|c|c|c|c|c|}
			\hline 
			Điểm & $[4{,}5;5{,}5)$ & $[5{,}5;6{,}5)$ & $[6{,}5;7{,}5)$ & $[7{,}5;8{,}5)$ &  $[8{,}5;9{,}5)$\\
			\hline 
			Số học sinh  & $3$ & $7$ & $9$ & $15$ & $6$\\
			\hline 
		\end{tabular}
	\end{center}
	Điểm trung bình giữa kỳ  I của lớp là
	 \choice
		{\True  $7{,}35$}
		{$7{,}45$} 
		 {$8{,}35$}
		{$6{,}45$} 
	\loigiai{ 
		\begin{center}
			\begin{tabular}{|c|c|c|c|c|c|}
				\hline 
				Điểm & $[4{,}5;5{,}5)$ & $[5{,}5;6{,}5)$ & $[6{,}5;7{,}5)$ &  $[7{,}5;8{,}5)$ & $[8{,}5;9{,}5)$\\
				 \hline
				Giá trị đại diện & $5$ &  $6$ & $7$ & $8$ & $9$\\
				 \hline
				Số học sinh & $3$  & $7$ & $9$ & $15$ & $6$\\
				 \hline
			\end{tabular}
		\end{center}
		Điểm trung bình của  lớp
		\[\dfrac{3\cdot5+7\cdot6+9\cdot7+15\cdot8+6\cdot9}{40}=7{,}35 
		\]. 
	 }
\end{ex}

% nguồn: Ok_Toan11_TracNghiem_Lan2
\begin{ex}%[1D5H1-3]%[ma:0a5003a6]
	Kiểm tra cân nặng $25$ quả dưa của hai nhà vườn $A$ và $B$  cho kết quả như bảng sau
	\begin{center}
		\begin{tabular}{|c|c|c|c|c|c|}
			 \hline
			Cân nặng (kg) & $[1{,}5;2{,}5)$ &  $[2{,}5;3{,}5)$ & $[3{,}5;4{,}5)$ & $[4{,}5;5{,}5)$ & $[5{,}5;6{,}5)$\\
			 \hline
			Số quả  của vườn $A$ &$2$ & $6$ & $4$ & $10$ & $3$\\
			 \hline
			 Số quả của vườn $B$ & $3$ & $7$ & $8$ &$5$ & $2$\\
			 \hline
		\end{tabular}
	\end{center}
	Cân nặng trung bình của mỗi quả dưa ở  vườn $A$ là
	 \choice
		 {$5{,}24$ kg}
		 {\True $4{,}24$ kg}
		{$3{,}24$  kg}
		 {$5{,}24$ kg}
	\loigiai{ 
		\begin{center}
			\begin{tabular}{|c|c|c|c|c|c|}
				 \hline
				Cân nặng(kg) & $[1{,}5;2{,}5)$ & $[2{,}5;3{,}5)$ & $[3{,}5;4{,}5)$  & $[4{,}5;5{,}5)$ & $[5{,}5;6{,}5)$\\
				 \hline
				Giá trị đại diện & $2$ & $3$ & $4$ &  $5$ & $6$\\
				 \hline
				 Số quả của vườn $A$ & $2$ & $6$ & $4$ & $10$ & $3$\\
				 \hline
				Số quả của vườn $B$ & $3$  & $7$ & $8$ & $5$ & $2$\\
				 \hline
			\end{tabular}
		\end{center}
		Trung bình mỗi quả dưa của vườn $A$ nặng số kg 
		\begin{center}
			$\dfrac{2\cdot2+6\cdot 3+4\cdot  4+10\cdot 5+3\cdot 6}{25}=4{,}24$ (kg).
		\end{center}
	} 
\end{ex}

% nguồn: Ok_Toan11_TracNghiem_Lan2
\begin{ex}%[1D5H1-3]%[ma:d8698297]
	Dưới đây  là một mẫu số liệu cho ở dạng bảng tần số ghép nhóm
	\begin{center}
		\begin{tabular}{|c|c|c|c|c|c|}
			 \hline
			Nhóm &  $[0;50)$ & $[50;100)$ & $[100;150)$ & $[150;200)$ &\\
			 \hline
			Tần số & $6$ & $8$  & $7$ & $6$ & $n=27$\\
			 \hline
		\end{tabular}
	\end{center}
	Số trung bình của  mẫu số liệu là một số thỏa mãn điều kiện
	 \choice
		 {$\overline{x}\approx 112{,}8$}
		{$\overline{x}\approx  107{,}8$}
		 {\True $\overline{x}\approx 99{,}1$}
		{$\overline{x}\approx  85{,}5$}
	 \loigiai{
		\[\overline{x}=\dfrac{25\cdot6+75\cdot8+125\cdot7+175\cdot6}{27}\approx  99{,}1.
		 \]
	 }
\end{ex}

% nguồn: Ok_Toan11_TracNghiem_Lan2
\begin{ex}%[1D5H1-3]%[ma:831976a0]
	Thống kê cân nặng (đơn vị kg) của học  sinh lớp $11D$ cho trong bảng dưới đây
	\begin{center}
		\begin{tabular}{|c|c|c|c|c|c|c|}
			 \hline
			Cân nặng & $[40{,}5;45{,}5)$ & $[45{,}5;50{,}5)$ & $[50{,}5;55{,}5)$  & $[55{,}5;60{,}5)$ & $[60{,}5;65{,}5)$ & $[65{,}5;70{,}5)$ \\
			\hline 
			Số học  sinh & $10$ & $7$ & $16$ & $4$ & $2$ & $3$ \\
			\hline 
		\end{tabular}
	\end{center}
	Tính cân nặng trung  bình của học sinh lớp $11D$.
	 \choice
		{\True  $51{,}81$}
		 {$50{,}1$}
		 {$52{,}83$}
		{$50{,}81$} 
	 \loigiai{
		Cân nặng trung bình  của học sinh lớp $11D$ là
		\[\overline{x}=\dfrac{10\cdot 43 + 7\cdot 48 + 16\cdot  53 + 4\cdot 58 + 2\cdot 63 + 3\cdot 68}{42} \approx 51{,}81 \text{ kg}.\]
	 }
\end{ex}

% nguồn: Ok_Toan11_TracNghiem_Lan2
\begin{ex}%[1D5H1-3]%[ma:7556a34d]
	Khi thống kê chiều cao  của $40$ bạn lớp $11A$, ta thu được mẫu số liệu ghép nhóm được cho ở bảng sau
	\begin{center}
		\begin{tabular}{|c|c|c|c|c|}
			 \hline
			Chiều cao & $[150;160)$  & $[160;170)$ & $[170;180)$ & $[180;190)$ \\
			\hline 
			Số bạn  & $5$ & $12$ & $16$ & $7$ \\
			\hline 
		\end{tabular}
	\end{center}
	Số trung bình cộng  bằng
	 \choice
		{\True $165{,}6$} 
		 {$156{,}6$}
		{$155{,}6$} 
		{$156{,}5$} 
	\loigiai{ 
		Số trung bình  cộng là
		\[\overline{x}=\dfrac{5\cdot 155 + 12\cdot 165 + 16\cdot 175 + 7\cdot 185}{40} \approx 165{,}6.\] 
	} 
\end{ex}

% nguồn: Ok_Toan11_TracNghiem_Lan2
\begin{ex}%[1D5H1-3]%[ma:071ba6bf]
	Khảo sát thời gian tập thể dục của một số học sinh  khối 11 thu được mẫu số liệu ghép nhóm sau
	\begin{longtable}{|c|c|c|c|c|c|}
		 \hline
		Thời gian (phút) & $[0; 20)$ & $[20; 40)$ & $[40; 60)$ & $[60; 80)$ & $[80;  100)$ \\
		 \hline
		Số học sinh & $5$ & $9$  & $12$ & $10$ & $6$ \\
		\hline 
	\end{longtable}
	Nhóm chứa trung vị của  mẫu số liệu trên là
	 \choice
		{\True  $\left[40; 60\right)$}
		{$\left[20;  40\right)$}
		 {$\left[60; 80\right)$}
		{$\left[80;  100\right)$}
	 \loigiai{
		Ta có $n=42$ nên trung vị của mẫu số  liệu trên là $Q_2=\dfrac{x_{21}+x_{22}}{2}$.\\
		Mà $x_{21}, x_{22} \in \left[40; 60\right)$.\\ 
		Vậy nhóm chứa trung vị của  mẫu số liệu trên là nhóm $[40; 60)$.
	 }
\end{ex}

% nguồn: Ok_TracNghiemLop11_Lan1
\begin{ex}%[1D5H1-3]%[ma:36f63384]
	Cho khảo sát thời gian tập  thể dục (đơn vị: phút) của một số học sinh khối $11$, thu được mẫu số liệu ghép nhóm sau:
	\begin{center}
		\begin{tabular}{|c|c|c|c|c|c|}
			\hline 
			Thời gian  & {$[0; 20)$} & {$[20; 40)$} & {$[40; 60)$} & {$[60; 80)$} & {$[80; 100)$} \\
			\hline 
			Số học sinh & $5$& $9$& $12$& $10$&  $6$\\
			 \hline
		\end{tabular}
	\end{center}
	Thời gian tập thể dục trung bình của học sinh khối $11$ trong mẫu số liệu gần nhất với số  nào sau đây?
	 \choice
		 {$50$ phút}
		 {$52$ phút}
		{\True  $51{,}43$ phút}
		{$51$ phút} 
	 \loigiai{
		\begin{center}
			\begin{tabular}{|c|c|c|c|c|c|}
				 \hline
				 Thời gian & {$[0; 20)$} & {$[20; 40)$} & {$[40; 60)$} & {$[60; 80)$} & {$[80; 100)$} \\
				\hline 
				Trung điểm $\left(x_i\right)$ & $10$  & $30$ & $50$ & $70$ & $90$\\
				 \hline
				Số học sinh $\left(f_i\right)$ & $5$& $9$&  $12$& $10$& $6$\\
				\hline 
			\end{tabular}
		\end{center}
		Thời gian tập  thể dục trung bình của học sinh khối $11$ là
		\[\overline{x}=\dfrac{\sum_{i=1}^nf_ix_i}{\sum_{i=1}^{n}f_i}=\dfrac{5\cdot 10+9\cdot30+12\cdot50+10\cdot70+6\cdot90}{5+9+12+10+6}=\dfrac{2160}{42}\approx 51{,}43 \text{ (phút)}.\] 
	} 
\end{ex}

% nguồn: Ok_TracNghiemLop11_Lan1
\begin{ex}%[1D5H1-3]%[ma:77cd1061]
	$100$ người thực hiện bài trắc  nghiệm để đo chỉ số $IQ$, kết quả thu được như sau:
	\begin{center}
		\begin{tabular}{|c|c|c|c|c|c|c|}
			\hline 
			Chỉ số $IQ$ & Dưới  $70$ & $\left[70;85\right)$ & $\left[85;115\right)$ & $\left[115;130\right)$ & $\left[130;145\right)$ & Từ $145$ trở lên \\
			 \hline
			Số người & $2$  & $15$ & $45$ & $20$ & $15$ & $3$ \\
			\hline 
		\end{tabular}
	\end{center}
	Người có chỉ số $IQ$ từ $85$ đến dưới $115$ là ở mức trung bình. Xác định tỉ lệ người có $IQ$ cao hơn  mức trung bình.
	\choice 
		{$41\%$} 
		 {$35\%$}
		 {$40\%$}
		{\True  $38\%$}
	\loigiai{ 
		\begin{listEX}
			\item Các nhóm số liệu  gồm Dưới $70$; $\left[70;85\right); \left[85;115\right); \left[115;130\right); \left[130;145\right)$;\\
			Từ $145$ trở lên với các tần số tương ứng là  $2, 15, 45, 20, 15, 3$.
			\item Số người có chỉ  số $IQ$ cao hơn mức trung bình là $20+15+3=38$.\\
			Vậy tỉ lệ người có chỉ số $IQ$ cao hơn mức trung bình là  $\dfrac{38}{100}=38\%$.
	\end{listEX}}
\end{ex}

% nguồn: Ok_TracNghiemLop11_Lan1
\begin{ex}%[1D5H1-3]%[ma:2d957a93]
	Quãng đường (km) các cầu thủ (không tính thủ môn) chạy trong một trận bóng đá tại giải ngoại hạng Anh được cho trong bảng thống kê  sau:
	\begin{center}
		\begin{tabular}{|c|c|c|c|c|c|}
			\hline Quãng đường & {$[2; 4)$} & {$[4; 6)$} & {$[6; 8)$}  & {$[8; 10)$} & {$[10; 12)$} \\
			\hline Số cầu thủ & 2 & 5 & 6  & 9 & 3 \\
			\hline 
		\end{tabular}
	\end{center}
	Tính quãng đường  trung bình một cầu thủ chạy trong trận đấu này.
	 \choice
		 {$7{,}25$}
		{$8{,}23$} 
		 {\True $7{,}48$}
		 {$7{,}67$}
	 \loigiai{
		Tổng  số cầu thủ là $n=2+5+6+9+3=25$.\\
		Quãng đường trung bình một  cầu thủ chạy trong trận đấu này là
		$$\overline{x}=\dfrac{2\cdot 3+5\cdot 5+6\cdot  7+9\cdot 9+3\cdot 11}{25}=7{,}48~(\mathrm{km}).$$
	} 
\end{ex}

% nguồn: Ok_TracNghiemLop11_Lan1
\begin{ex}%[1D5H1-3]%[ma:676ffe01]
	Cho bảng số liệu ghép nhóm về chiều cao  cm của học sinh của một lớp như sau
	\begin{center}
		\begin{tabular}{|c|c|c|c|c|c|}
			\hline Chiều cao $(\mathrm{cm})$ & {$[160; 165)$} & {$[165; 170)$} & {$[170; 175)$} & {$[175; 180)$} & {$[180; 185)$}  \\
			\hline Số học sinh & $13$ & $12$ & $10$ & $7$ & $3$ \\ 
			\hline 
		\end{tabular}
	\end{center}
	Tính chiều cao  trung bình của học sinh lớp trên (kết quả làm tròn đến hàng phần trăm).
	\choice 
		{$169{,}73$} 
		 {$169{,}27$}
		{$168{,}37$} 
		 {\True $169{,}72$}
	\loigiai{ 
		Chiều cao trung bình  của học sinh lớp trên là
		$$\dfrac{162{,}5\cdot13+167{,}5\cdot12+172{,}5\cdot10+177{,}5\cdot7+182{,}5\cdot3}{13+12+10+7+3}=\dfrac{3055}{18}\approx 169{,}72.$$ 
	 }
\end{ex}

% nguồn: Ok_TracNghiemLop11_Lan1
\begin{ex}%[1D5H1-3]%[ma:bb81bd85]
	Dưới đây là một mẫu số liệu cho ở dạng bảng tần số ghép  nhóm
	\begin{center}
		\begin{tabular}{|c|c|c|c|c|c|}
			 \hline
			Nhóm &  $[0;50)$ & $[50;100)$ & $[100;150)$ & $[150;200)$ & \\
			 \hline
			 Tần số & $6$ & $8$ & $7$ & $6$ & $N=27$ \\
			\hline 
		\end{tabular}
	\end{center}
	Số trung bình của mẫu  số liệu là (kết quả làm tròn đến hàng phần chục).
	 \choice
		{$\overline{x}\approx 112{,}8$} 
		{$\overline{x}\approx  107{,}8$}
		{\True $\overline{x}\approx 99{,}1$} 
		 {$\overline{x}\approx 85{,}5$}
	\loigiai{ 
		Ta có $\overline{x}=\dfrac{25\cdot 6+75\cdot 8+125\cdot 7+175\cdot 6}{27}\approx  99{,}1$.
	} 
\end{ex}

% nguồn: Ok_TracNghiemLop11_Lan1
\begin{ex}%[1D5H1-3]%[ma:bad0b02b]
	Cho mẫu số liệu ghép nhóm về thống kê nhiệt độ tại một  địa điểm trong $30$ ngày, ta có bảng số liệu sau:
	\begin{center}
		\begin{tabular}{|l|c|c|c|c|}
			\hline Nhiệt  độ $\left({ }^{\circ} \mathrm{C}\right)$ & {$[18; 21)$} & {$[21; 24)$} & {$[24; 27)$} & {$[27; 30)$}\\
			\hline Số ngày & $6$ & $12$ & $9$ & $3$  \\
			 \hline
		\end{tabular}
	\end{center}
	Nhiệt  độ trung bình trong $30$ ngày trên là
	 \choice
		{$24^{\circ}  \mathrm{C}$}
		 {$25{,}4^{\circ} \mathrm{C}$}
		{$24{,}3^{\circ}  \mathrm{C}$}
		{\True $23{,}4^{\circ} \mathrm{C}$} 
	\loigiai{ 
		Giá trị đại diện của mỗi nhóm số liệu là  trung bình cộng của hai đầu mút.\\
		Ta có bảng tần số ghép nhóm  theo giá trị đại diện của mỗi nhóm:
		\begin{center}
			\begin{tabular}{|l|c|c|c|c|}
				\hline Nhóm & {$[18; 21)$} & {$[21; 24)$} & {$[24; 27)$}  & {$[27; 30)$}\\
				\hline Giá trị đại diện &  $19{,}5$ & $22{,}5$ & $25{,}5$ & $28{,}5$ \\
				\hline Tần số & $6$ & $12$ & $9$ &  $3$ \\
				\hline 
			\end{tabular}
		\end{center}
		 Nhiệt độ trung bình trong 30 ngày trên là
		$$\overline{x}=\dfrac{6 \cdot 19{,}5+12 \cdot 22{,}5+9\cdot25{,}5+3 \cdot  28{,}5}{30}=23{,}4^{\circ} \mathrm{C}.$$
	 }
\end{ex}

% nguồn: Ok_TracNghiemLop11_Lan1
\begin{ex}%[1D5H1-3]%[ma:00b8b5ea]
	Điểm trung bình các môn học kì $I$ của bạn An được cho bởi bảng sau 
	\begin{center}
		\begin{tabular}{|c|c|c|c|c|c|c|c|c|}
			\hline Môn & Toán  & Vật Lý & Hóa học & Ngữ văn & Lịch sử & Địa lý & Tin học & Tiếng Anh \\
			\hline Điểm & $9{,}4$ & $8{,}8$ & $9{,}2$ & $6{,}8$ & $8{,}0$ & $7{,}8$ & $8{,}4$ &  8{,}6 \\
			 \hline
		\end{tabular}
	\end{center}
	 Tính điểm trung bình môn học kì $I$ của bạn An.
	 \choice
		{$8{,}0$} 
		{\True  ${8{,}375}$}
		{${8{,}2}$} 
		 {$8{,}5$}
	\loigiai{ 
		Điểm trung bình môn học kì I của bạn  An là
		\[\dfrac{9{,}4+8{,}8+9{,}2+6{,}8+8{,}0+7{,}8+8{,}4+8{,}6}{8}=8{,}375.\] 
	} 
\end{ex}

% nguồn: Ok_TracNghiemLop11_Lan1
\begin{ex}%[1D5H1-3]%[ma:bcf030f2]
	Các bạn học sinh lớp $11$A$1$ trả lời $40$ câu hỏi trong một bải kiểm tra. Kết quả được thống kê ở bảng sau. Hãy ước lượng trung bình số câu trả lời đúng  của các học sinh lớp $11$A$1$.
	\begin{center}
		\begin{tabular}{|c|c|c|c|c|c|c|}
			\hline 
			Số câu trả lời đúng & $\left[16;21\right)$  & $\left[21;26\right)$ & $\left[26;31\right)$ & $\left[31;36\right)$ & $\left[36;41\right)$ \\
			 \hline
			Số học  sinh & $4$ & $6$ & $8$ & $18$ & $4$ \\
			 \hline
		\end{tabular}
	\end{center}
	\choice 
		{\True $30$} 
		 {$32$}
		{$29$} 
		{$31$} 
	\loigiai{ 
		Trong mỗi khoảng, giá trị đại diện là trung bình cộng của giá trị hai đầu mút nên  ta có bảng sau
		\begin{center}
			\begin{tabular}{|c|c|c|c|c|c|c|}
				\hline 
				Số câu trả lời đúng & $\left[16;21\right)$ & $\left[21;26\right)$ &  $\left[26;31\right)$ & $\left[31;36\right)$ & $\left[36;41\right)$ \\
				 \hline
				Số  học sinh & $4$ & $6$ & $8$ & $18$ & $4$ \\
				\hline 
				Giá trị đại diện & $18{,}5$ & $23{,}5$ & $28{,}5$ & $33{,}5$ & $38{,}5$  \\
				 \hline
			\end{tabular}
		\end{center}
		Trung bình số câu trả lời đúng của các học sinh lớp  11A1 là
		$$\bar{x}=\dfrac{4\cdot 18{,}5+6\cdot  23{,}5+8\cdot 28{,}5+18\cdot 33{,}5+4\cdot 38{,}5}{4+6+8+18+4}=30.$$
	 }
\end{ex}

% nguồn: Ok_TracNghiemLop11_Lan1
\begin{ex}%[1D5H1-3]%[ma:7ed45f43]
	Số trung bình của mẫu số liệu trên thuộc khoảng nào trong các khoảng dưới  đây?
	\choice 
		 {$[7;9)$}
		{\True $[9;11)$} 
		 {$[11;13)$}
		{$[13;15)$} 
	\loigiai{ 
		Bảng tần số  ghép nhóm theo giá trị đại diện là
		\begin{center}
			\begin{tabular}{|c|c|c|c|c|c|}
				\hline Doanh thu &{$[5; 7)$} &{$[7; 9)$} &{$[9; 11)$} &{$[11; 13)$} &{$[13;  15)$} \\
				\hline Giá trị đại  diện &$6$ &$8$ &$10$ &$12$ &$14$ \\
				\hline Số ngày & 2 & 7 & 7 &  3 & 1 \\
				\hline 
			\end{tabular}
		\end{center}
		Số trung bình: $\overline{x}=\dfrac{2\cdot 6+7\cdot 8+7\cdot 10+3\cdot 12+1\cdot 14}{20}=9{,}4$. 
	} 
\end{ex}

% nguồn: Ok_TracNghiemLop11_Lan1
\begin{ex}%[1D5H1-3]%[ma:7a0a6dfc]
	Một bưu tá thống kê lại số bưu phẩm gửi đến một cơ  quan mỗi ngày trong tháng 6/2022 trong bảng sau:
	\begin{center}
		\begin{tabular}{|l|c|c|c|c|c|}
			\hline Số bưu phẩm & $[20; 24]$ & $[25; 29]$ & $[30; 34]$ & $[35; 39]$ & $[40;  44]$ \\
			\hline Số ngày & $4$ & $6$ & $10$ &  $6$ & $4$ \\
			 \hline
		\end{tabular}
	\end{center}
	Số trung bình của mấu số liệu  là
	 \choice
		 {$30$}
		 {$31$}
		 {$30$}
		{\True $32$} 
	\loigiai{ 
		Do số bưu phẩm là số nguyên nên ta  hiệu chỉnh lại
		\begin{center}
			\begin{tabular}{|l|c|c|c|c|c|}
				 \hline
				Số bưu  phẩm & $[19{,}5; 24{,}5)$ & $[24{,}5; 29{,}5)$ & $[29{,}5; 34{,}5)$ & $[34{,}5; 39{,}5)$ & $[39{,}5; 44{,}5)$ \\
				\hline 
				Giá  trị đại diện & $22$ & $27$ & $32$ & $37$ & $42$ \\
				 \hline
				Số ngày & $4$ & $6$  & $10$ & $6$ & $4$ \\
				\hline 
			\end{tabular}
		\end{center}
		Số trung bình của  mẫu số liệu ghép nhóm là
		 $\overline{x}=\dfrac{4\cdot 22+6\cdot 27+10\cdot 32+6\cdot 37+4\cdot 42}{30}{=~32}$.
	} 
\end{ex}

% nguồn: Ok_TracNghiemLop11_Lan1
\begin{ex}%[1D5H1-3]%[ma:d00a906f]
	Nồng độ cồn trong hơi thở (đơn vị tính là miligam/1 lít khí thở) của $20$ lái xe ô tô vi phạm được cho như sau 
	\begin{center}
		\begin{longtable}{|c|c|c|c|c|c|c|c|c|c|}
			 \hline
			$0{,}09$ & $0{,}18$ & $0{,}47$ & $1{,}20$ & $0{,}28$ & $0{,}45$ & $0{,}72$ & $0{,}15$ &  $0{,}75$ & $0{,}36$\\
			 \hline
			$0{,}21$ & $0{,}15$ & $0{,}23$ & $0{,}30$ & $0{,}41$ & $0{,}13$ & $0{,}05$ & $0{,}38$ & $0{,}42$ &  $0{,}79$\\
			\hline 
		\end{longtable}
	\end{center}
	 Theo quy định, mức phạt nồng độ cồn đối với lái xe ô tô như sau:
	\begin{description}
		 \item[Mức 1.] Nồng độ cồn trong hơi thở chưa vượt quá $0{,}25$ phạt từ $6$ đến $8$ triệu đồng;
		\item[Mức 2.] Nồng độ cồn trong hơi  thở từ trên $0{,}25$ đến $0{,}4$ phạt từ $16$ đến $18$ triệu đồng;
		\item[Mức 3.] Nồng độ cồn trong hơi thở  vượt quá $0{,}4$ phạt từ $30$ đến $40$ triệu đồng.
	\end{description}
	Trung bình mỗi lái xe bị phạt  bao nhiêu tiền?
	 \choice
		{$23{,}5$} 
		 {$22{,}1$}
		 {$24{,}4$}
		{\True  $20{,}2$}
	 \loigiai{
		Ta  có bảng số liệu ghép nhóm như sau
		\begin{center}
			\begin{longtable}{|c|c|c|c|}
				 \hline
				Tiền phạt (triệu đồng) & $6-8$ &  $16-18$ & $30-40$\\
				\hline 
				 Người vi phạm & $8$ & $4$ & $8$\\
				 \hline
			\end{longtable}
		\end{center}
		Cỡ mẫu $n=20$. Số tiền trung bình  một người bị phạt là
		$$\overline{x}=\dfrac{8\cdot 7+4\cdot 17+8\cdot 35}{20}=20{,}2\text{ (triệu đồng).}$$ 
	} 
\end{ex}

% nguồn: Ok_TracNghiemLop11_Lan1
\begin{ex}%[1D5H1-3]%[ma:29dd4513]
	Điều tra về chiều cao của  học sinh khối lớp $10$, ta có kết quả sau
	\begin{center}
		\begin{tabular}{|c|c|c|}
			\hline Nhóm & Chiều cao (cm) & Số học  sinh \\
			\hline 1  & {$[150; 152)$} & $10$ \\
			\hline 2 & {$[152; 154)$} & $12$ \\ 
			\hline 3 &  {$[154; 156)$} & $35$ \\
			\hline 4 & {$[156; 158)$} &  $20$ \\
			\hline 5 & {$[158; 160)$} &  $10$ \\
			\hline  6 & {$[160; 162)$} & $13$ \\
			\hline \multicolumn{2}{|c|}{} & $\mathrm{N}=100$ \\ 
			 \hline
		\end{tabular}
	\end{center}
	Số trung bình là 
	 \choice
		 {$154{,}94$}
		{$159{,}54$} 
		 {$152{,}45$}
		{\True $155{,}94$} 
	\loigiai{ 
		Số trung  bình là
		$$\dfrac{151\cdot 10+153\cdot 12+155\cdot 35+157\cdot  20+159\cdot 10+161\cdot 13}{100}=155{,}94.$$
	 }
\end{ex}

% nguồn: Ok_TracNghiemLop11_Lan1
\begin{ex}%[1D5H1-3]%[ma:9a63f2d1]
	 Số nguyện vọng đăng kí vào đại học của các bạn trong lớp được thống kê trong bảng sau
	\begin{center}
		\begin{longtable}{|c|c|c|c|c|}
			 \hline
			Số nguyện vọng & $1-3$ &  $4-6$ & $7-9$ & $10-12$\\
			\hline 
			Số học sinh & $5$  & $18$ & $13$ & $7$\\
			 \hline
		\end{longtable}
	\end{center}
	Trung bình một bạn trong lớp đăng kí bao nhiêu  nguyện vọng.
	\choice 
		{$7{,}45$} 
		{$6{,}67$} 
		{$7{,}42$} 
		{\True $6{,}53$} 
	\loigiai{ 
		 Số trung bình của mẫu số liệu là
		$$\overline{x}=\dfrac{5\cdot 2+18\cdot 5+13\cdot 8+7\cdot 11}{43}\approx 6{,}53.$$ 
	 }
\end{ex}

% nguồn: Ok_TracNghiemLop11_Lan1
\begin{ex}%[1D5H1-3]%[ma:69537eae]
	Diện tích các tỉnh và thành phố khu vực Nam Bộ được thống  kê ở bảng sau
	\begin{center}
		\begin{tabular}{|c|c|c|c|}
			 \hline Tỉnh/ thành phố & Diện tích $\left(\mathbf{k m}^2\right)$ & Tỉnh/ thành phố & Diện tích $\left.\mathbf{( k m}^{\mathbf{2}}\right)$ \\
			\hline Bình Phước & 6877 & Vĩnh Long &  1526 \\
			\hline Tây Ninh & 4041 & Đồng Tháp & 3384  \\
			\hline Bình Dương & 2695 &  An Giang & 3537 \\
			\hline Đồng Nai & 5864 & Kiên Giang & 6349  \\
			\hline Bà Rịa - Vũng Tàu & 1981 & Cần Thơ  & 1439 \\
			\hline TP.Hồ Chí Minh &  2061 & Hậu Giang & 1622 \\
			\hline Long An & 4495 &  Sóc Trăng & 3312 \\
			\hline Tiền Giang & 2511  & Bạc Liêu & 2669 \\
			\hline Bến Tre &  2395 & Cà Mau & 5221 \\
			\hline Trà Vinh &  2358 & & \\
			\hline 
		\end{tabular}
		 \newline
		(Nguồn: Tổng  cục Thống kê)
	\end{center}
	Hãy ước lượng số trung bình của mẫu số liệu ghép nhóm  trên (làm tròn đến hàng phần trăm).
	 \choice
		{\True $3407{,}89$} 
		 {$3321{,}35$}
		{$3514{,}26$} 
		{$3612{,}15$} 
	\loigiai{ 
		Bảng tần số ghép  nhóm bao gồm giá trị đại diện của các nhóm như sau
		\begin{center}
			\begin{tabular}{|l|c|c|c|c|}
				\hline Diện tích $\left(\mathrm{km}^2\right)$ & {$[1000; 2500)$} & {$[2500; 4000)$} &  {$[4000; 5500)$} & {$[5500; 7000)$} \\
				\hline Giá trị đại diện & 1750 &  3250 & 4750 & 6250 \\
				\hline Số tỉnh/ thành phố & 7 &  6 & 3 & 3 \\
				\hline 
			\end{tabular}
		\end{center}
		Khi đó, số trung bình của mẫu số liệu ghép nhóm là 
		$$\bar{x}=\dfrac{7\cdot 1750+6 \cdot  3250+3 \cdot 4750+3 \cdot 6250}{19} \approx 3407{,}89.$$
	 }
\end{ex}

% nguồn: Ok_TracNghiemLop11_Lan1
\begin{ex}%[1D5H1-3]%[ma:d4576baa]
	Doanh thu bán hàng trong $20$ ngày được lựa chọn  ngẫu nhiên của một của hàng được ghi lại ở bảng sau (đơn vị: triệu đồng).
	\begin{center}
		\begin{tabular}{|l|c|c|c|c|c|}
			\hline Doanh thu & {$[5; 7)$} & {$[7; 9)$} & {$[9; 11)$}  & {$[11; 13)$} & {$[13; 15)$} \\
			\hline Số ngày &  $2$ & $7$ & $7$ & $3$ & $1$ \\
			 \hline
		\end{tabular}
	\end{center}
	Số trung bình của mẫu số liệu trên thuộc khoảng nào trong các khoảng dưới đây? 
	 \choice
		{$[7;  9)$}
		{\True $[9; 11)$} 
		{$[11;  13)$}
		{$[13;  15)$}
	 \loigiai{
		Bảng tần số ghép nhóm theo  giá trị đại diện là
		\begin{center}
			\begin{tabular}{|l|c|c|c|c|c|}
				\hline Doanh thu & {$[5;  7)$} & {$[7; 9)$} & {$[9; 11)$} & {$[11; 13)$} & {$[13; 15)$} \\
				\hline Giá trị đại diện  & $6$ & $8$ & $10$ & $12$ & $14$ \\
				\hline Số ngày & $2$ & $7$  & $7$ & $3$ & $1$ \\
				\hline 
			\end{tabular}
		\end{center}
		Số  trung bình là $\overline{x}=\dfrac{2\cdot6+7\cdot8+7\cdot10+3\cdot12+1\cdot14}{20}=9{,}4$.
	 }
\end{ex}

% nguồn: Ok_TracNghiemLop11_Lan1
\begin{ex}%[1D5H1-3]%[ma:25bf1cee]
	Cho  mẫy số liệu ghép nhóm về tuổi thọ (đơn vị tính là năm) của một loại bóng đèn mới như sau.
	\begin{center}
		\begin{longtable}{|c|c|c|c|c|}
			\hline 
			Tuổi thọ  (năm) & $[2;3{,}5)$ & $[3{,}5;5)$ & $[5;6{,}5)$ & $[6{,}5;8)$\\
			\hline 
			Số bóng đèn &  $8$ & $22$ & $35$ & $15$\\
			\hline 
		\end{longtable}
	\end{center}
	Số  trung bình của mẫu số liệu là
	 \choice
		 {$5{,}0$}
		 {\True $5{,}32$}
		{$5{,}75$} 
		 {$6{,}5$}
	\loigiai{ 
		Ta có  bảng sau
		\begin{center}
			\begin{longtable}{|c|c|c|c|c|}
				\hline 
				Tuổi thọ (năm) & $[2;3{,}5)$ & $[3{,}5;5)$ &  $[5;6{,}5)$ & $[6{,}5;8)$\\
				\hline 
				Giá trị đại  diện & $2{,}75$ & $4{,}25$ & $5{,}75$ & $7{,}25$\\
				Số bóng đèn & $8$ & $22$  & $35$ & $15$\\
				 \hline
			\end{longtable}
		\end{center}
		Tổng số bóng đèn  là $n=80$.\\
		Số trung bình của mẫu số liệu  là
		 $$\overline{x}=\dfrac{2{,}75\cdot 8+4{,}25\cdot 22+5{,}75\cdot 35+7{,}25\cdot 15}{80}=5{,}31875\approx 5{,}32.$$
	} 
\end{ex}

% nguồn: Ok_TracNghiemLop11_Lan1
\begin{ex}%[1D5V1-3]%[ma:7a80bda9]
	Một của hàng bán $3$ loại hoa quả nhập khẩu: Nho Mỹ, Lê Hàn Quốc và Táo New Zealand. Sau khi giảm giá mỗi loại lần lượt  là $x$, $y$, $z$ trên $1kg$ thì số liệu tính toán được ghi lại bởi bảng sau
	\begin{center}
		\begin{tabular}{|c|c|c|c|}
			\hline Loại quả & Nho Mỹ & Lê Hàn Quốc & Táo  New Zealand \\
			\hline Giá bán (nghìn $/ \mathrm{kg}$)  & $250-x$ & $200-y$ & $180-z$ \\
			\hline Số  lượng bán ( $\mathrm{kg}$) & $250+x$ & $200+y$ & $180+z$ \\
			\hline 
		\end{tabular}
	\end{center}
	Biết rằng $x+y+z=120$. Tính giá trị $x$, $y$, $z$ để  lợi nhuận bình quân của $1kg$ hoa quả đạt được cao nhất.
	 \choice
		{\True  $x=y=z=40$}
		{$x=50$; $y=30$;  $z=40$}
		{$x=30$; $y=50$;  $z=40$}
		 {$x=20$; $y=60$; $z=40$}
	 \loigiai{
		Vì số lượng hoa quả bán được là $250+x+200+y+180+z=750$ là cố định nên bình  quân mỗi $kg$ hoa quả có giá cao nhất khi số tiền thu được là cao nhất.\\
		Gọi $P$ là  tổng số tiền thu được.\\
		Khi  đó
		\begin{eqnarray*}
			 P&=&\left(250-x\right)(250+x)+(200-y)(200+y)+(180-z)(180+z)\\
			&=&250^2-x^2+200^2-y^2+180^2-z^2=134900-x^2-y^2-z^2. 
		\end{eqnarray*}
		Ta có bất đẳng  thức $x^2+y^2+z^2\ge \dfrac{1}{3} \left(x+y+z\right)^2=4800$.\\
		 Do đó $P\le 130100$.\\
		Vậy $P$ lớn nhất  $\Leftrightarrow \heva{&{x+y+z=120} \\
		&{x=y=z}}\Leftrightarrow  x=y=z=40$.
	} 
\end{ex}


%%%%%%%%%%%%%%%% ĐÚNG/SAI %%%%%%%%%%%%%%%%

% nguồn: Ok_dungsai_lop11_Lan1
\begin{ex}%[1D5N1-3]%[ma:0a7420cf]
 Số khách hàng mua bảo  hiểm ở từng độ tuổi được thông kê như sau
 \begin{center}
   \renewcommand{\arraystretch}{1.5}
  \begin{tabular}{|c|c|c|c|c|c|}
   \hline 
   Độ tuổi & {$[20; 30)$} & {$[30; 40)$} & {$[40; 50)$} & {$[50; 60)$} &  {$[60; 70)$}\\
   \hline 
   Số khách hàng & 3 & 9 & 6 &  4 & 2 \\
    \hline
  \end{tabular}
 \end{center}
  \choiceTF
  {\True Độ tuổi trung bình của khách hàng (làm tròn kết quả đến hàng phần trăm) là  $42{,}08$}
  {\True Mốt của mẫu số liệu  trên thuộc nhóm $2$}
  {Tổng số  khách hàng được khảo sát là $25$}
  {\True Mẫu số liệu trên có $5$  nhóm}
  \loigiai{
  \begin{center}
    \renewcommand{\arraystretch}{1.5}
   \begin{tabular}{|c|c|c|c|c|c|}
    \hline 
    Độ tuổi & {$[20; 30)$} & {$[30; 40)$} & {$[40; 50)$} & {$[50;  60)$} & {$[60; 70)$}\\
     \hline
    Giá trị đại diện&25&35&45&55&65\\ 
     \hline
    Số khách  hàng & 3 & 9 & 6 & 4 & 2 \\
     \hline
   \end{tabular}
  \end{center}
   \vspace*{0.3cm}
  \begin{itemize}
   \item Độ tuổi trung bình của khách hàng $\bar X=\dfrac{3\cdot 25+9\cdot 35+6\cdot 45+4\cdot 55+2\cdot  65}{24}=42{,}08$.
   \item Một  của mẫu số liệu thuộc nhóm $[30;40)$.
   \item Tổng số khách hàng được khảo sát  là $3+9+6+4+2=24$.
   \item Mẫu số liệu trên có $5$ nhóm. 
  \end{itemize}
  }
\end{ex}

% nguồn: Ok_dungsai_lop11_Lan1
\begin{ex}%[1D5H1-3]%[ma:edd40f0c]
 Số cuộc điện thoại một người  thực hiện trong $30$ ngày được lựa chọn ngẫu nhiên được thống kê trong bảng sau:
 \begin{center}
  \begin{tabular}{|c|c|c|c|c|c|}
    \hline
   Số cuộc gọi & $[2{,}5;5{,}5)$ & $[5{,}5;8{,}5)$ & $[8{,}5;11{,}5)$ & $[11{,}5;14{,}5)$  & $[14{,}5;17{,}5)$ \\
   \hline 
   Số ngày & $5$ & $13$ & $7$ & $3$ & $2$ \\ 
   \hline 
  \end{tabular}
 \end{center}
 \choiceTF 
  {Số cuộc gọi  trung bình mỗi ngày là $8{,}1$}
   {\True Nhóm chứa mốt là $[5{,}5;8{,}5)$}
  {\True Mốt của mẫu số  liệu ghép nhóm là $\approx 7{,}21$}
  {Người đó thực hiện tối đa khoảng $8$ cuộc gọi mỗi  ngày}
 \loigiai{ 
  Ta viết lại bảng tần số ghép nhóm  theo giá trị đại diện là
  \begin{center}
   \begin{tabular}{|c|c|c|c|c|c|}
     \hline
    Số cuộc gọi & $[2{,}5;5{,}5)$ & $[5{,}5;8{,}5)$ &  $[8{,}5;11{,}5)$ & $[11{,}5;14{,}5)$ & $[14{,}5;17{,}5)$ \\
     \hline
    Giá trị đại diện & $4$ &  $7$ & $10$ & $13$ & $16$ \\
     \hline
    Số ngày & $5$ & $13$ & $7$  & $3$ & $2$ \\
    \hline 
   \end{tabular}
  \end{center}
  \begin{itemchoice}
   \itemch  Sai. Số cuộc gọi trung bình mỗi ngày là
   \[ 
   \overline{x} = \dfrac{4 \cdot 5 + 7 \cdot 13 + 10 \cdot 7 + 13 \cdot 3+  16 \cdot 2}{30} = 8{,}4.
    \]
   \itemch Đúng. Nhóm chứa mốt  kaf $[5{,}5;8{,}5)$.
   \itemch Sai. Ta có $u_m  = 5{,}5$; $u_{m+1}=8{,}5$ $\Rightarrow u_{m+1}-u_m = 3$; $n_m = 13$; $n_{m-1}=5$; $n_{m+1}=7$.\\
   Vì  vậy mốt của mẫu số liệu ghép nhóm là
    \[
    M_0 = 5{,}5 + \dfrac{13-5}{(13-5)+(13-7)} \cdot 3 = \dfrac{101}{14} \approx 7{,}21.
   \] 
    \itemch Sai. Người đó thực hiện tối đa $7$ cuộc gọi mỗi ngày.
  \end{itemchoice}
  }
\end{ex}

% nguồn: Ok_dungsai_lop11_Lan1
\begin{ex}%[1D5H1-3]%[ma:9685e841]
 Kết quả khảo sát cân nặng của $25$ quả cam ở mỗi lô hàng $A, B$  được cho ở bảng sau:
 \begin{center}
  \begin{tabular}{|c|c|c|c|c|c|}
   \hline 
   Cân nặng (gam) & $[150; 155)$ & $[155; 160)$ & $[160; 165)$ & $[165; 170)$ &  $[170; 175)$ \\
    \hline
   Số quả cam ở lô hàng $A$ & $2$ & $6$ &  $12$ & $4$ & $1$ \\
   \hline 
   Số quả cam ở lô hàng $B$ & $1$ & $3$  & $7$ & $10$ & $4$ \\
   \hline 
  \end{tabular}
 \end{center}
 Khi đó 
 \choiceTF 
  {\True Giá trị đại  diện nhóm $[150; 155)$ bằng $152{,}5$}
  {Cân nặng trung bình của mỗi quả cam ở lô $A$ là  $163{,}7$ (gam)}
  {Cân nặng trung bình của mỗi quả  cam ở lô $B$ là $162{,}1$ (gam)}
   {\True Theo số trung bình thì cam ở lô hàng $B$ nặng hơn cam ở lô hàng $A$}
 \loigiai{ 
  Bảng thống kê số lượng  cam theo giá trị đại diện:
  \begin{center}
   \begin{tabular}{|c|c|c|c|c|c|}
    \hline 
    Cân nặng đại diện (gam)  & $152{,}5$ & $157{,}5$ & $162{,}5$ & $167{,}5$ & $172{,}5$ \\
     \hline
    Số quả cam ở lô hàng $A$ & $2$ & $6$ & $12$ & $4$ & $1$  \\
     \hline
    Số quả cam ở lô hàng $B$  & $1$ & $3$ & $7$ & $10$ & $4$ \\
     \hline
   \end{tabular}
  \end{center}
  \begin{itemchoice}
    \itemch \textbf{Đúng}.\\
   Giá trị đại diện nhóm $[150; 155)$ bằng  $152{,}5$.
    \itemch \textbf{Sai}.\\
   Cân nặng trung bình của mỗi quả cam ở lô A là  $163{,}7$ (gam).\\
   Cân nặng  trung bình của mỗi quả cam ở lô $A$ là
   \[\overline{x}_{A}=\dfrac{152{,}5 \cdot 2+157{,}5  \cdot 6+162{,}5 \cdot 12+167{,}5 \cdot 4+172{,}5 \cdot 1}{25}=161{,}7 (\text{gam}).\]
   \itemch \textbf{Sai}.\\ 
   Cân nặng trung bình của mỗi quả cam ở lô B là $162{,}1$ (gam).\\ 
   Cân nặng  trung bình của mỗi quả cam ở lô $B$ là
   \[ 
   \overline{x}_{B}=\dfrac{152{,}5 \cdot 1+157{,}5 \cdot 3+162{,}5 \cdot 7+167{,}5 \cdot 10+172{,}5 \cdot 4}{25}=165{,}1 (\text{gam}).\] 
   \itemch \textbf{Đúng}.\\ 
   Theo số trung bình thì cam ở lô hàng $B$ nặng hơn cam  ở lô hàng $A$.\\
    Ta thấy $\overline{x}_{A}<\overline{x}_{B}$.\\
   Vậy nếu so sánh theo  số trung bình thì cam ở lô hàng $B$ nặng hơn cam ở lô hàng $A$.
  \end{itemchoice}
 } 
\end{ex}

% nguồn: Ok_dungsai_lop11_Lan1
\begin{ex}%[1D5H1-3]%[ma:3eb5a66a]
 Thống kê số giờ tự học của học sinh tại một trường THPT trong một tuần, người ta  ghi được kết quả sau (đơn vị: giờ).
 $$\begin{array}{*{35}{l}}
 14 & 12{,}5  & 15 & 16{,}5 & 17 & 14{,}5 & 13 & 15{,}5 & 16{,}5 & 17{,}5\\
 16{,}5 & 18{,}5 & 19  & 20 & 19{{,}}5 & 17 & 16{,}5 & 14 & 18 & 21\\
 15{,}5 & 13{,}5 & 12{,}5 & 14{,}5 & 17{,}5 & 19 & 19{{,}}5 & 20{,}5 & 20 &  17\\
 14{,}5 & 13  & 14{,}5 & 18 & 16 & 15 & 13 & 18{,}5 & 14{,}5 & 12{,}5\\
\end{array}$$
Khi  đó
\choiceTF 
 {Tổng số giờ tự  học của học sinh là $829{,}5$ giờ}
 {\True Tổng hợp kết quả tự học của học sinh vào bảng tần số ghép nhóm  theo mẫu sau:\\
 \centerline{\begin{tabular}{|c|c|c|c|c|c|}
 \hline 
 Giờ tự học &$[12{,}5;14{,}5)$ &$[14{,}5;16{,}5)$ &$[16{,}5;18{,}5)$ &$[18{,}5;20{,}5)$  &$[20{,}5;22{,}5)$\\
 \hline 
 Số học sinh  & $9$ & $13$ & $17$ & $9$ & $2$\\
 \hline 
\end{tabular}}}
{\True Thời gian trung bình học của mỗi học sinh dựa vào bảng tần  số ghép nhóm tìm được ở ý (b) là $16{,}78$}
{Khả năng tự  học trong tuần cao nhất của học sinh là $16{,}28$ giờ}
\loigiai{\begin{itemchoice}
 \itemch Sai. Tổng số giờ tự học của  học sinh là $829{,}5$ giờ.\\
 Tổng số giờ tự học của học sinh  là $819{,}5$ giờ.\\
 Số giờ tự học trung bình của học  sinh $\overline{x}=\dfrac{891{,}5}{50}=16{,}39$ giờ.
 \itemch Đúng. Tổng hợp kết quả tự học của học sinh vào bảng tần số ghép nhóm  theo mẫu sau:\\
 \centerline{\begin{tabular}{|c|c|c|c|c|c|}
 \hline 
  Giờ tự học & $[12{,}5;14{,}5)$ & $[14{,}5;16{,}5)$ & $[16{,}5;18{,}5)$ & $[18{,}5;20{,}5)$ & $[20{,}5;22{,}5)$\\
 \hline 
 Số học sinh & $9$ & $13$ & $17$ & $9$  & $2$\\
  \hline
 \end{tabular}}
 \itemch Đúng. Thời gian trung bình học của  mỗi học sinh dựa vào bảng tần số ghép nhóm tìm được ở ý (b) là $16{,}78.$\\
 \centerline{\begin{tabular}{|c|c|c|c|c|c|}
  \hline
 Giờ tự học & $[12{,}5;14{,}5)$ & $[14{,}5;16{,}5)$ & $[16{,}5;18{,}5)$  & $[18{,}5;20{,}5)$ & $[20{,}5;22{,}5)$\\
 \hline 
 Giá trị đại diện & $13{,}5$ & $15{,}5$ &  $17{,}5$ & $19{,}5$ & $21{,}5$\\
  \hline
 Số học sinh & $9$ &  $13$ & $17$ & $9$ & $2$\\
  \hline
 \end{tabular}}\\
 Thời gian  tự học trung bình của học sinh
  $$\overline{x}=\dfrac{9\cdot13{,}5+13\cdot 15{,}5+17\cdot 17{,}5+9\cdot19{,}5+2\cdot21{,}5}{50}=16{,}78.$$
 \itemch  Sai. Khả năng tự học trong tuần cao nhất của học sinh là $16{,}28$ giờ.\\
 Khả năng  tự học trong tuần cao nhất của học sinh là $16{,}78$ giờ.
 \end{itemchoice}
 }
\end{ex}

% nguồn: Ok_dungsai_lop11_Lan1
\begin{ex}%[1D5H1-3]%[ma:96a97045]
 Số câu trả lời đúng một bài thi trắc nghiệm môn Sinh học gồm 50 câu của  nhóm $11A$ ở một trường THPT như sau.\\
 \centerline{\begin{tabular}{|c|c|c|c|c|c|}
 \hline 
 Số câu đúng & $[14;21)$ & $[21;28)$ & $[28;35)$ & $[35;42)$ & $[42;49)$\\ 
  \hline
 Số học  sinh & $4$ & $8$ & $25$ & $6$ & $7$\\
 \hline 
\end{tabular}}\\
 Khi đó
\choiceTF 
 {\True Giá trị đại diện  của nhóm $[14;21)$ là $17{,}5$}
 {\True Giá trị đại diện của nhóm  $[21;28)$ là $24{,}5$}
 {\True Giá trị đại diện của  nhóm $[42;49)$ là $45{,}5$}
 {Số câu đúng trung  bình là $32{,}26$}
 \loigiai{
 \centerline{\begin{tabular}{|c|c|c|c|c|c|}
  \hline
 Số câu đúng & $[14;21)$ & $[21;28)$ &  $[28;35)$ & $[35;42)$ & $[42;49)$\\
  \hline
 Giá trị đại diện & $17{,}5$ &  $24{,}5$ & $31{,}5$ & $38{,}5$ & $45{,}5$\\
  \hline
 Số học sinh &  $4$ & $8$ & $25$ & $6$ & $7$\\
 \hline 
 \end{tabular}}\\
 \begin{itemchoice}
  \itemch Đúng. Giá trị đại  diện của nhóm $[14;21)$ là $17{,}5$.
  \itemch Đúng.  Giá trị đại diện của nhóm $[21;28)$ là $24{,}5$.
  \itemch  Đúng. Giá trị đại diện của nhóm $[42;49)$ là $45{,}5$.
  \itemch Sai. Số câu đúng  trung bình là $32{,}26$.\\
  Số  câu đúng trung bình
  $$\overline{x}=\dfrac{4\cdot17{,}5+8\cdot24{,}5+25\cdot31{,}5+6\cdot38{,}5+7\cdot45{,}5}{50}=32{,}06.$$ 
 \end{itemchoice}
 }
\end{ex}

% nguồn: Ok_dungsai_lop11_Lan1
\begin{ex}%[1D5H1-3]%[ma:0c51b8fd]
 Thâm niên giảng dạy của một số giáo viên trường THPT được ghi lại ở  bảng sau
 \begin{center}
  \begin{tabular}{|c|c|c|c|c|c|}
    \hline
   Thâm niên (số năm) &  $[1;5)$ & $[5;10)$ & $[10;15)$ & $[15;20)$ & $[20;25)$ \\
    \hline
   Số giáo viên & $4$ & $12$ & $16$ & $8$  & $3$ \\
    \hline
  \end{tabular}
 \end{center}
 \choiceTF 
  {\True Cỡ mẫu  của mẫu số liệu bằng $50$}
  {\True  Số trung bình của mẫu ghép nhóm là $11{,}84$}
  {\True Nhóm chứa mốt của mẫu số liệu trên là  nhóm $[10;15)$}
  {Mốt  của mẫu số liệu ghép nhóm bằng $11{,}74$}
  \loigiai{
   Ta có bảng sau
  \begin{center}
   \begin{tabular}{|c|c|c|c|c|c|}
    \hline 
    Thâm niên (số năm) & $[1;5)$ & $[5;10)$ & $[10;15)$ & $[15;20)$  & $[20;25)$ \\
    \hline 
    Giá trị đại diện & $3$ & $7{,}5$ & $12{,}5$ & $17{,}5$ & $22{,}5$  \\
     \hline
    Số giáo viên & $4$ & $12$ & $16$  & $8$ & $3$ \\
    \hline 
   \end{tabular}
  \end{center}
  \begin{itemchoice}
   \itemch Đúng. Cỡ mẫu của mẫu  số liệu bằng $50$.
   \itemch Đúng. Ước lượng số trung bình của mẫu ghép  là
   \[ 
   \overline{x} = \dfrac{3 \cdot 4+ 7{,}5 \cdot  12 + 12{,}5 \cdot 16 + 17{,}5 \cdot 8 + 22{,}5 \cdot 3}{50} = 11{,}84.
    \]
   \itemch Đúng. Nhóm chứa mốt của mẫu số  liệu trên là nhóm $[10;15)$.
   \itemch Sai. Do đó $u_m = 10$, $n_{m-1}=12$, $n_m = 16$, $n_{m+1}=8$, $n_{m+1}-u_m =  15-10=5$.
    \[
   M_0 = 10 + \dfrac{16-12}{(16-12)+(16-8)} \cdot  5 = 11{,}66.
    \]
  \end{itemchoice}
  }
\end{ex}

% nguồn: Ok_dungsai_lop11_Lan2
\begin{ex}%[1D5H1-3]%[ma:1d32db6b]
	Thống kê số giờ tự học của học sinh tại một trường THPT trong một tuần, người ta ghi được kết quả sau (đơn vị:  giờ).
	\begin{center}
		\begin{tabular}{|l|c|c|c|c|c|c|c|c|c|c|}
			\hline $14$ &  $12{,}5$&$15$&$16{,}5$& $17$& $14{,}5$& $13$& $15{,}5$& $16{,}5$& $17{,}5$\\
			\hline $16{,}5$&  $18{,}5$&$19$&$20$&$19{,}5$&$17$&$16{,}5$&$14$&$18$&$21$\\
			\hline $15{,}5$&$13{,}5$&$12{,}5$&$14{,}5$&$17{,}5$&$19$&$19{,}5$&$20{,}5$&$20$&$17$\\ 
			\hline $14{,}5$&$13$&$14{,}5$&$18$&$16$&$15$&$13$&$18{,}5$&$14{,}5$&$12{,}5$\\ 
			\hline 
		\end{tabular}
	\end{center}
	 \choiceTF
		{Tổng số giờ tự học của học sinh là $829{,}5$ giờ} 
		{\True Tổng hợp kết quả tự học của học sinh vào bảng tần số ghép nhóm theo mẫu  sau
		\begin{center}
			\begin{tabular}{|l|c|c|c|c|c|c|c|}
				\hline Số cuộc gọi & {$[12{,}5;14{,}5)$} & {$ [14{,}5;16{,}5)$} &  {$[16{,}5;18{,}5) $} & {$[18{,}5;20{,}5)$} & {$[20{,}5;22{,}5)$}\\
				\hline Số học sinh& $9$ & $13$ & $17$ & $9$&  $2$\\
				 \hline
		\end{tabular}
	\end{center}}
	{\True Khi mẫu số liệu được ghép nhóm thì thời gian trung bình học của mỗi học sinh  là $16{,}78$}
	{Khả năng tự học trong tuần cao nhất của học sinh là $16{,}28$  giờ}
	\loigiai{ 
		\begin{itemchoice}
			\itemch Sai. Tổng số giờ tự học của học sinh là  $819{,}5$ giờ.\\
			Số giờ tự học trung  bình của học sinh là $\overline{x}=\dfrac{891{,}5}{50}=16{,}39$ giờ.
			\itemch Đúng. Tổng hợp kết quả tự học của học sinh vào  bảng tần số ghép nhóm theo mẫu sau
			\begin{center}
				\begin{tabular}{|l|c|c|c|c|c|c|c|}
					\hline Số cuộc gọi & {$[12{,}5;14{,}5)$} & {$ [14{,}5;16{,}5)$} & {$[16{,}5;18{,}5) $} & {$[18{,}5;20{,}5)$} &  {$[20{,}5;22{,}5)$}\\
					\hline Số học sinh&  $9$ & $13$ & $17$ & $9$& $2$\\
					\hline 
				\end{tabular}
			\end{center}
			\itemch Đúng. Thời gian tự học trung bình là $$\overline{x}=\dfrac{9\cdot 13{,}5+13\cdot  15{,}5+17\cdot 17{,}5+9\cdot 19{,}5+2\cdot 21{,}5}{50}=16{,}78.$$
			\itemch Sai.  Khả năng tự học trong tuần cao nhất của học sinh là $16{,}78$ giờ.
		\end{itemchoice}
	} 
\end{ex}

% nguồn: Ok_dungsai_lop11_Lan2
\begin{ex}%[1D5H1-3]%[ma:bbcd6d02]
	Kết quả khảo sát cân nặng của $25$ quả cam ở mỗi lô hàng $A$, $B$ được cho ở bảng sau 
	\begin{center}
		\begin{tabular}{|l|c|c|c|c|c|c|c|}
			\hline Cân nặng (gam) &  {$[150;155)$} & {$ [155;160)$} & {$[160;165) $} & {$[165;170)$} & {$[170;175)$}\\
			\hline Số quả cam ở lô hàng $A$ & $2$  & $6$ & $12$ & $4$& $1$\\
			\hline Số  quả cam ở lô hàng $B$ & $1$ & $3$ & $7$ & $10$& $4$\\
			 \hline
		\end{tabular}
	\end{center}
	 \choiceTF
		{\True Giá  trị đại diện nhóm $[150;155)$ bằng $152{,}5$}
		{Cân nặng trung bình  của mỗi quả cam ở lô $A$ là $163{,}7$ (gam)}
		 {Cân nặng trung bình của mỗi quả cam ở lô $B$ là $162{,}1$ (gam)}
		{\True Theo số trung bình thì  cam ở lô hàng $B$ nặng hơn cam ở lô hàng $A$}
	\loigiai{ 
		\begin{itemchoice}
			\itemch Đúng. Bảng thống kê số lượng cam theo  giá trị đại diện
			\begin{center}
				\begin{tabular}{|l|c|c|c|c|c|c|c|}
					\hline Cân nặng đại diện (gam) & {$152{,}5$} & {$157{,}5$} & {$162{,}5$} & {$167{,}5$} & {$172{,}5$}\\ 
					\hline Số quả cam ở lô hàng  A & $2$ & $6$ & $12$ & $4$& $1$\\
					\hline Số quả cam ở lô hàng B & $1$ & $3$ & $7$ & $10$& $4$\\ 
					\hline 
				\end{tabular}
			\end{center}
			 Vậy giá trị đại diện nhóm $[150;155)$ bằng $152{,}5$.
			\itemch Sai. Cân nặng trung bình của mỗi  quả cam ở lô $A$ là
			$$\overline{x}_A=\dfrac{152{,}5\cdot 2+157{,}5\cdot 6+162{,}5\cdot  12+167{,}5\cdot 4+172{,}5\cdot 1}{25}=161{,}7\text{ (gam)}.$$
			\itemch Sai. Cân nặng trung bình của mỗi quả cam ở lô $B$  là
			$$\overline{x}_B=\dfrac{152{,}5\cdot 1+157{,}5\cdot 3+162{,}5\cdot 7+167{,}5\cdot 10+172{,}5\cdot  4}{25}=165{,}1\text{ (gam)}.$$
			\itemch Đúng. Ta thấy $\overline{x}_A<\overline{x}_B$  nên nếu so sánh theo số trung bình thì cam ở lô hàng $B$ nặng hơn cam ở lô hàng $A$.
		\end{itemchoice}
	} 
\end{ex}

% nguồn: Ok_dungsai_lop11_Lan2
\begin{ex}%[1D5H1-3]%[ma:eeee3edb]
	Khảo  sát thời gian tập thể dục trong ngày của một số học sinh khối $11$ thu được mẫu số liệu ghép nhóm sau
	\begin{center}
		\begin{tabular}{|c|c|c|c|c|c|}
			\hline 
			Thời gian (phút) & $[0; 20)$ &  $[20; 40)$ & $[40; 60)$ & $[60; 80)$ & $[80; 100)$ \\
			\hline 
			Số học sinh & $5$ &  $9$ & $12$ & $10$ & $6$\\
			\hline 
		\end{tabular}
	\end{center}
	 \choiceTF
		{\True Tổng số học sinh được khảo sát là $42$ học sinh} 
		{Giá trị  đại diện của nhóm $[20;40)$ là $25$}
		{\True Có $16$ học sinh tập thể dục ít nhất $1$  giờ trong ngày}
		{Số  trung bình của mẫu số liệu trên thuộc nhóm $[20;40)$}
	\loigiai{ 
		Bảng tần số ghép nhóm theo  giá trị đại diện là
		\begin{center}
			\begin{tabular}{|c|c|c|c|c|c|}
				 \hline
				Thời gian (phút)  & $[0; 20)$ & $[20; 40)$ & $[40; 60)$ & $[60; 80)$ & $[80; 100)$ \\
				\hline 
				Giá  trị đại diện & $10$ & $30$ & $50$ & $70$ & $90$ \\
				 \hline
				Số học  sinh & $5$ & $9$ & $12$ & $10$ & $6$\\
				\hline 
			\end{tabular}
		\end{center}
		\begin{itemchoice}
			\itemch  \textbf{Đúng.}\\
			Tổng số học  sinh được khảo sát là $5+9+12+10+6=42$.
			\itemch  \textbf{Sai.}\\
			Giá trị đại diện của nhóm $[20;40)$  là $\dfrac{20+40}{2}=30$.
			\itemch  \textbf{Đúng.}\\
			Số học sinh tập thể dục ít nhất  $1$ giờ trong ngày $10+6=16$ học sinh.
			 \itemch \textbf{Sai.}\\
			Số  trung bình của mẫu số liệu trên là
			\[\overline{x}=\dfrac{5\cdot 10+9\cdot 30+12\cdot 59+10\cdot 70+6\cdot  90}{42}=\dfrac{360}{7} \in [40;60).\]
		\end{itemchoice}
	} 
\end{ex}


%%%%%%%%%%%%%%%% TRẢ LỜI NGẮN %%%%%%%%%%%%%%%%

% nguồn: Ok_TraLoiNgan_Lop11_Lan1
\begin{ex}%[1D5N1-3]%[ma:6dd56dd1]
	Phòng tập thể dục thể thao của một trường THPT đã ghi lại số giờ học  sinh khối 11 sử dụng cơ sở vật chất để tập luyện trong một tháng. Dữ liệu thu được trong bảng dưới đây:
	\begin{center}
		\begin{tabular}{|>{\centering\arraybackslash}p{4cm}|>{\centering\arraybackslash}p{1.5cm}|>{\centering\arraybackslash}p{1.5cm}|>{\centering\arraybackslash}p{1.5cm}|>{\centering\arraybackslash}p{1.5cm}|>{\centering\arraybackslash}p{1.5cm}|>{\centering\arraybackslash}p{1.5cm}|}
			\hline 
			Thời gian (giờ)&$[1;5)$ & $[5;9)$ & $[9;13)$  &$[13;17)$ & $[17;21)$ &$[21;25)$ \\
			\hline 
			Tần số (số học sinh) &$10$ & $14$ & $31$ &  $2$ & $5$ & $23$ \\
			\hline 
		\end{tabular}
	\end{center}
	Tính số trung bình của mẫu số liệu (làm tròn kết quả đến hàng phần mười). 
	\shortans{13,2} 
	\loigiai{ 
		Số trung bình của mẫu số liệu  ghép nhóm là\\
		$$\overline{x}=\dfrac{10\cdot3+14\cdot7+31\cdot11+2\cdot15+5\cdot19+23\cdot23}{85} \approx  13{,}2.$$
	} 
\end{ex}

% nguồn: Ok_TraLoiNgan_Lop11_Lan1
\begin{ex}%[1D5N1-3]%[ma:adc64d1a]
	Bảng sau  cho ta cân nặng của học sinh một lớp $11$:\\
	\begin{center}
		\begin{tabular}{|c|c|c|c|c|c|}
			\hline 
			Cân nặng  (kg)&$[40{,}5;45{,}5)$&$[45{,}5;50{,}5)$&$[50{,}5;55{,}5)$&$[55{,}5;60{,}5)$&$[60{,}5;65{,}5)$\\
			\hline 
			Số học  sinh&$10$&$7$&$16$&$4$&$2$\\
			\hline 
		\end{tabular}
	\end{center}
	Tìm cân nặng (kg) trung bình của học sinh lớp $11$ đó (kết quả làm tròn đến hàng phần mười). 
	\shortans{50,6} 
	 \loigiai{
		Bảng thống  kê cân nặng của học sinh theo giá trị đại diện
		\begin{center}
			\begin{tabular}{|c|c|c|c|c|c|}
				 \hline
				Cân  nặng (kg)&$43$&$48$&$53$&$58$&$63$\\
				\hline 
				Số  học sinh&$10$&$7$&$16$&$4$&$2$\\
				 \hline
			\end{tabular}
		\end{center}
		Cân  nặng trung bình của học sinh lớp $11$ đó là
		\[\overline{x}=\dfrac{43\cdot 10+48\cdot 7+53\cdot 16+58\cdot 4+63\cdot 2}{39}\approx  50{,}6\,\text{ (kg).}\]
	} 
\end{ex}

% nguồn: Ok_TraLoiNgan_Lop11_Lan1
\begin{ex}%[1D5H1-3]%[ma:069e481e]
	Khảo sát chiều cao  của học sinh lớp $11 \mathrm{C} 1$, thu được bảng số liệu ghép nhóm như sau
	\begin{center}
		\begin{tabular}{|l|c|c|c|c|c|}
			\hline Khoàng chiều cao $(\mathrm{cm})$ & {$[150; 155)$} & {$[155;  160)$} & {$[160; 165)$} & {$[165; 170)$} & {$[170; 175)$} \\
			 \hline Số học sinh & $7$ & $15$ & $12$ & $8$ & $3$ \\
			 \hline
		\end{tabular}
	\end{center}
	Chiều cao trung bình của học sinh $11 \mathrm{C} 1$ là bao nhiêu? (Kết quả làm tròn đến hàng  đơn vị).
	 \shortans{161}
	\loigiai{ 
		\begin{center}
			\begin{tabular}{|l|c|c|c|c|c|}
				\hline Khoàng chiều cao $(\mathrm{cm})$ & {$[150; 155)$} & {$[155; 160)$} & {$[160; 165)$} & {$[165; 170)$}  & {$[170; 175)$} \\
				\hline Giá trị đại diện & $152{,}5$  & $157{,}5$ & $162{,}5$ & $167{,}5$ & $172{,}5$ \\
				\hline Số học sinh & $7$ & $15$ & $12$ & $8$  & $3$ \\
				\hline 
			\end{tabular}
		\end{center}
		Chiều cao trung bình của học sinh lớp 11C1  là
		$$\overline{x}=\dfrac{1}{45}(7\cdot 152{,}5+15\cdot 157{,}5+12\cdot 162{,}5+4 \cdot  167{,}5+3\cdot 172{,}5)=160{,}8\approx 161 \mathrm{~cm}.$$
	 }
\end{ex}

% nguồn: Ok_TraLoiNgan_Lop11_Lan1
\begin{ex}%[1D5H1-3]%[ma:39287bd7]
	 \immini{
	Bảng bên cho ta bảng tần số  ghép nhóm số liệu thống kê chiều cao của $40$ mẫu cây ở một vườn thực vật (đơn vị: centimét).
	Xác định số trung bình cộng của mẫu số liệu ghép nhóm trên. 
	 }
	{ 
	\begin{tabular}{|c|c|c|}
		 \hline
		\textbf{Nhóm} & \textbf{Tần số}  & \textbf{Tần số tích luỹ}\\
		 \hline
		$\left[30;40\right)$ & $4$  & $4$\\
		$\left[40;50\right)$ & $10$ & $14$\\ 
		$\left[50;60\right)$ & $14$ &  $28$\\
		$\left[60;70\right)$ & $6$  & $34$\\
		$\left[70;80\right)$ & $4$ & $38$\\ 
		 $\left[80;90\right)$ & $2$ & $40$\\
		\hline 
		& $n = 40$  &\\
		\hline 
	\end{tabular}
	 }
	\shortans{55,5} 
	 \loigiai{
		Trung bình cộng của mẫu số  liệu trên là
		 \[
		\overline{x} = \dfrac{35 \cdot 4 + 45 \cdot 10 + 55 \cdot 14 + 65 \cdot 6 + 75 \cdot 4 + 85 \cdot 2}{40} = 55{,}5\text{  (cm)}.
		\] 
	 }
\end{ex}

% nguồn: Ok_TraLoiNgan_Lop11_Lan1
\begin{ex}%[1D5H1-3]%[ma:67023404]
	Mẫu số  liệu sau ghi lại cân nặng của $30$ bạn học sinh (đơn vị: kilôgam)
	\[ 
	\begin{array}{cccccccccc}
		17 & 40 & 39 & 40{,}5 & 42 & 51 & 41{,}5 & 39 & 41 &  30\\
		40 & 42 & 40{,}5 & 39{,}5 & 41 & 40{,}5 & 37  & 39{,}5 & 40 & 41\\
		38{,}5 & 39{,}5 & 40 & 41 & 39 & 40{,}5 & 40 & 38{,}5 & 39{,}5 &  41{,}5
	\end{array}
	 \]
	Xác định số trung bình cộng của mẫu số liệu ghép nhóm trên. 
	\shortans{40} 
	 \loigiai{
		Trung bình cộng của mẫu  số liệu trên là
		 \[
		\overline{x} = \dfrac{17{,}5 \cdot 1 + 32{,}5 \cdot 1 + 37{,}5 \cdot 10+ 42{,}5 \cdot 17+ 52{,}5  \cdot 1}{30} = 40\text{ (kg)}.
		 \]
	 }
\end{ex}

% nguồn: Ok_TraLoiNgan_Lop11_Lan1
\begin{ex}%[1D5H1-3]%[ma:94715f35]
	Trong một hoạt động thiện nguyện,  các bạn học sinh lớp $11$B đã tiến hành tặng sách cho các trẻ em vùng khó khăn. Kết quả hoạt động được ghi nhận ở bảng sau
	\begin{center}
		\begin{tabular}{|l|c|c|c|c|c|}
			\hline Số sách (quyển) & { $[0; 2)$ } & { $[2; 4)$ } & { $[4; 6)$ } & { $[6; 8)$ } & { $[8; 10]$  } \\
			\hline Số học sinh & $10$ & $13$ &  $8$ & $5$ & $4$ \\
			\hline 
		\end{tabular}
	\end{center}
	Trung bình mỗi bạn của lớp $11$B đem tặng bao  nhiêu quyển sách?
	\shortans{40} 
	 \loigiai{
		Cỡ mẫu $n = 40$. Trung bình mỗi bạn của lớp  11B đem tặng là
		$$ \overline{x} = \dfrac{1\cdot 10+3\cdot 13+5\cdot 8+7\cdot 5+9\cdot 4}{40} = 4 \text{ (quyển sách)}.$$ 
	 }
\end{ex}

% nguồn: Ok_TraLoiNgan_Lop11_Lan1
\begin{ex}%[1D5H1-3]%[ma:21643cf2]
	 Bảng sau cho ta cân nặng của học sinh một nhóm 11.\\
	\centerline{\begin{tabular}{|c|c|c|c|c|c|}
	 \hline
	Cân nặng (kg) & $[40{,}5;45{,}5)$ & $[45{,}5;50{,}5)$ & $[50{,}5;55{,}5)$ & $[55{,}5;60{,}5)$ &  $[60{,}5;65{,}5)$\\
	 \hline
	Số học sinh & $10$ & $7$ &  $16$ & $4$ & $2$\\
	\hline 
\end{tabular}}\\
Tìm cân nặng (kg) trung bình của học sinh nhóm 11 đó. Kết quả làm tròn  đến hàng phần mười.
\shortans{50,6} 
 \loigiai{
	Bảng thống kê cân nặng của học sinh theo giá trị đại  diện.\\
	\centerline{\begin{tabular}{|c|c|c|c|c|c|}
	\hline 
	Cân nặng (kg) & $43$ & $48$ & $53$ &  $58$ & $63$\\
	 \hline
	 Số học sinh & $10$ & $7$ & $16$ & $4$ & $2$\\
	 \hline
	\end{tabular}}\\
	Cân nặng trung bình của  học sinh nhóm 11 đó là
 $$\overline{x}=\dfrac{43\cdot 10+48\cdot 7+53\cdot 16+58\cdot 4+63\cdot 2}{39}\approx 50{,}6\, (\text{ kg}).$$}
\end{ex}

% nguồn: Ok_TraLoiNgan_Lop11_Lan1
\begin{ex}%[1D5H1-3]%[ma:c3b63bd3]
	Khi đo cân nặng của học sinh lớp 11D, y tá lập được bảng số  liệu ghép nhóm sau đây:
	\begin{center}
		\begin{tabular}{|l|c|c|c|c|c|c|}
			\hline Cân nặng  (kg) & {$[40{,}5; 45{,}5)$} & {$[45{,}5; 50{,}5)$} & {$[50{,}5; 55{,}5)$} & {$[55{,}5; 60{,}5)$} & {$[60{,}5; 65{,}5)$} & {$[65{,}5; 70{,}5)$} \\
			\hline  Số học sinh & $10$ & $7$ & $16$ & $4$ & $2$ & $3$ \\
			 \hline
		\end{tabular}
	\end{center}
	Cân nặng trung bình của học sinh  lớp 11D là bao nhiêu? (Kết quả làm tròn đến hàng phần chục)
	\shortans{51,8} 
	 \loigiai{
		Ta có  bảng thống kê
		\begin{center}
			\begin{tabular}{|l|c|c|c|c|c|c|}
				\hline Cân nặng $(\mathrm{kg})$ & {$[40{,}5; 45{,}5)$} & {$[45{,}5; 50{,}5)$} & {$[50{,}5; 55{,}5)$} & {$[55{,}5; 60{,}5)$} & {$[60{,}5; 65{,}5)$}  & {$[65{,}5; 70{,}5)$} \\
				\hline Giá  trị đại diện & $43$ & $48$ & $53$ & $58$ & $63$ & $68$ \\
				\hline Số học sinh & $10$  & $7$ & $16$ & $4$ & $2$ & $3$ \\
				 \hline
			\end{tabular}
		\end{center}
		Cân nặng trung bình của học sinh là $$\overline{x}=\dfrac{43\cdot 10+48\cdot 7+53\cdot 16+58\cdot 4+63\cdot  2+68\cdot 3}{42}=51{,}8.$$
	 }
\end{ex}

% nguồn: Ok_TraLoiNgan_Lop11_Lan1
\begin{ex}%[1D5H1-3]%[ma:ef7a44af]
	Cơ cấu dân số Việt Nam năm 2020 theo độ tuổi được cho trong  bảng sau
	\begin{center}
		\begin{tabular}{|c|c|c|c|c|c|}
			\hline 
			 Độ tuổi & Dưới $5$ tuổi &$5-14$&$15-24$&$25-64$& Trên $65$\\
			\hline 
			 Số người (triệu)&$7{,}89$&$14{,}68$&$13{,}32$&$53{,}78$&$7{,}66$\\
			\hline 
		\end{tabular}
	\end{center}
	Chọn $80$ là giá trị đại diện cho nhóm trên $65$ tuổi. Tính tuổi trung bình của người Việt Nam năm  $2020$. Làm tròn đến hàng phần mười.
	 \shortans{35,2}
	 \loigiai{
		Trong mỗi khoảng độ tuổi, giá trị đại diện là trung bình cộng của giá trị hai đầu mút  nên ta có bảng sau
		\begin{center}
			\begin{tabular}{|c|c|c|c|c|c|}
				 \hline
				Độ tuổi & $2{,}5$ &$9{,}5$&$19{,}5$&$44{,}5$&  $80$\\
				\hline 
				Số  người (triệu)&$7{,}89$&$14{,}68$&$13{,}32$&$53{,}78$&$7{,}66$\\
				\hline 
			\end{tabular}
		\end{center}
		Tổng số dân là $n=97{,}33$ triệu người. Do đó tuổi  trung bình của người dân Việt Nam trong năm $2020$ là
		$$\overline{x}=\dfrac{2{,}5\cdot  7{,}89+9{,}5\cdot 14{,}68+19{,}5\cdot 13{,}32+44{,}5\cdot 53{,}78+80\cdot 7{,}66}{97{,}33}\approx 35{,}2.$$
	 }
\end{ex}

% nguồn: Ok_TraLoiNgan_Lop11_Lan1
\begin{ex}%[1D5H1-3]%[ma:99baa103]
	Người ta đếm số xe ô tô đi qua một trạm  thu phí mỗi phút trong khoảng thời gian từ $9$ giờ đến $9$ giờ $30$ phút sáng. Kết quả được ghi lại ở bảng sau
	\begin{center}
		\begin{tabular}{|c|c|c|c|c|c|c|c|c|c|c|c|c|c|c|}
			\hline $15$ & $16$ & $13$ & $21$ &  $17$ & $23$ & $15$ & $21$ & $6$ & $11$ & $12$ & $23$ & $19$ & $25$ & $11$ \\
			\hline $25$ & $7$ & $29$ & $10$ & $28$ & $29$ & $24$ & $6$ & $11$ & $23$ & $11$ & $21$ & $9$ &  $27$ & $15$ \\
			\hline 
		\end{tabular}
	\end{center}
	Hãy ước lượng trung bình số xe đi qua trạm thu phí trong mỗi phút (làm tròn đến hàng  phần mười).
	 \shortans{17,7}
	\loigiai{ 
		Bảng  tần số ghép nhóm
		\begin{center}
			\begin{tabular}{|c|c|c|c|c|c|}
				\hline Số xe &{$[6; 10]$} &{$[11; 15]$} &{$[16; 20]$}  &{$[21; 25]$} &{$[26; 30]$} \\
				\hline Số lần & $5$ &  $9$ & $3$ & $9$ & $4$ \\
				\hline 
			\end{tabular}
		\end{center}
		Bảng tần số ghép nhóm (theo giá trị đại diện) được hiệu chỉnh lại như sau 
		\begin{center}
			\begin{tabular}{|c|c|c|c|c|c|}
				\hline Số xe &{$[5{,}5; 10{,}5)$} &{$[10{,}5; 15{,}5)$}  &{$[15{,}5; 20{,}5)$} &{$[20{,}5; 25{,}5)$} &{$[25{,}5; 30{,}5)$} \\
				\hline Giá trị đại diện &{$8$} &{$13$} &{$18$}  &{$23$} &{$28$} \\
				\hline Số lần & $5$ & $9$ & $3$ & $9$  & $4$ \\
				 \hline
			\end{tabular}
		\end{center}
		Số xe trung bình đi qua trạm qua bảng tần  số ghép nhóm là
		\[\overline{x}=\dfrac{8\cdot 5+13\cdot 9+18\cdot 3+23\cdot 9+28\cdot  4}{30}\approx 17{,}7\ (\text{xe}). \]
	 }
\end{ex}

% nguồn: Ok_TraLoiNgan_Lop11_Lan1
\begin{ex}%[1D5H1-3]%[ma:219a9bb9]
	\immini{Một nhà thực vật học đo chiều dài của $74$ lá cây (đơn vị: milimét) và thu được bảng tần số như bảng bên. Tính chiều dài trung bình của  $74$ lá cây trên theo đơn vị milimét (làm tròn kết quả đến hàng phần trăm).}
	{ 
	\begin{tabular}{|c|c|c|}
		\hline 
		\textbf{Nhóm} & \textbf{Giá trị đại diện}  & \textbf{Tần số}\\
		\hline 
		$\left[5{,}45;5{,}85\right)$ & $5{,}65$ &  $5$\\
		$\left[5{,}85;6{,}25\right)$ & $6{,}05$ &  $9$\\
		 $\left[6{,}25;6{,}65\right)$ & $6{,}45$ & $15$\\
		$\left[6{,}65;7{,}05\right)$ & $6{,}85$ &  $19$\\
		 $\left[7{,}05;7{,}45\right)$ & $7{,}25$ & $16$\\
		$\left[7{,}45;7{,}85\right)$  & $7{,}65$ & $8$\\
		$\left[7{,}85;8{,}25\right)$  & $8{,}05$ & $2$\\
		 \hline
		&  & $n = 74$\\
		 \hline
	\end{tabular}
	} 
	 \shortans{6,80}
	 \loigiai{
		Chiều dài trung bình  của $74$ lá cây mà nhà thực vật học đo xấp xỉ là
		\[ 
		\overline{x} = \dfrac{5\cdot 5{,}65 + 9 \cdot 6{,}05 + 15\cdot 6{,}45 + 19\cdot 6{,}85 + 16 \cdot 7{,}25 + 8\cdot 7{,}65 + 2\cdot 8{,}05}{74} \approx 6{,}80\ (\text{mm}). 
		\] 
	} 
\end{ex}

% nguồn: Ok_TraLoiNgan_Lop11_Lan1
\begin{ex}%[1D5H1-3]%[ma:93026a9c]
	Khảo sát chiều cao của học sinh lớp $11 \mathrm{C} 1$, thu được bảng số  liệu ghép nhóm như sau
	\begin{center}
		\begin{tabular}{|l|c|c|c|c|c|}
			\hline Khoàng chiều cao $(\mathrm{cm})$  & {$[150; 155)$} & {$[155; 160)$} & {$[160; 165)$} & {$[165; 170)$} & {$[170; 175)$} \\
			\hline Số học sinh & $7$ & $15$ & $12$ & $8$  & $3$ \\
			\hline 
		\end{tabular}
	\end{center}
	Chiều cao trung bình  của học sinh 11C1 có dạng $abc{,}d$, với $a,b,c,d$ là các số tự nhiên. Tính $S=a+b+c+d$.
	\shortans{15} 
	 \loigiai{
		\begin{center}
			\begin{tabular}{|l|c|c|c|c|c|}
				\hline Khoàng chiều cao $(\mathrm{cm})$ & {$[150; 155)$} & {$[155; 160)$} & {$[160; 165)$} & {$[165; 170)$} & {$[170; 175)$}  \\
				\hline  Giá trị đại diện & $152{,}5$ & $157{,}5$ & $162{,}5$ & $167{,}5$ & $172{,}5$ \\
				\hline Số học sinh  & $7$ & $15$ & $12$ & $8$ & $3$ \\
				 \hline
			\end{tabular}
		\end{center}
		 Chiều cao trung bình của học sinh lớp 11C1 là
		$$\overline{x}=\dfrac{1}{45}(7\cdot 152{,}5+15\cdot 157{,}5+12\cdot 162{,}5+4 \cdot 167{,}5+3\cdot  172{,}5)=160{,}8 \mathrm{~cm}.$$
		Khi đó  $S=1+6+0+8=15$.
	} 
\end{ex}

% nguồn: Ok_TraLoiNgan_Lop11_Lan1
\begin{ex}%[1D5H1-3]%[ma:75313d53]
	Để kiểm tra thời gian sử dụng pin của một chiếc điện thoại mới, chị  An thống kê thời gian sử dụng điện thoại của mình từ lúc sạc đầy pin cho đến khi hết pin ở bảng sau:
	\begin{center}
		\begin{tabular}{|c|c|c|c|c|c|} \hline
			\text{Thời gian sử dụng (giờ)} & $\left[7;9\right)$ & $\left[9;11\right)$ &  $\left[11;13\right)$ & $\left[13,15\right)$ & $\left[15;17\right)$ \\ \hline
			\text{Số lần}  & 2 & 5 & 7 & 6 & 3 \\ \hline
		\end{tabular}
	\end{center}
	Hãy ước lượng thời gian sử dụng trung bình từ lúc chị An sạc đầy pin điện  thoại cho tới khi hết pin (làm tròn đến hàng phần mười).
	 \shortans{12,3}
	 \loigiai{
		Thời gian sử  dụng trung bình
		$$\overline{x}=\dfrac{2\cdot8 +5\cdot10 +7\cdot12 +6\cdot14 +3\cdot16}{2+5+7+6+3}=\dfrac{282}{23}\approx 12{,}3.$$ 
	 }
\end{ex}

% nguồn: Ok_TraLoiNgan_Lop11_Lan1
\begin{ex}%[1D5H1-3]%[ma:2c1e4bfc]
	Thống  kê điểm trung bình môn Toán của một số học sinh lớp $11$ được cho ở bảng sau
	\begin{center}
		\begin{tabular}{|c|c|c|c|c|c|c|c|}
			\hline Khoảng điểm &{$[6{,}5; 7)$} &{$[7; 7{,}5)$} &{$[7{,}5; 8)$} &{$[8; 8{,}5)$} &{$[8{,}5;9)$} &{$[9; 9{,}5)$}  &{$[9{,}5; 10)$} \\
			\hline Tần số & $8$ & $10$ & $16$ & $24$ & $13$ & $7$  & $4$ \\
			\hline 
		\end{tabular}
	\end{center}
	Hãy ước lượng  số trung bình của mẫu số liệu ghép nhóm trên (làm tròn đến hàng phần trăm).
	\shortans{8,12} 
	 \loigiai{
		Bảng tần  số ghép nhóm theo giá trị đại diện
		\begin{center}
			\begin{tabular}{|c|c|c|c|c|c|c|c|}
				\hline Khoảng điểm &{$[6{,}5; 7)$} &{$[7; 7{,}5)$} &{$[7{,}5; 8)$} &{$[8; 8{,}5)$} &{$[8{,}5;9)$} &{$[9; 9{,}5)$}  &{$[9{,}5; 10)$} \\
				\hline Giá trị đại diện  &$6{,}75$ &$7{,}25$ &$7{,}75$ &$8{,}25$ &$8{,}75$ &$9{,}25$ &$9{,}75$ \\
				\hline Tần số & $8$ & $10$ & $16$ & $24$ & $13$  & $7$ & $4$ \\
				 \hline
			\end{tabular}
		\end{center}
		Điểm trung bình môn Toán  của một số học sinh lớp $11$ là
		\[\overline{x}=\dfrac{8\cdot 6{,}75+10\cdot 7{,}25+16\cdot 7{,}75+24\cdot  8{,}25+13\cdot 8{,}75+7\cdot 9{,}25+4\cdot 9{,}75}{82}\approx 8{,}12.\]
	 }
\end{ex}

% nguồn: Ok_TraLoiNgan_Lop11_Lan1
\begin{ex}%[1D5H1-3]%[ma:1196098b]
	Doanh thu bán hàng  trong 20 ngày được lựa chọn ngẫu nhiên của một cửa hàng được ghi lại ở bảng sau (đơn vị: triệu đồng):
	\begin{center}
		\begin{tabular}{|l|c|c|c|c|c|}
			\hline Doanh thu & {$[5; 7)$} & {$[7; 9)$} & {$[9; 11)$}  & {$[11; 13)$} & {$[13; 15)$} \\
			\hline Số ngày & $2$ &  $7$ & $7$ & $3$ & $1$ \\
			\hline 
		\end{tabular}
	\end{center}
	Doanh  số trung bình của cửa hàng trên trong $20$ ngày là bao nhiêu? (Kết quả làm tròn đến hàng phần chục)
	\shortans{9,4} 
	 \loigiai{
		Bảng  tần số ghép nhóm theo giá trị đại diện là
		\begin{center}
			\begin{tabular}{|l|c|c|c|c|c|}
				\hline Doanh thu & {$[5; 7)$} & {$[7; 9)$} & {$[9; 11)$} & {$[11; 13)$} &  {$[13; 15)$} \\
				\hline Giá trị đại diện & $6$  & $8$ & $10$ & $12$ & $14$ \\
				 \hline Số ngày & $2$ & $7$ & $7$ & $3$ & $1$ \\
				\hline 
			\end{tabular}
		\end{center}
		Số trung bình $\overline{x}=\dfrac{2\cdot 6+7\cdot 8+7\cdot 10+3\cdot 12+1\cdot  14}{20}=9{,}4$.
	 }
\end{ex}

% nguồn: Ok_TraLoiNgan_Lop11_Lan2
\begin{ex}%[1D5H1-3]%[ma:6b94845c]
	Dưới đây là một mẫu số liệu cho ở dạng bảng tần số ghép  nhóm
	\begin{center}
		\begin{tabular}{|c|c|c|c|c|c|}
			 \hline
			Nhóm & [0; 50) & [50; 100) & [100; 150) &  [150; 200) & \\
			\hline 
			Tần  số & 6 & 8 & 7 & 6 & $N=27$\\
			 \hline
		\end{tabular}
	\end{center}
	Trung  bình của mẫu số liệu xấp xỉ bằng (kết quả làm tròn đến hàng phần chục).
	 \shortans{99,1}
	\loigiai{ 
		 $$
		\overline{x} = \frac{25 \cdot 6 + 75  \cdot 8 + 125 \cdot 7 + 175 \cdot 6}{27} \approx 99{,}1
		$$ 
	 }
\end{ex}

% nguồn: Ok_TraLoiNgan_Lop11_Lan2
\begin{ex}%[1D5H1-3]%[ma:537c248f]
	Người ta tiến hành  phỏng vấn $50$ người về phim chiếu rạp Lật mặt $6$ của Lý Hải. Người điều tra yêu cầu cho điểm phim theo thang điểm $100$. Kết quả được trình bày trong bảng phân bố tần số ghép lớp sau đây
	\begin{center}
		\begin{tabular}{|c|c|c|c|c|c|}
			 \hline
			Số điểm  & $[50; 60)$ & $[60; 70)$ & $[70; 80)$ & $[80; 90)$ & $[90; 100)$\\
			 \hline
			Số người & $4$ & $7$ & $9$ &  $18$ & $12$\\
			 \hline
		\end{tabular}
	\end{center}
	Hãy  tìm số trung bình của mẫu số liệu ghép nhóm ở trên (làm tròn đến hàng phần mười). \shortans{80,4}
	\loigiai{ 
		\begin{center}
			\begin{tabular}{|c|c|c|c|c|c|}
				 \hline
				Số  điểm & $[50; 60)$ & $[60; 70)$ & $[70; 80)$ & $[80; 90)$ & $[90; 100)$\\
				\hline 
				Giá trị đại diện & $55$ & $65$ & $75$ & $85$ &  $95$\\
				 \hline
				Số người  & $4$ & $7$ & $9$ & $18$ & $12$\\
				 \hline
			\end{tabular}
		\end{center}
		Điểm trung bình đánh giá phim Lật  Mặt $6$ bằng
		\[\overline{x}=\dfrac{55\cdot 4+65\cdot 7+75\cdot 9+85\cdot 18+95\cdot 12}{50}=80{,}4.\] 
	} 
\end{ex}

% nguồn: Ok_TraLoiNgan_Lop11_Lan2
\begin{ex}%[1D5H1-3]%[ma:dd974bca]
	Kết quả đo chiều  cao của $250$ cây dừa đột biến $3$ năm tuổi ở một viện nghiên cứu được tổng hợp ở bảng sau
	\begin{center}
		\begin{tabular}{|c|c|c|c|c|c|}
			 \hline
			Chiều cao $(\text{m})$ & $[8{,}5; 8{,}8)$  & $[8{,}8; 9{,}1)$ & $[9{,}1; 9{,}4)$ & $[9{,}4; 9{,}7)$ & $[9{,}7; 10)$\\
			\hline 
			Số cây & $36$ & $45$ &  $83$ & $65$ & $21$\\
			\hline 
		\end{tabular}
	\end{center}
	Hãy tìm số trung bình của mẫu số liệu ghép nhóm trên  (đơn vị: mét, làm tròn đến hàng phần mười). \shortans{9,2}
	 \loigiai{
		\begin{center}
			\begin{tabular}{|c|c|c|c|c|c|}
				\hline 
				Chiều cao $(\text{m}^2)$ & $[8{,}5; 8{,}8)$ & $[8{,}8; 9{,}1)$ & $[9{,}1;  9{,}4)$ & $[9{,}4; 9{,}7)$ & $[9{,}7; 10)$\\
				 \hline
				Giá trị  đại diện & $8{,}65$ & $8{,}95$ & $9{,}25$ & $9{,}55$ & $9{,}85$\\
				\hline 
				Số cây & $36$ & $45$ & $83$ & $65$ & $21$\\ 
				 \hline
			\end{tabular}
		\end{center}
		Chiều cao trung bình của 250 cây  dừa đột biến xấp xỉ bằng
	\[\overline{x}=\dfrac{8{,}65  \cdot 36+8{,}95 \cdot 45+9{,}25 \cdot 83+9{,}55 \cdot 65+9{,}85 \cdot 21}{250}\approx 9{,}2\,\, (\text{m}).\]}
\end{ex}

% nguồn: Ok_TraLoiNgan_Lop11_Lan2
\begin{ex}%[1D5H1-3]%[ma:30762162]
	Bảng  sau cho ta cân nặng của học sinh một lớp $11$
	\begin{center}
		\begin{tabular}{|c|c|c|c|c|c|}
			 \hline
			Cân  nặng (kg) & $[40{,}5;45{,}5)$ & $[45{,}5;50{,}5)$ & $[50{,}5;55{,}5)$ & $[55{,}5;60{,}5)$ & $[60{,}5;65{,}5)$\\
			 \hline
			 Số học sinh & $10$ & $7$ & $16$ & $4$ & $2$\\
			\hline 
		\end{tabular}
	\end{center}
	Tìm cân nặng trung bình của học sinh lớp 11 đó  (làm tròn đến hàng phần mười). \shortans{50,6}
	\loigiai{ 
		 Bảng thống kê cân nặng của học sinh theo giá trị đại diện
		\begin{center}
			\begin{tabular}{|c|c|c|c|c|c|}
				 \hline
				Cân nặng (kg) & $43$ &  $48$ & $53$ & $58$ & $63$\\
				 \hline
				Số học sinh & $10$ & $7$ & $16$ & $4$  & $2$\\
				\hline 
			\end{tabular}
		\end{center}
		Cân nặng trung  bình của học sinh lớp 11 đó là
	\[\overline{x}=\dfrac{43\cdot 10+48\cdot 7+53\cdot 16+58\cdot 4+63\cdot 2}{39}\approx  50{,}6\,\, \text{(kg)}.\]}
\end{ex}

% nguồn: Ok_TraLoiNgan_Lop11_Lan2
\begin{ex}%[1D5H1-3]%[ma:db7be7a7]
	Đo chiều cao của $500$ học sinh trong một trường  THPT X ta thu được kết quả như sau
	\begin{center}
		\begin{tabular}{|c|c|c|c|c|c|}
			\hline 
			 Chiều cao & $\left[150;\,154\right)$ & $\left[154;\,158\right)$ & $\left[158;\,162\right)$ & $\left[162;\,166\right)$ & $\left[166;\,170\right)$\\
			 \hline
			Tần số & $25$ &  $50$ & $200$ & $175$ & $50$\\
			\hline 
		\end{tabular}
	\end{center}
	Tính chiều cao trung bình của $500$ học sinh trong một trường THPT X (làm tròn đến hàng đơn vị).  \shortans{161}
	 \loigiai{
		Ta có bảng sau\\ 
		\centerline{\begin{tabular}{|c|c|c|}
		 \hline
		Lớp chiều cao &  Giá trị đại diện & Tần số\\
		\hline 
		$\left[150;\,154\right)$ & $152$ &  $25$\\
		 \hline
		$\left[154;\,158\right)$ & $156$  & $50$\\
		\hline 
		$\left[158;\,162\right)$ & $160$  & $200$\\
		\hline 
		$\left[162;\,166\right)$ & $164$ & $175$\\ 
		 \hline
		$\left[166;\,170\right)$ &  $168$ & $50$\\
		 \hline
	\end{tabular}}\\
	Ta có chiều  cao trung bình
\[\overline{x}=\dfrac{1}{500}\left(152\cdot 25+156\cdot 50+160\cdot 200+164\cdot 175+168\cdot 50\right)\approx 161.\]} 
\end{ex}

% nguồn: Ok_TraLoiNgan_Lop11_Lan2
\begin{ex}%[1D5H1-3]%[ma:7902452e]
	Kết quả khảo sát cân nặng của $25$ quả cam ở lô hàng A được cho  ở bảng sau:
	\begin{center}
		\begin{tabular}{|c|c|c|c|c|c|}
			\hline 
			Cân nặng (g) & [150; 155) & [155; 160) & [160;  165) & [165; 170) & [170; 175) \\
			 \hline
			Số quả cam ở lô hàng A & 2 & 6 & 12 & 4 & 1  \\
			\hline 
		\end{tabular}
	\end{center}
	Cân nặng trung bình của mỗi quả cam ở lô hàng A xấp xỉ bằng  (kết quả làm tròn đến hàng đơn vị).
	 \shortans{162}
	 \loigiai{
		Ta có bảng thống kê số lượng cam theo  giá trị đại diện:
		\begin{center}
			\begin{tabular}{|c|c|c|c|c|c|}
				\hline 
				Cân nặng đại diện (g) &$152{,}5$ & $157{,}5$&$ 162{,}5$ & $167{,}5$  & $172{,}5$ \\
				 \hline
				Số quả cam ở lô hàng A & 2 & 6 & 12 & 4 &  1 \\
				 \hline
			\end{tabular}
		\end{center}
		Cân nặng trung bình của mỗi quả cam ở lô hàng A  xấp xỉ bằng
		$$\overline{x}=\frac{2 \cdot 152{,}5 + 6 \cdot 157{,}5 + 12 \cdot 162{,}5 + 4 \cdot 167{,}5  + 1 \cdot 172{,}5}{25} \approx 161{,}7 \,\, \text{(g)}$$.
	} 
\end{ex}

% nguồn: Ok_TraLoiNgan_Lop11_Lan1
\begin{ex}%[1D5V1-3]%[ma:9151f3fb]
	Thời gian truy cập internet mỗi buổi trưa của một số học sinh được  cho trong bảng sau:
	\begin{center}
		\begin{tabular}{|l|c|c|c|c|c|}
			\hline Thời gian (phút) & {$[9{,}5; 12{,}5)$} & {$[12{,}5; 15{,}5)$} & {$[15{,}5; 18{,}5)$} & {$[18{,}5; 21{,}5)$}  & {$[21{,}5; 24{,}5)$} \\
			\hline Số học sinh & $3$ & $12$ & $15$ & $24$  & $2$ \\
			 \hline
		\end{tabular}
	\end{center}
	Số trung vị của mẫu số liệu trên là bao  nhiêu? (Kết quả làm tròn đến hàng phần chục)
	\shortans{18,1} 
	 \loigiai{
		Cỡ mẫu  là $n=3+12+15+24+2=56$.\\
		Gọi $x_1; \ldots; x_{56}$ là thời gian truy cập internet của $56$ học sinh và giả  sử dãy này đã được sắp xếp theo thứ tự tăng dần.\\
		Khi đó, số trung vị là $\dfrac{x_{28}+x_{29}}{2}$.\\ 
		Do hai giá trị $x_{28}; x_{29}$ thuộc nhóm $[15{,}5; 18{,}5)$  nên nhóm này chứa trung vị.\\
		Do đó $p=3$; $a_3=15{,}5$; $m_3=15$;  $m_1+m_2=15$; $a_4-a_3=3$.
		 $$M_e=15{,}5+\dfrac{\dfrac{56}{2}-15}{15}\cdot 3=18{,}1.$$
	} 
\end{ex}

% nguồn: Ok_TraLoiNgan_Lop11_Lan1
\begin{ex}%[1D5V1-3]%[ma:d557aca8]
	Lớp $11$A trường THPT  có $40$ học sinh tham gia trả lời $40$ câu hỏi trong một bài kiểm tra khảo sát. Kết quả được thống kê trong bảng tần số ghép nhóm sau
	\begin{center}
		\begin{tabular}{|c|c|c|c|c|c|}
			 \hline
			 Số câu trả lời đúng & $\left[16;21\right)$ & $\left[21;26\right)$ & $\left[26;31\right)$ & $\left[31;36\right)$ & $\left[36;41\right)$ \\
			\hline 
			Số học sinh & $4$ &  $m$ & $8$ & $18$ & $n$ \\
			 \hline
		\end{tabular}
	\end{center}
	Hỏi giá trị $m-n$ bằng bao nhiêu, biết rằng số câu trả lời trung bình  của lớp $11A$ là $30$?
	 \shortans{2}
	\loigiai{ 
		Lớp $11$A có  $40$ học sinh nên
		$$4+m+8+18+n=40  \Rightarrow m+n=10.\qquad(1)$$
		Ta có bảng sau 
		\begin{center}
			\begin{tabular}{|c|c|c|c|c|c|}
				\hline 
				Số câu trả lời đúng & $\left[16;21\right)$ & $\left[21;26\right)$ & $\left[26;31\right)$  & $\left[31;36\right)$ & $\left[36;41\right)$ \\
				 \hline
				Giá trị đại diện & $18{,}5$ & $23{,}5$ & $28{,}5$  & $33{,}5$ & $38{,}5$ \\
				 \hline
				Số  học sinh & $4$ & $m$ & $8$ & $18$ & $n$ \\
				 \hline
			\end{tabular}
		\end{center}
		 Khi đó
		\[\overline{x} = \dfrac{4 \cdot 18{,}5+m \cdot 23{,}5+8\cdot 28{,}5+18 \cdot 33{,}5+n\cdot 38{,}5}{40}  = 30.\]
		Do  đó
		$$23{,}5 \cdot  m + 38{,}5 \cdot n=295. \qquad(2)$$
		Từ (1) và (2) ta suy  ra $m=6$ và $n=4$. Do đó $m-n=2$.
	 }
\end{ex}

% nguồn: Ok_TraLoiNgan_Lop11_Lan1
\begin{ex}%[1D5V1-3]%[ma:d0fa9fe6]
	Mẫu số  liệu dưới đây ghi lại tốc độ của $40$ ô tô khi đi qua một trạm đo tốc độ (đơn vị: km/h).
	\begin{center}
		\begin{tabular}{cccccccccc}
			$49$ & $42$  & $51$ & $55$ & $45$ & $60$ &$53$ & $55$& $44$ & $65$ \\
			$5$2  &$62$ &$41$ &$44$& $57$ &$56$& $68$&$48$ &$46$&$53$ \\
			$63$  &$49$& $54$ &$61$ &$59$& $57$& $47$& $50$&$60$&$62$\\
			$48$ &$52$& $58$& $47$& $60$& $55$&  $45$&$47$& $48$&$61$\\
		\end{tabular}
	\end{center}
	Sau khi ghép  nhóm mẫu số liệu trên thành sáu nhóm ứng với sáu nửa khoảng
	 \[
	[40;45), \,\,  [45;50),\,\, [50;55),\,\, [55;60),\,\, [60;65),\,\, [65;70)
	 \]
	Hỏi tốc độ trung bình thông qua mẫu số liệu ghép nhóm trên là bao nhiêu? (Làm tròn kết quả đến hàng phần  mười).
	 \shortans{53,9}
	\loigiai{ 
		 Ta có bảng tần số
		\begin{center}
			\begin{tabular}{|c|c|c|c|c|c|c|}
				\hline 
				 Nhóm & $[40;45)$& $[45;50)$&$[50;55)$&$[55;60)$&$[60;65)$&$[65;70)$\\
				\hline 
				Giá  trị đại diện & $42{,}5$& $47{,}5$&$52{,}5$&$57{,}5$&$62{,}5$&$67{,}5$\\
				 \hline
				Tần số &$4$&  $11$&$7$&$8$&$8$&$2$\\
				\hline 
			\end{tabular}
		\end{center}
		 Khi đó
		$$\overline{x}=\dfrac{42{,}5 \cdot 4 + 47{,}5\cdot 11+  52{,}5\cdot 7+57{,}5\cdot8+62{,}5\cdot 8+ 67{,}5\cdot 2}{40}=53{,}875\approx 53{,}9.$$
	 }
\end{ex}

% nguồn: Ok_TraLoiNgan_Lop11_Lan1
\begin{ex}%[1D5V1-3]%[ma:74f918a3]
	Thống  kê điểm trung bình môn Toán của một số học sinh lớp $11$ được cho ở bảng sau:
	\begin{center}
		\begin{tabular}{|c|c|c|c|c|c|c|c|}
			\hline Khoảng  điểm &{$[6{,}5; 7)$} &{$[7; 7{,}5)$} &{$[7{,}5; 8)$} &{$[8; 8{,}5)$} &{$[8{,}5;9)$} &{$[9; 9{,}5)$} &{$[9{,}5; 10)$} \\
			\hline Tần số & $8$ & $10$ & $16$ & $24$  & $13$ & $7$ & $4$ \\
			 \hline
		\end{tabular}
	\end{center}
	Hãy ước lượng số  trung bình của mẫu số liệu ghép nhóm trên.
	 \shortans{8,15}
	\loigiai{ 
		Bảng tần số ghép nhóm  theo giá trị đại diện
		\begin{center}
			\begin{tabular}{|c|c|c|c|c|c|c|c|}
				\hline Khoảng điểm &{$[6{,}5; 7)$} &{$[7; 7{,}5)$} &{$[7{,}5; 8)$} &{$[8;  8{,}5)$} &{$[8{,}5;9)$} &{$[9; 9{,}5)$} &{$[9{,}5; 10)$} \\
				\hline Giá trị đại diện &$6{,}75$ &$7{,}25$  &$7{,}75$ &$8{,}25$ &$8{,}75$ &$9{,}25$ &$9{,}75$ \\
				\hline Tần số & $8$ & $10$ & $16$  & $24$ & $13$ & $7$ & $4$ \\
				\hline 
			\end{tabular}
		\end{center}
		Tứ  phân vị thứ hai là $\dfrac{x_{41}+x_{42}}{2}\in [8;8{,}5)$. Do đó
		\[Q_2=u_m+\dfrac{\dfrac{n}{2}-C}{n_m}\cdot\left(u_{m+1}-u_m\right)=8+\dfrac{\dfrac{82}{2}-34}{24}\cdot  (8{,}5-8)=8{,}15. \]
	} 
\end{ex}

% nguồn: Ok_TraLoiNgan_Lop11_Lan1
\begin{ex}%[1D5V1-3]%[ma:73de069f]
	Số liệu thống kê $100$ học sinh tham gia kì thi học sinh giỏi toán (thang điểm $20$).  Kết quả được thống kê trong bảng sau:
	\begin{center}
		\begin{tabular}{|c|c|c|c|c|c|c|c|c|c|c|c|c|}
			\hline Điểm & $9$ & $10$ & $11$ &  $12$ & $13$ & $14$ & $15$ & $16 $& $17$ & $18$ & $19$ & \\
			\hline Tần số & $1$ & $1$ & $3$ & $5$ & $8$ & $13$ & $19$ & $24$ & $14$ &  $10$ & $2$ & $N=100$ \\
			 \hline
		\end{tabular}
	\end{center}
	Phương sai của bảng số liệu thống kê  là (làm tròn đến hàng phần trăm)
	\shortans{3,96} 
	\loigiai{Điểm số trung bình của các học  sinh tham gia thi học sinh giỏi là\\
		 $$\overline{x}=\dfrac{1\cdot9+1\cdot10+3\cdot11+\ldots+14\cdot17+10\cdot18+2\cdot19}{100}=15{,}23.$$
		 Phương sai của số liệu thống kê là
	$$S_x^2=\dfrac{(\overline{x}-9)^2+(\overline{x}-10)^2+3(\overline{x}-11)^2+5(\overline{x}-12)^2+\ldots+2(\overline{x}-19)^2}{100}  \approx 3{,}96.$$}
\end{ex}


%%%%%%%%%%%%%%%% TỰ LUẬN %%%%%%%%%%%%%%%%

% nguồn: Ok_TuanLuan_Toan11_Lan1
\begin{bt}%[1D5N1-3]%[ma:5ee20a0f]
	Chiều cao của $20$  cây giống được cho bởi bảng sau
	\begin{center}
		\begin{tabular}{|c|c|}
			\hline 
			Lớp (cm)& Tần số \\ 
			\hline 
			$[40;44]$& $2$ \\ 
			 \hline
			 $[45;49]$& $5$ \\
			 \hline
			$[50;54]$& $3$  \\
			 \hline
			 $[55;59]$& $4$ \\
			\hline 
			 $[60;64]$& $3$ \\
			 \hline
			 $[64;69]$& $3$ \\
			 \hline
			& $N=20$  \\
			\hline 
		\end{tabular}
	\end{center}
	Tính số trung bình (làm tròn đến chữ  số thứ hai sau dấu phẩy) của bảng nói trên.
	\loigiai{Bảng tần số ghép lớp  của bảng nói trên là
		\begin{center}
			\begin{tabular}{|c|c|c|}
				 \hline
				Lớp (cm) & Giá trị đại diện& Tần số  \\
				\hline 
				$[40;44]$ &$42$& $2$ \\ 
				\hline 
				$[45;49]$ &$47$& $5$ \\ 
				\hline 
				 $[50;54]$ &$52$& $3$ \\
				\hline 
				$[55;59]$ &$57$&  $4$ \\
				\hline 
				 $[60;64]$ &$62$& $3$ \\
				 \hline
				$[64;69]$  &$67$& $3$ \\
				 \hline
				& &  $N=20$ \\
				\hline 
			\end{tabular}
		\end{center}
		Áp dụng công  thức tính số trung bình cho bảng tần số ghép lớp ta có
		$$\overline{x}= \dfrac{42\cdot2+47\cdot5+52\cdot3+57\cdot4+62\cdot3+67\cdot3}{20}\approx54,5  \text{ (cm).}$$
	 }
\end{bt}

% nguồn: Ok_TuanLuan_Toan11_Lan1
\begin{bt}%[1D5N1-3]%[ma:d724d4c8]
	Tìm hiểu thời gian xem tivi trong tuần trước (đơn vị: giờ) của một số học sinh thu được kết  quả sau
	\begin{center}
		\begin{tabular}{|c|c|c|c|c|c|}
			 \hline
			Thời gian& $[0;5)$  & $[5;10)$ & $[10;15)$ & $[15;20)$ & $[20;25)$ \\
			\hline 
			Số học sinh& $8$ & $16$ & $4$ & $2$ & $2$  \\
			\hline 
		\end{tabular}
	\end{center}
	Tính thời  gian xem tivi trung bình trong tuần trước của các bạn học sinh này.
	\loigiai{Trong mỗi khoảng thời gian, giá trị đại diện trung bình  cộng của giá trị hai đầu mút nên ta có bảng sau
		\begin{center}
			\begin{tabular}{|c|c|c|c|c|c|}
				 \hline
				Thời gian&  $2{,}5$ & $7{,}5$ & $12{,}5$ & $17{,}5$ & $22{,}5$ \\
				 \hline
				Số học  sinh& $8$ & $16$ & $4$ & $2$ & $2$ \\
				\hline 
			\end{tabular}
		\end{center}
		Tổng số học sinh là  $n=32$.\\
		Thời gian  xem tivi trung bình của học sinh là
		$$\overline{x}=\dfrac{2,5.8+7,5.16+12,5.4+17,5.2+22,5.2}{32}\approx  8{,}44\ \text{ (giờ)}.$$
	 }
\end{bt}

% nguồn: Ok_TuanLuan_Toan11_Lan2
\begin{ex}%[1D5N1-3]%[ma:b9eca157]
	Chiều cao của  $20$ cây giống được cho bởi bảng sau
	\begin{center}
		\begin{tabular}{|c|c|}
			\hline 
			 Lớp (cm)& Tần số \\
			\hline 
			$[40;44]$& $2$  \\
			 \hline
			$[45;49]$&  $5$ \\
			\hline 
			$[50;54]$& $3$ \\ 
			 \hline
			$[55;59]$&  $4$ \\
			\hline 
			$[60;64]$& $3$  \\
			 \hline
			$[64;69]$& $3$  \\
			\hline 
			& $N=20$  \\
			\hline 
		\end{tabular}
	\end{center}
	 Tính số trung bình (làm tròn đến chữ số thứ hai sau dấu phẩy) của bảng nói trên.
	\loigiai{Bảng tần số  ghép lớp của bảng nói trên là
		\begin{center}
			\begin{tabular}{|c|c|c|}
				 \hline
				Lớp (cm) & Giá trị đại diện&  Tần số \\
				 \hline
				$[40;44]$  &$42$& $2$ \\
				\hline 
				$[45;49]$  &$47$& $5$ \\
				 \hline
				$[50;54]$ &$52$& $3$  \\
				\hline 
				$[55;59]$ &$57$&  $4$ \\
				\hline 
				$[60;64]$  &$62$& $3$ \\
				\hline 
				$[64;69]$ &$67$&  $3$ \\
				\hline 
				& & $N=20$ \\ 
				 \hline
			\end{tabular}
		\end{center}
		Áp dụng công thức tính số trung bình cho bảng tần số ghép lớp ta  có
		 $$\overline{x}= \dfrac{42\cdot2+47\cdot5+52\cdot3+57\cdot4+62\cdot3+67\cdot3}{20}\approx54,5 \text{ (cm).}$$
	} 
\end{ex}

% nguồn: Ok_TuanLuan_Toan11_Lan2
\begin{ex}%[1D5N1-3]%[ma:55a8ec0f]
	Tìm hiểu thời gian xem tivi trong tuần trước (đơn vị: giờ) của một số học sinh  thu được kết quả sau
	\begin{center}
		\begin{tabular}{|c|c|c|c|c|c|}
			\hline 
			Thời gian& $[0;5)$  & $[5;10)$ & $[10;15)$ & $[15;20)$ & $[20;25)$ \\
			\hline 
			Số học sinh& $8$ & $16$ & $4$ & $2$ & $2$  \\
			\hline 
		\end{tabular}
	\end{center}
	Tính thời gian xem tivi trung bình  trong tuần trước của các bạn học sinh này.
	\loigiai{Trong mỗi khoảng thời gian, giá trị đại  diện trung bình cộng của giá trị hai đầu mút nên ta có bảng sau
		\begin{center}
			\begin{tabular}{|c|c|c|c|c|c|}
				 \hline
				Thời gian& $2{,}5$ & $7{,}5$ & $12{,}5$  & $17{,}5$ & $22{,}5$ \\
				 \hline
				Số học sinh& $8$ &  $16$ & $4$ & $2$ & $2$ \\
				 \hline
			\end{tabular}
		\end{center}
		Tổng  số học sinh là $n=32$.\\
		Thời gian xem tivi trung bình của học  sinh là
		$$\overline{x}=\dfrac{2,5.8+7,5.16+12,5.4+17,5.2+22,5.2}{32}\approx 8{,}44\  \text{ (giờ)}.$$
	 }
\end{ex}

% nguồn: Ok_TuanLuan_Toan11_Lan2
\begin{ex}%[1D5N1-3]%[ma:e1ea3b63]
	Thống kê chiều cao của  $60$ học sinh lớp 11 ở một trường trung học phổ thông (đơn vị: cm) thu được kết quả như sau
	\begin{center}
		\begin{tabular}{|l|c|c|c|c|c|c|c|c|c|c|c|c|c|}
			\hline $157$ & $163$& $146$& $150$& $176$& $162$& $147$& $151$& $176$& $161$& $149$& $152$& $177$\\ 
			\hline $165$ & $148$& $153$& $176$& $169$& $152$& $155$&  $170$& $144$& $168$& $160$& $170$& $143$\\
			 \hline $167$ & $160$& $170$& $142$& $166$& $160$& $165$& $141$& $174$& $161$& $166$& $144$& $173$\\
			\hline $162$ & $175$& $156$& $161$& $172$&  $175$& $157$& $162$& $171$& $176$& $160$& $158$& $170$\\
			\hline $176$ & $164$& $159$& $170$& $175$& $163$& $160$& $170$& $ $& $ $& $ $& $ $&  $ $\\
			\hline 
		\end{tabular}
	\end{center}
	\begin{enumEX}{1}
		\item Hãy thu gọn dữ liệu bằng cách lập bảng tần số ghép  nhóm với các nhóm ghép lần lượt là $[135;145);[145;155)$; $[155;165)$; $[165;175)$,$[175;185)$.
		\item Hãy tìm chiều cao trung bình của học sinh trong lớp học đó từ bảng  tần số ghép nhóm đã lập ở câu a.
	\end{enumEX}
	\loigiai{ 
		\begin{enumEX}{1}
			\item Bảng tần số ghép nhóm chiều cao của $60$ học sinh lớp 11 là 
			\begin{center}
				\begin{tabular}{|l|c|c|c|c|c|c|}
					\hline Chiều cao (cm) & {$[135;145)$} & {$ [145;155)$} & {$[155;165) $} &  {$[165;175)$}& {$[175;185)$}\\
					\hline Tần số & $5$ & $9$ &  $19$ & $17$& $10$\\
					 \hline
				\end{tabular}
			\end{center}
			\item Ta có bảng thống  kê thời gian chờ khám bệnh theo giá trị đại diện
			\begin{center}
				\begin{tabular}{|l|c|c|c|c|c|c|}
					\hline Chiều cao (cm) & {$[135;145)$}  & {$ [145;155)$} & {$[155;165) $} & {$[165;175)$}& {$[175;185)$}\\
					\hline Chiều cao đại diện  & $140$ & $150$ & $160$ & $170$& $180$\\
					\hline Tần  số & $5$ & $9$ & $19$ & $17$& $10$\\
					\hline 
				\end{tabular}
			\end{center}
			Chiều cao trung bình là $$\overline{x}=(140\cdot 5+150\cdot 9+160\cdot 19+170\cdot  17+180\cdot 10):60=163\text{(cm)}.$$
		\end{enumEX}
	} 
\end{ex}

% nguồn: Ok_TuanLuan_Toan11_Lan2
\begin{ex}%[1D5N1-3]%[ma:9ed90dbf]
	Kết quả điều tra thời  gian chờ khám (đơn vị: phút) của bệnh nhân tại phòng khám của một bệnh viện ở Thành phố Hồ Chí Minh được cho trong bảng sau
	\begin{center}
		\begin{tabular}{|l|c|c|c|c|c|c|}
			\hline Thời gian (phút) & {$[20;25)$} & {$  [25;30)$} & {$[30;35) $} & {$[35;40)$}\\
			\hline Tần số & $4$ & $5$ & $8$ &  $3$\\
			\hline 
		\end{tabular}
	\end{center}
	Hãy tính  thời gian trung bình chờ khám bệnh của mẫu số liệu ghép nhóm trên
	 \loigiai{
		Ta có bảng  thống kê thời gian chờ khám bệnh theo giá trị đại diện
		\begin{center}
			\begin{tabular}{|l|c|c|c|c|c|c|}
				\hline Thời gian (phút) &  {$[20;25)$} & {$ [25;30)$} & {$[30;35) $} & {$[35;40)$}\\
				\hline Thời gian đại diện &  $22{,}5$ & $27{,}5$ & $32{,}5$ & $37{,}5$\\
				\hline Tần số & $4$ & $5$ & $8$ &  $3$\\
				\hline 
			\end{tabular}
		\end{center}
		Thời gian trung bình chờ khám bệnh  $(22{,}5\cdot 4+27{,}5.5+32{,}5\cdot 8+37{,}5\cdot3):20=30$.
	 }
\end{ex}

% nguồn: Ok_TuanLuan_Toan11_Lan2
\begin{ex}%[1D5N1-3]%[ma:9689ef58]
	Thống kê cân nặng (kg) của $30$ học sinh trong  một lớp học thu được kết quả như sau
	\begin{center}
		\begin{tabular}{|l|c|c|c|c|c|c|c|c|c|c|}
			\hline $54$ &  $42$& $53$& $63$& $58$& $54$& $63$& $70$& $59$& $49$\\
			\hline $73$ & $49$& $66$& $60$& $55$& $48$& $65$&  $63$& $56$& $51$\\
			\hline $61$ & $57$& $69$& $67$&  $63$& $52$& $60$& $56$& $49$& $44$\\
			 \hline
		\end{tabular}
	\end{center}
	\begin{enumEX}{1}
		\item Hãy tính cân nặng trung bình của  học sinh lớp học trên (kết quả làm tròn đến hàng phần trăm);
		\item Hãy thu gọn dữ liệu bằng cách lập bảng tần số ghép nhóm  với các nhóm ghép lần lượt là $[40;50);[50;60)$; $[60;70)$; $[70;80)$.
		\item Hãy ước lượng cân nặng trung bình của học sinh trong lớp từ  bảng tần số ghép nhóm trên.
	\end{enumEX}
	 \loigiai{
		\begin{enumEX}{1}
			\item Cân nặng  trung bình của học sinh lớp học trên là $\overline{x}=\dfrac{1729}{30}\approx 57{,}63$.
			\item Bảng tần số  ghép nhóm
			\begin{center}
				\begin{tabular}{|l|c|c|c|c|c|c|}
					\hline Cân nặng(kg)  & {$[40;50)$} & {$ [50;60)$} & {$[60;70) $} & {$[70;80)$}\\
					\hline Số học sinh & $6$ & $12$ &  $10$ & $2$\\
					 \hline
				\end{tabular}
			\end{center}
			\item Ta có bảng thống kê cân nặng học  sinh theo giá trị đại diện
			\begin{center}
				\begin{tabular}{|l|c|c|c|c|c|c|}
					\hline Cân nặng(kg) & {$[40;50)$} & {$ [50;60)$} & {$[60;70)  $} & {$[70;80)$}\\
					\hline cân nặng đại diện &  $45$ & $55$ & $65$ & $75$\\
					\hline Số học  sinh & $6$ & $12$ & $10$ & $2$\\
					 \hline
				\end{tabular}
			\end{center}
			Cân nặng trung bình của học sinh lớp học trên  là $(45\cdot 6+55\cdot 12+65\cdot 10+75\cdot 2):30=57{,}67$.
		\end{enumEX}
	} 
\end{ex}

% nguồn: Ok_TuanLuan_Toan11_Lan1
\begin{bt}%[1D5H1-3]%[1D5H1-4]%[1D5H2-2]%[1D5H2-3]%[ma:21c6c5de]
	Thống kê điểm trung bình môn Toán của một số học sinh lớp $11$ được cho bởi  bảng sau:
	\begin{center}
		\resizebox{\linewidth}{!}{ 
		\begin{tabular}{|c|c|c|c|c|c|c|c|c|}
			\hline 
			\textbf{Khoảng  điểm} &
			$\left[6;6{,}5\right)$  &
			$\left[6{,}5;7\right)$ & 
			 $\left[7;7{,}5\right)$ &
			$\left[7{,}5;8\right)$ & 
			$\left[8;8{,}5\right)$ & 
			$\left[8{,}5;9\right)$  &
			$\left[9;9{},5\right)$  &
			$\left[9{,}5;10\right)$ \\ 
			 \hline
			\textbf{Tần  số} &
			$6$  &
			$8$  &
			$28$  &
			$24$  &
			 $30$ &
			 $12$ &
			$12$ & 
			$7$  \\
			 \hline
		\end{tabular}}
	\end{center}
	Hãy ước lượng \textbf{số trung bình, mốt, trung vị} và \textbf{tứ phân vị thứ ba} của  mẫu số liệu ghép nhóm trên.
	\loigiai{ 
		\begin{itemize}
			\item Bảng tần số ghép nhóm theo giá trị đại diện 
			\begin{center}
				 \resizebox{\linewidth}{!}{
				\begin{tabular}{|c|c|c|c|c|c|c|c|c|}
					 \hline
					\textbf{Khoảng  điểm} &
					 $\left[6;6{,}5\right)$ &
					$\left[6{,}5;7\right)$  &
					$\left[7;7{,}5\right)$ & 
					$\left[7{,}5;8\right)$  &
					$\left[8;8{,}5\right)$ & 
					$\left[8{,}5;9\right)$  &
					$\left[9;9{},5\right)$ & 
					$\left[9{,}5;10\right)$ \\ 
					\hline 
					 \textbf{Tần số} &
					 $6$ &
					 $8$ &
					$28$  &
					$24$  &
					$30$  &
					$12$  &
					$12$ & 
					$7$ \\ 
					 \hline
					\textbf{Giá trị đại diện} & 
					$6{,}25$ & 
					$6{,}75$  &
					$7{,}25$ & 
					$7{,}75$ & 
					 $8{,}25$ &
					$8{,}75$  &
					 $9{,}25$ &
					$9{,}5$ \\ 
					 \hline
				\end{tabular}}
			\end{center}
			\item 
			Điểm trung bình môn Toán của  một số học sinh lớp $11$
			$$\bar{x}=\dfrac{6{,}25\cdot6+6{,}75\cdot8+\cdots  +9{,}25\cdot12+9{,}75\cdot7}{127}\approx 7{,}97$$
			 \item
			Nhóm chứa mốt của mẫu  số liệu là $\left[8;8,5 \right)$, do đó
			$${{M}_{0}}=8+\dfrac{30-24}{\left( 30-24 \right)+\left( 30-12 \right)}\left( 8{,}5-8 \right)=8{,}125$$ 
			 \item
			Nhóm chứa trung vị của mẫu số liệu  là $\left[7{,}5;8 \right)$\\
			Trung  vị của mẫu số liệu ghép nhóm:
			${{M}_{e}}=7{,}5+\dfrac{\dfrac{127}{2}-\left( 6+8+28 \right)}{24}\left( 8-7{,}5 \right)\approx 7{,}95$ 
			\item Nhóm chứa tứ phân vị thứ ba của  mẫu số liệu là $\left[8;8{,}5 \right)$\\
			Tứ phân vị thứ ba của mẫu số liệu ghép  nhóm:
			 $${{Q}_{3}}=8+\dfrac{\dfrac{3\cdot127}{4}-\left( 6+8+28+24 \right)}{30}\left( 8{,}5-8 \right)=8{,}4875$$
		\end{itemize}
	} 
\end{bt}

% nguồn: Ok_TuanLuan_Toan11_Lan2
\begin{ex}%[1D5H1-3]%[1D5H1-4]%[1D5H2-2]%[1D5H2-3]%[ma:3358bb47]
	Thống kê điểm trung bình môn Toán của một  số học sinh lớp $11$ được cho bởi bảng sau:
	\begin{center}
		\resizebox{\linewidth}{!}{ 
		\begin{tabular}{|c|c|c|c|c|c|c|c|c|}
			 \hline
			 \textbf{Khoảng điểm} &
			$\left[6;6{,}5\right)$ & 
			$\left[6{,}5;7\right)$  &
			 $\left[7;7{,}5\right)$ &
			$\left[7{,}5;8\right)$ & 
			 $\left[8;8{,}5\right)$ &
			$\left[8{,}5;9\right)$  &
			$\left[9;9{},5\right)$  &
			 $\left[9{,}5;10\right)$ \\
			\hline 
			\textbf{Tần số} & 
			$6$  &
			$8$  &
			$28$ & 
			$24$  &
			 $30$ &
			$12$  &
			$12$  &
			$7$ \\ 
			 \hline
		\end{tabular}}
	\end{center}
	Hãy ước lượng \textbf{số trung bình, mốt, trung vị} và \textbf{tứ phân vị thứ ba}  của mẫu số liệu ghép nhóm trên.
	\loigiai{ 
		\begin{itemize}
			\item  Bảng tần số ghép nhóm theo giá trị đại diện
			\begin{center}
				\resizebox{\linewidth}{!}{ 
				\begin{tabular}{|c|c|c|c|c|c|c|c|c|}
					\hline 
					\textbf{Khoảng điểm} & 
					$\left[6;6{,}5\right)$ & 
					 $\left[6{,}5;7\right)$ &
					$\left[7;7{,}5\right)$ & 
					$\left[7{,}5;8\right)$ & 
					 $\left[8;8{,}5\right)$ &
					$\left[8{,}5;9\right)$  &
					$\left[9;9{},5\right)$  &
					 $\left[9{,}5;10\right)$ \\
					\hline 
					\textbf{Tần số}  &
					 $6$ &
					$8$  &
					$28$  &
					$24$ & 
					$30$ & 
					$12$ & 
					$12$ & 
					$7$  \\
					 \hline
					 \textbf{Giá trị đại diện} &
					 $6{,}25$ &
					$6{,}75$  &
					$7{,}25$ & 
					 $7{,}75$ &
					 $8{,}25$ &
					$8{,}75$ & 
					$9{,}25$  &
					$9{,}5$  \\
					 \hline
				\end{tabular}}
			\end{center}
			\item 
			Điểm trung bình môn Toán của một số học sinh lớp  $11$
			$$\bar{x}=\dfrac{6{,}25\cdot6+6{,}75\cdot8+\cdots +9{,}25\cdot12+9{,}75\cdot7}{127}\approx 7{,}97$$ 
			 \item
			Nhóm chứa  mốt của mẫu số liệu là $\left[8;8,5 \right)$, do đó
			$${{M}_{0}}=8+\dfrac{30-24}{\left( 30-24 \right)+\left( 30-12  \right)}\left( 8{,}5-8 \right)=8{,}125$$
			 \item
			Nhóm chứa trung vị của mẫu  số liệu là $\left[7{,}5;8 \right)$\\
			Trung vị của mẫu số  liệu ghép nhóm:
			${{M}_{e}}=7{,}5+\dfrac{\dfrac{127}{2}-\left( 6+8+28 \right)}{24}\left(  8-7{,}5 \right)\approx 7{,}95$
			\item Nhóm chứa tứ phân vị thứ ba  của mẫu số liệu là $\left[8;8{,}5 \right)$\\
			Tứ phân vị  thứ ba của mẫu số liệu ghép nhóm:
			$${{Q}_{3}}=8+\dfrac{\dfrac{3\cdot127}{4}-\left(  6+8+28+24 \right)}{30}\left( 8{,}5-8 \right)=8{,}4875$$
		\end{itemize}
	 }
\end{ex}

% nguồn: Ok_TuanLuan_Toan11_Lan2
\begin{ex}%[1D5H1-3]%[ma:7ab1d3c7]
	Người ta đếm số xe ô tô đi qua một trạm thu phí mỗi phút trong  khoảng thời gian từ $9$ giờ đến $9$ giờ $30$ phút sáng. Kết quả được ghi lại ở bảng sau:
	\begin{center}
		\begin{tabular}{|c|c|c|c|c|c|c|c|c|c|c|c|c|c|c|}
			 \hline
			$15$ & $16$ & $13$ & $21$ & $17$ & $23$ & $15$ & $21$ & $6$ & $11$ & $12$ & $23$  & $19$ & $25$ & $11$ \\
			\hline 
			$25$ & $7$ & $29$ &  $10$ & $28$ & $29$ & $24$ & $6$ & $11$ & $23$ & $11$ & $21$ & $9$ & $27$ & $15$ \\
			 \hline
		\end{tabular}
	\end{center}
	\begin{enumerate}
		\item Tính số xe trung bình đi qua trạm  thu phí trong mỗi phút.
		\item Tổng hợp lại số liệu trên vào bảng tần số ghép nhóm theo mẫu sau: 
		\begin{center}
			\begin{tabular}{|c|c|c|c|c|c|}
				\hline 
				Số xe & $[6;10]$ & $[11;15]$ & $[16;20]$ & $[21;25]$ & $[26;30]$ \\ 
				 \hline
				Số lần &? &?  &? &? &? \\
				 \hline
			\end{tabular}
		\end{center}
		\item Hãy ước lượng trung bình số xe đi qua trạm thu phí trong mỗi phút từ  bảng tần số ghép nhóm trên.
	\end{enumerate}
	 \loigiai{
		\begin{enumerate}
			\item Số xe trung bình đi qua  trạm thu phí trong mỗi phút là $17{,}4$ xe.
			 \item Bảng tần số ghép nhóm
			\begin{center}
				\begin{tabular}{|c|c|c|c|c|c|}
					 \hline
					Số  xe & $[6;10]$ & $[11;15]$ & $[16;20]$ & $[21;25]$ & $[26;30]$ \\
					\hline 
					Số lần & $5$ & $9$ & $3$  & $9$ & $4$ \\
					 \hline
				\end{tabular}
			\end{center}
			\item  Hiệu chỉnh lại bảng số liệu ta có
			\begin{center}
				\begin{tabular}{|c|c|c|c|c|c|}
					\hline 
					Số  xe & $[5{,}5;10{,}5)$ & $[10{,}5;15{,}5)$ & $[15{,}5;20{,}5)$ & $[20{,}5;25{,}5)$ & $[25{,}5;30{,}5)$ \\
					 \hline
					Giá trị đại  diện & $8$ & $13$ & $18$ & $23$ & $28$ \\
					\hline 
					Số lần & $5$ & $9$ & $3$ & $9$ & $4$  \\
					 \hline
				\end{tabular}
			\end{center}
			Trung bình số xe  đi qua trạm thu phí mỗi phút xấp xỉ bằng
			\[\dfrac{8\cdot 5+13\cdot 9+18\cdot 3+23\cdot 9+28\cdot  4}{30}=17{,}7.\]
		\end{enumerate}
	} 
\end{ex}

% nguồn: Ok_TuanLuan_Toan11_Lan2
\begin{ex}%[1D5H1-3]%[1D5H2-3]%[ma:605a9c10]
	Cân nặng của một số con lợn con mới sinh thuộc hai giống A và B được cho ở biểu đồ dưới đây (đơn vị:  kg).
	\begin{center}
		\begin{tikzpicture}[>=stealth,declare function={ox=-2; kc=3; rc=1.0; kcy=1.5; }]
				 \foreach[count=\i] \yt in {10,20,30,40}{
					\draw[gray!35]  (0,{\i*kcy}) node[left,text=black]{$\yt$}--+(14,0);
				} 
				\foreach[count=\i] \yt/\yh/\xt in {0.8/1.3/{[1{,}0;  1{,}1)},
					 2.8/1.4/{[1{,}1; 1{,}2)},
					3.2/2.4/{[1{,}2;  1{,}3)},
					1.7/1.4/{[1{,}3; 1{,}4)} 
					 }{
					\pgfmathsetmacro{\gtm}{int  (\yt*10)}
					 \pgfmathsetmacro{\gth}{int (\yh*10)}
					\fill[blue!75!black] ({ox+\i*kc},0) rectangle +(rc,{\yt*kcy}) 
					node[above,shift={(-0.5*rc,0)}]{$\gtm$}; 
					\fill[orange] ({ox+\i*kc+rc+0.25},0)  rectangle +(rc,{\yh*kcy})
					node[above,shift={(-0.5*rc,0)}]{$\gth$};; 
					\path ({ox+\i*kc+rc},0) node[below]{$\xt$}; 
				} 
				\draw[->]  (0,0)--(14,0) node[anchor=west,shift={(-15pt,-10pt)}]{(kg)};
				 \draw[->] (0,0)--(0,6.75) node[above]{(Số con)};
				\path (0,0)  node[left]{$0$};
				\path (current bounding box.north) node[below,font=\bfseries\sffamily\large]{Cân nặng của  một số lợn cong mới sinh};
			\end{tikzpicture}
	\end{center}
	\begin{enumerate}
		\item Hãy so sánh cân năng của lợn con mới sinh giống A và giống B theo số trung  bình và trung vị.
		\item Hãy ước lượng tứ phân vị thứ nhất và thứ ba của cân năng lợn con mới sinh giống  A và của cân năng lợn con mới sinh giống B.
	\end{enumerate}
	\loigiai{ 
		\begin{enumerate}
			\item Ta có bảng tần  số ghép nhóm như sau
			\begin{longtable}{|*{5}{c|}}
				\hline 
				Cân nặng  (cm) & [$1{,}0$; $1{,}1$) & [$1{,}1$; $1{,}2$) & [$1{,}2$; $1{,}3$) & [$1{,}3$; $1{,}4$) \\
				\hline 
				Cân nặng đại diện &  $1{,}05$ & $1{,}15$ & $1{,}25$ & $1{,}35$ \\
				\hline 
				Giống A & $8$ & $28$ & $32$ & $17$  \\
				 \hline
				Giống B & $13$ & $14$  & $24$ & $14$ \\
				\hline 
			\end{longtable}
			Theo số trung  bình\\
			Cân nặng trung bình của các lợn con mới  sinh giống A xấp xỉ bằng\\
			$\overline{x_1}=\dfrac{8\cdot 1{,}05+28\cdot  1{,}15+32\cdot 1{,}25+17\cdot 1{,}35}{8+28+32+17} \approx 1{,}218$ (kg).\\
			Cân nặng trung  bình của các lợn con mới sinh giống B xấp xỉ bằng\\
			$\overline{x_2}=\dfrac{13\cdot 1{,}05+14\cdot 1{,}15+24\cdot  1{,}25+14\cdot 1{,}35}{13+14+24+14}=1{,}21$ (kg).\\
			Vậy theo cân nặng trung bình giống A nặng hơn giống B.\\ 
			 Theo trung vị:\\
			Gọi $x_1, x_2, \ldots, x_{85}$ là cân nặng của $85$ con lợn  giống A mới sinh xếp theo thứ tự không giảm.\\
			Số  trung vị của mẫu số liệu trên là $x_{43}$.\\
			Ta có $x_1, x_2, \ldots, x_8 \in \left[1{,}0; 1{,}1\right)$; $x_9, x_{10}, \ldots, x_{36} \in \left[1{,}1; 1{,}2\right)$; $x_{37},  x_{38}, \ldots, x_{68} \in \left[1{,}2; 1{,}3\right)$.\\
			Nhận  thấy $x_{43}$ thuộc $\left[1{,}2; 1{,}3\right)$.\\
			Số trung vị của mẫu số liệu ghép nhóm giống A là 
			\begin{align*}
				M_o=1{,}2+\dfrac{\dfrac{85}{2}-\left(8+28\right)}{32} \left(1{,}3-1{,}2\right)\approx  1{,}22.
			\end{align*}
			Gọi $y_1, y_2, \ldots, y_{65}$ là cân nặng của 65 con lợn giống B mới sinh xếp  theo thứ tự không giảm.\\
			Số trung vị của mẫu số liệu trên  là $y_{33}$.\\
			Ta có $y_1, y_2, \ldots, y_{13} \in \left[1{,}0; 1{,}1\right)$; $y_{14}, y_{15}, \ldots, y_{27} \in \left[1{,}1; 1{,}2\right)$;  $y_{28}, y_{29}, \ldots, y_{51} \in \left[1{,}2; 1{,}3\right)$.\\
			Nhận  thấy $y_{33}$ thuộc $\left[1{,}2; 1{,}3\right)$\\
			Số trung vị của mẫu số liệu ghép nhóm giống  B là
			\begin{align*}
				M'_o =1{,}2+\dfrac{\dfrac{65}{2}-\left(13+14\right)}{24} \left(1{,}3-1{,}2\right)\approx 1{,}155. 
			\end{align*}
			Vậy theo trung vị, cân nặng của lợn con giống A  nặng hơn cân nặng của lợn con giống B.
			 \item Xét giống A:\\
			 Tứ phân vị thứ nhất của dãy số liệu $x_1, x_2, \ldots, x_{85}$ là $x_{21}$ và $x_{22}$.\\
			Vì  $x_{21}, x_{22} \in \left[1{,}1; 1{,}2\right)$ nên tứ phân vị thứ nhất của mẫu số liệu là
			\begin{align*}
				Q_1=1{,}1+\dfrac{\dfrac{1\cdot 85}{4}-8}{28} \left(1{,}2-1{,}1\right)\approx  1{,}147.
			\end{align*}
			Tứ phân vị thứ ba của dãy  số liệu $x_1, x_2, \ldots, x_{85}$ là $x_{64}$ và $x_{65}$.\\
			Vì $x_{64}, x_{65} \in  \left[1{,}2; 1{,}3\right)$ nên tứ phân vị thứ ba của mẫu số liệu là
			\begin{align*}
				Q_3=1{,}2+\dfrac{\dfrac{3\cdot 85}{4}-\left(8+28\right)}{35} \left(1{,}3-1{,}2\right)\approx  1{,}279.
			\end{align*}
			 Xét giống B:\\
			Tứ phân vị thứ nhất của dãy số liệu $y_1, y_2, \ldots, y_{65}$  là $y_{16}$ và $y_{17}$.\\
			Vì $y_{16}, y_{17} \in \left[1{,}1; 1{,}2\right)$ nên tứ  phân vị thứ nhất của mẫu số liệu là
			\begin{align*}
				 Q_1=1{,}1+\dfrac{\dfrac{1\cdot 65}{4}-13}{14} \left(1{,}2-1{,}1\right)\approx 1{,}123.
			\end{align*}
			Tứ phân vị thứ ba của dãy  số liệu $y_1, y_2, \ldots, y_{65}$ là $y_{49}$ và $y_{50}$.\\
			Vì $y_{49}, y_{50} \in \left[1{,}2; 1{,}3\right)$ nên tứ phân  vị thứ ba của mẫu số liệu là
			\begin{align*}
				Q_3=1{,}2+\dfrac{\dfrac{3\cdot 65}{4}-\left(13+14\right)}{24} \left(1{,}3-1{,}2\right)\approx 1{,}29. 
			\end{align*}
		\end{enumerate}
	 }
\end{ex}

% nguồn: Ok_TuanLuan_Toan11_Lan2
\begin{ex}%[1D5H1-3]%[ma:c3e071db]
	Một kĩ thuật viên ghi lại cân  nặng của $20$ chi tiết máy ở bảng sau (đơn vị: gam):
	\begin{center}
		\begin{tabular}{|*{10}{c|}}
			 \hline
			5{,}63 & 5{,}58 & 5{,}42 & 5{,}58 & 5{,}56 & 5{,}54 & 5{,}55 & 5{,}40 & 5{,}60 &  5{,}56 \\
			 \hline
			5{,}46 & 5{,}51 &  5{,}58 & 5{,}48 & 5{,}61 & 5{,}50 & 5{,}54 & 5{,}64 & 5{,}43 & 5{,}63 \\
			 \hline
		\end{tabular}
	\end{center}
	\begin{listEX}
		\item Tính cân nặng trung  bình của mỗi chi tiết máy.
		\item Lập bảng tần số ghép nhóm của mẫu số liệu trên với  nhóm đầu tiên là $[5{,}40; 5{,}45)$ và ước lượng số trung bình của mẫu số liệu ghép nhóm trên.
	\end{listEX}
	\loigiai{ 
		\begin{listEX}
			\item Cân nặng trung bình của  mỗi chi tiết máy là $5{,}54$\,g.
			\item Bảng số liệu ghép nhóm của mẫu số liệu đã cho như sau: 
			\begin{center}
				\begin{tabular}{|*6{c|}}
					 \hline
					Cân nặng & $[5{,}40; 5{,}45)$& $[5{,}45; 5{,}50)$  & $[5{,}50; 5{,}55)$ & $[5{,}55; 5{,}60)$& $[5{,}60; 5{,}65)$ \\
					 \hline
					Số chi tiết & 3  & 2 & 4 & 6 & 5 \\
					\hline 
				\end{tabular}
			\end{center}
			Số trung bình của mẫu số liệu ghép nhóm  là $5{,}545$.
		\end{listEX}
	} 
\end{ex}

% nguồn: Ok_TuanLuan_Toan11_Lan2
\begin{ex}%[1D5H1-3]%[ma:d8c11b21]
	Cho bảng tần số ghép nhóm số liệu thống kê chiều cao của $40$  mẫu cây ở một vườn thực vật (đơn vị: centimét).
	\begin{center}
		\begin{tabular}{|c|c|c|c|c|c|c|}
			 \hline
			Chiều cao(m) & $[30;40)$ & $[40;50)$ &  $[50;60)$ & $[60;70)$ & $[70;80)$ & $[80;90)$ \\
			 \hline
			Số cây &$4$ &$10$  &$14$ &$6$ &$4$ &$2$ \\
			 \hline
		\end{tabular}
	\end{center}
	\begin{enumerate}
		\item Xác định số trung bình cộng của mẫu số  liệu ghép nhóm trên.
		\item Mốt của mẫu số liệu trên là bao  nhiêu?
	\end{enumerate}
	 \loigiai{
		\begin{enumerate}
			\item\, 
			\begin{center}
				\begin{tabular}{|c|c|c|c|c|c|c|}
					 \hline
					Chiều cao(m) & $[30;40)$ & $[40;50)$  & $[50;60)$ & $[60;70)$ & $[70;80)$ & $[80;90)$ \\
					 \hline
					Giá trị  đại diện & $35$ & $45$ & $55$ & $65$ & $75$ & $85$ \\
					\hline 
					Số cây & $4$ &  $10$ & $14$ & $6$ & $4$ & $2$ \\
					\hline 
				\end{tabular}
			\end{center}
			Giá trị  trung bình của mẫu số liệu là
			$$\overline{x}=\dfrac{35\cdot 4+45\cdot 10+55\cdot 14+65\cdot 6+75\cdot 4+85\cdot  2}{40} \approx 63{,}43.$$
			 \item Nhóm chứa mốt: $\left[50;60\right)$.\\
			Mốt của mẫu số  liệu là
			$$M_0=50+\dfrac{4}{4+8}\cdot 10\approx 53{,}3.$$ 
		\end{enumerate}
	 }
\end{ex}

% nguồn: Ok_TuanLuan_Toan11_Lan2
\begin{ex}%[1D5H1-3]%[ma:56aa324f]
	Xét mẫu số liệu được cho dưới dạng bảng tần số ghép nhóm  như dưới đây
	\begin{center}
		\begin{tabular}{|c|c|c|}
			\hline 
			 $\textbf{Nhóm}$ & $\textbf{Giá trị đại diện}$ & \textbf{Tần số} \\
			\hline $[40;47)$ & $x_1=?$  & $n_1=6$\\
			 \hline $[47;54)$ & $x_2=?$ & $n_2=12$\\
			\hline $[54;61)$ & $x_3=?$ & $n_3=10$\\ 
			\hline $[61;68)$ & $x_4=?$ & $n_4=5$\\ 
			\hline  $[68;75)$ & $x_5=?$ & $n_5=3$\\
			 \hline & &$n=36$\\
			 \hline
		\end{tabular}
	\end{center}
	Tính giá trị trung bình  cho bởi công thức sau $\overline{x}=\dfrac{m_1x_1+\ldots+m_5x_5}{n}$.
	 \loigiai{
		\begin{center}
			\begin{tabular}{|c|c|c|}
				\hline 
				 $\textbf{Nhóm}$ & $\textbf{Giá trị đại diện}$ & \textbf{Tần số} \\
				\hline $[40;47)$ &  $x_1=43{,}5$ & $n_1=6$\\
				 \hline $[47;54)$ & $x_2=50{,}5$ & $n_2=12$\\
				\hline $[54;61)$ & $x_3=57{,}5$  & $n_3=10$\\
				\hline $[61;68)$ & $x_4=64{,}5$ & $n_4=5$\\ 
				\hline $[68;75)$ & $x_5=71{,}5$ &  $n_5=3$\\
				 \hline & &$n=36$\\
				 \hline
			\end{tabular}
		\end{center}
		Số trung bình của mẫu số liệu  trên là
		\begin{eqnarray*}
			\overline{x}&=&\dfrac{m_1x_1+\ldots+m_5x_5}{n}\\ 
			&=& \dfrac{43{,}5\cdot  6+50{,}5\cdot 12+57{,}5\cdot 10+64{,}5\cdot 5+71{,}5\cdot 3}{36}\\
			 &\approx& 55.
		\end{eqnarray*}
	 }
\end{ex}

% nguồn: Ok_TuanLuan_Toan11_Lan2
\begin{ex}%[1D5H1-3]%[ma:17ca0e56]
	Diện tích các tỉnh và thành phố khu vực Nam Bộ được thống  kê ở bảng sau:
	\begin{center}
		\begin{tabular}{|*4{c|}}
			\hline 
			Tỉnh/ thành  phố & Diện tích (km$^2$) & Tỉnh/ thành phố & Diện tích (km$^2$)\\
			\hline 
			Bình Phước & $6\,877$ & Vĩnh  Long & $1\,526$ \\
			 \hline
			 Tây Ninh & $4\,041$ & Đồng Tháp & $3\,384$ \\
			\hline 
			Bình Dương & $2\,695$ & An  Giang & $3\,537$ \\
			 \hline
			Đồng Nai & $5\,864$ & Kiên Giang &  $6\,349$ \\
			\hline 
			Bà Rịa - Vũng Tàu & $1\,981$ & Cần Thơ & $1\,439$  \\
			\hline 
			TP.Hồ Chí Minh & $2\,061$ & Hậu Giang & $1\,622$ \\ 
			 \hline
			Long An &  $4\,495$ & Sóc Trăng & $3\,312$ \\
			 \hline
			Tiền Giang & $2\,511$ & Bạc  Liêu & $2\,669$ \\
			 \hline
			Bến Tre  & $2\,395$ & Cà Mau & $5\,221$ \\
			 \hline
			 Trà Vinh & $2\,358$ & & \\
			\hline 
		\end{tabular}
	\end{center}
	\begin{flushright}
		(Nguồn: Tổng cục  Thống kê)
	\end{flushright}
	\begin{listEX}
		\item Hãy tính diện tích trung bình của mỗi tỉnh/thành phố khu vực  Nam Bộ.
		\item Dựa vào số liệu trên, hãy hoàn thiện bảng tần số ghép nhóm về diện tích các tỉnh khu vực Nam Bộ theo mẫu sau: 
		\begin{center}
			\begin{tabular}{|*5{c|}}
				\hline 
				Diện tích (km$^2$) & $[1\,000; 2\,500)$ & $[2\,500; 4\,000)$  & $[4\,000; 5\,500)$ & $[5\,500; 7\,000)$ \\
				 \hline
				Số tỉnh/thành phố  & $?$ & $?$ & $?$ & $?$ \\
				 \hline
			\end{tabular}
		\end{center}
		\item Hãy ước lượng số  trung bình và mốt của mẫu số liệu ghép nhóm trên.
	\end{listEX}
	 \loigiai{
		\begin{listEX}
			\item Ta có $\begin{aligned} 6\,877+4\,041+2\,695+5\,864+1\,981+2\,061+4\,495+2\,511+2\,395&\\
			{}+2\,358+1\,526+3\,384+3\,537+6\,349+1\,439+1\,622+3\,312+2\,669+5\,221 & =  64\,337.
		\end{aligned}$\\
		Diện tích trung bình của mỗi tỉnh/thành phố khu vực Nam Bộ  là
		$\overline{x} = \dfrac{64\,337}{19} \approx 3386{,}16\,\text{(km$^2$)}$. 
		 \item Bảng tần số ghép nhóm:
		\begin{center}
			\begin{tabular}{|*5{c|}}
				\hline 
				Diện  tích (km$^2$) & $[1\,000; 2\,500)$ & $[2\,500; 4\,000)$ & $[4\,000; 5\,500)$ & $[5\,500; 7\,000)$ \\
				\hline 
				Số tỉnh/thành phố & 7 & 6 & 3  & 3 \\
				\hline 
			\end{tabular}
		\end{center}
		\item Bảng tần số ghép nhóm  bao gồm giá trị đại diện của các nhóm như sau:
		\begin{center}
			\begin{tabular}{|*5{c|}}
				\hline 
				Diện tích (km$^2$) & $[1\,000; 2\,500)$ &  $[2\,500; 4\,000)$ & $[4\,000; 5\,500)$ & $[5\,500; 7\,000)$ \\
				 \hline
				Giá trị đại diện & $1\,750$ & $3\,250$ &  $4\,750$ & $6\,250$ \\
				\hline 
				Số tỉnh/ thành phố & 7 & 6 & 3 & 3  \\
				 \hline
			\end{tabular}
		\end{center}
		Khi đó, số trung bình của mẫu số liệu ghép nhóm  là
		\[\overline{x}=\dfrac{7\cdot 1\,750+6\cdot 3\,250+3 \cdot  4\,750+3\cdot 6\,250}{19} \approx 3\,407{,}89.\]
		Nhóm chứa mốt của mẫu số liệu trên là nhóm $\left[1\,000;2\,500  \right)$.\\
		Do đó $u_m=1\,000$; $n_{m-1}=0$;  $n_m=7$; $n_{m+1}=6$; $u_{m+1}-u_m=2\,500-1\,000=1\,500$.\\
		Mốt của mẫu số liệu ghép nhóm là 
		\[M_o=1\,000+\dfrac{7-0}{(7-0)+(7-6)} \cdot  1\,500=2\,312{,}5.\]
\end{listEX}}
\end{ex}

% nguồn: Ok_TuanLuan_Toan11_Lan2
\begin{ex}%[1D5H1-3]%[ma:b19ea8a2]
	An tìm hiểu hàm lượng chất béo (đơn vị: g) có trong $100$ g  mỗi loại thực phẩm. Sau khi thu thập dữ liệu về $60$ loại thực phẩm, An lập được bảng thống kê
	\begin{center}
		\begin{tabular}{|c|c|c|c|c|c|c|}
			\hline 
			Hàm lượng chất béo (g) & $[2;6)$ & $[6;10)$ & $[10;14)$ & $[14;18)$ &  $[18;22)$ & $[22;26)$ \\
			 \hline
			 Tần số & $2$ & $6$ & $10$ & $13$ & $16$ & $13$ \\
			\hline 
		\end{tabular}
	\end{center}
	\begin{enumerate}
		\item Xác định giá  trị trung bình của mẫu số liệu.
		\item Tìm mốt của  mẫu số liệu.
	\end{enumerate}
	 \loigiai{
		\begin{center}
			\begin{tabular}{|c|c|c|c|c|c|c|}
				 \hline
				Hàm lượng chất béo (g) & $[2;6)$ & $[6;10)$ & $[10;14)$ & $[14;18)$ & $[18;22)$ &  $[22;26)$ \\
				\hline 
				Giá trị  đại diện & $4$ & $8$ & $12$ & $16$ & $20$ & $24$ \\
				\hline 
				Tần số & $2$ & $6$ & $10$ &  $13$ & $16$ & $13$ \\
				\hline 
			\end{tabular}
		\end{center}
		\begin{enumerate}
			 \item Giá trị trung bình: $\bar{x}=\dfrac{4\cdot 2+8\cdot 6+12\cdot 10+16\cdot 13+20\cdot 16+24\cdot 13}{60}=16{,}9$.
			\item Mốt của mẫu số  liệu $M_0=18+\dfrac{3}{3+3}\cdot 4=20$.
		\end{enumerate}
	} 
\end{ex}

% nguồn: Ok_TuanLuan_Toan11_Lan2
\begin{ex}%[1D5H1-3]%[1D5H1-4]%[ma:e0df54f2]
	Bảng sau thống kê số lượt chở khách  mỗi ngày của một lái xe taxi trong 30 ngày.
	\begin{center}
		\begin{tabular}{|*{15}{c|}}
			\hline 
			15 & 13 & 7  & 5 & 18 & 13 & 11 & 9 & 10 & 8 & 14 & 11 & 16 & 10 & 9 \\
			\hline 
			13 & 11 & 12 & 13 & 15 & 12 & 13 & 6 & 8  & 17 & 13 & 6 & 18 & 12 & 13 \\
			\hline 
		\end{tabular}
	\end{center}
	\begin{listEX}
		\item Hãy tính số trung bình và mốt của  mẫu số liệu trên.
		\item Hãy lập bảng tần số  ghép nhóm của mẫu số liệu trên với nhóm đầu tiên là $[4{,}5; 7{,}5)$.
		\item Hãy ước lượng số  trung bình và mốt của mẫu số liệu ghép nhóm đó.
	\end{listEX}
	 \loigiai{
		\begin{listEX}
			\item Số trung bình  mẫu số liệu là $11{,}7$. Mốt của mẫu số liệu là $13$.
			\item Bảng tần số ghép  nhóm của mẫu số liệu đã cho như sau:
			\item Số trung  bình của mẫu số liệu ghép nhóm là $11{,}6$.
			\begin{center}
				\begin{tabular}{|*{15}{c|}}
					\hline 
					Số lượt khách & $[4{,}5; 7{,}5)$& $[7{,}5;  10{,}5)$& $[10{,}5; 13{,}5)$& $[13{,}5; 16{,}5)$ & $[16{,}5; 19{,}5)$ \\
					\hline 
					Số ngày & 4 & 6  & 13 & 4 & 3 \\
					 \hline
				\end{tabular}
			\end{center}
			Mốt của mẫu số liệu ghép nhóm  là $\dfrac{189}{16} \approx 11{,}8$.
		\end{listEX}
	 }
\end{ex}

% nguồn: Ok_TuanLuan_Toan11_Lan2
\begin{ex}%[1D5H1-3]%[1D5H1-4]%[ma:66460868]
	Quãng đường (km) các cầu thủ (không tính thủ  môn) chạy trong một trận bóng đá tại giải ngoại hạng Anh được cho trong bảng thống kê sau:
	\begin{center}
		\begin{tabular}{|*6{c|}}
			 \hline
			Quãng đường & $[2; 4)$ &$[4; 6)$ & $[6; 8)$ & $[8; 10)$ & $[10; 12)$  \\
			\hline 
			Số cầu thủ & 2 & 5 & 6 & 9 & 3  \\
			 \hline
		\end{tabular}
	\end{center}
	\begin{listEX}
		 \item Tính quãng đường trung bình một cầu thủ chạy trong trận đấu này.
		\item Tính mốt của mẫu số liệu và giải thích ý nghĩa của giá trị thu được. 
	\end{listEX}
	\loigiai{ 
		\begin{listEX}
			\item Ta có giá trị đại diện của  các nhóm
			\begin{center}
				\begin{tabular}{|*6{c|}}
					 \hline
					\makecell{Quãng đường\\(giá trị đại diện)} & $3$ &$5$  & $7$ & $9$ & $11$ \\
					\hline 
					Số cầu thủ & 2 & 5 & 6  & 9 & 3 \\
					\hline 
				\end{tabular}
			\end{center}
			 Quãng đường trung bình một cầu thủ chạy trong trận đấu này là
			\[\overline{x}=\dfrac{2\cdot 3 +5\cdot 5+6\cdot 7+9\cdot 9+3\cdot 11}{2+5+6+9+3}= 7{,}48\,\text{(km)}.\] 
			\item Nhóm chứa mốt  là $[8; 10)$. Mốt là $M_{\circ}=8+\dfrac{(9-6)}{(9-6)+(9-3)} \cdot 2 \approx 8{,}67$. \\
			Số cầu thủ chạy khoảng 8{,}67 km là nhiều  nhất.
		\end{listEX}
	} 
\end{ex}

% nguồn: Ok_TuanLuan_Toan11_Lan2
\begin{ex}%[1D5H1-3]%[1D5H1-4]%[ma:033b2db7]
	Tuổi thọ (năm) của 50 bình ắc quy  ô tô được cho như sau:
	\begin{center}
		\begin{tabular}{|*7{c|}}
			\hline 
			Tuổi thọ (năm) & $[2; 2{,}5)$ & $[2{,}5; 3)$&$[3; 3{,}5)$& $[3{,}5; 4)$ & $[4; 4{,}5)$&$[4{,}5;  5)$ \\
			 \hline
			Tần số &  4 & 9 & 14 & 11 & 7 & 5 \\
			\hline 
		\end{tabular}
	\end{center}
	\begin{listEX}
		\item Xác định mốt và giải thích ý  nghĩa.
		\item Tính tuổi thọ trung bình của $50$ bình ắc quy ô tô này. 
	\end{listEX}
	 \loigiai{
		\begin{listEX}
			\item Tần số lớn nhất là $14$ nên nhóm chứa  mốt là nhóm $\left[3;3{,}5 \right)$. \\
			Ta có, $j=3$,  $a_3=3$, $m_3=14$, $m_2=9$, $m_4=11$, $h=0{,}5$. Do đó mốt của mẫu số liệu ghép nhóm là $M_o=3+\dfrac{14-9}{(14-9)+(14-11)} \cdot 0{,}5=3{,}3125$.\\
			Ý nghĩa: Tuổi thọ của bình ắc quy  ô tô khoảng $3{,}3125$ năm là nhiều nhất hay tuổi thọ chủ yếu của bình ắc quy ô tô khoảng $3{,}3125$ năm.
			\item Trong mỗi khoảng tuổi  thọ, giá trị đại diện là trung bình cộng của giá trị hai đầu mút nên ta có bảng sau:
			\begin{center}
				\begin{tabular}{|*7{c|}}
					 \hline
					Tuổi thọ (năm) & $2{,}25$ & $2{,}75$ &  $3{,}25$ & $3{,}75$ & $4{,}25$ & $4{,}75$ \\
					\hline 
					Tần  số & $4$ & $9$ & $14$ & $11$ & $7$ & $5$ \\
					 \hline
				\end{tabular}
			\end{center}
			Tổng số ắc quy ô tô là $50$. Tuổi thọ trung bình của $50$ ắc quy ô tô  này là
			 \[\overline{x}=\dfrac{4.2{,}25+9.2{,}75+14.3{,}25+11.3{,}75+7.4{,}25+5.4{,}75}{50}=3{,}48.\]
	\end{listEX}}
\end{ex}

% nguồn: Ok_TuanLuan_Toan11_Lan2
\begin{ex}%[1D5H1-3]%[1D5H1-4]%[ma:c100f093]
	Bảng sau thống kê chiều cao (đơn vị: cm) của một số cây giống sau khi nảy mầm  được $2$ tuần.
	\begin{center}
		\begin{tabular}{|*{6}{c|}}
			 \hline
			Chiều cao (cm) & $[6{,}2; 6{,}7)$ &$[6{,}7; 7{,}2)$ &$[7{,}2; 7{,}7)$ &$[7{,}7; 8{,}2)$  &$[8{,}2; 8{,}7)$ \\
			\hline 
			Số cây & 10 & 21  & 28 & 12 & 9 \\
			 \hline
		\end{tabular}
	\end{center}
	Hãy ước lượng  chiều cao trung bình và mốt của mẫu số liệu ghép nhóm trên.
	\loigiai{ 
		Số  trung bình của mẫu số liệu ghép nhóm là khoảng $7{,}38$.\\
		Mốt của mẫu số liệu ghép nhóm là khoảng $7{,}35$. 
	} 
\end{ex}

% nguồn: Ok_TuanLuan_Toan11_Lan2
\begin{ex}%[1D5H1-3]%[1D5H1-4]%[ma:9b0b5bad]
	Một bưu tá  thống kê lại số bưu phẩm gửi đến một cơ quan mỗi ngày trong tháng 6/2022 ở bảng sau:
	\begin{center}
		\begin{tabular}{*{10}{c}}
			30 &  32 & 28 & 34 & 37 & 26 & 44 & 24 & 22 & 38 \\
			34 & 20 & 30  & 27 & 28 & 34 & 38 & 32 & 42 & 39 \\
			43 & 42 &  32 & 26 & 36 & 32 & 37 & 24 & 29 & 32
		\end{tabular}
	\end{center}
	\begin{listEX}
		\item Tính số trung bình và mốt của mẫu số liệu  trên.
		\item Tổng hợp lại số liệu trên vào bảng tần số  ghép nhóm theo mẫu sau:
		\begin{center}
			\begin{tabular}{|*6{c|}}
				 \hline
				Số bưu phẩm &  $[20; 24]$ &$[25; 29]$ & $[30; 34]$ & $[35; 39]$ & $[40; 44]$ \\
				 \hline
				Số ngày & $?$ & $?$ & $?$ & $?$  & $?$ \\
				 \hline
			\end{tabular}
		\end{center}
		\item Hãy ước lượng số trung bình và mốt của mẫu số liệu ghép  nhóm trên.
	\end{listEX}
	 \loigiai{
		\begin{listEX}
			\item Số trung bình của mẫu số liệu là 32{,}4. Mốt của mẫu  số liệu là 32.
			\item Bảng tần  số ghép nhóm.
			\begin{center}
				\begin{tabular}{|*6{c|}}
					 \hline
					Số bưu phẩm & $[20;  24]$ & $[25; 29]$ & $[30; 34]$ & $[35; 39]$ & $[40; 44]$ \\
					 \hline
					Số ngày & 4 & 6 & 10 & 6 & 4  \\
					 \hline
				\end{tabular}
			\end{center}
			\item Do số bưu phẩm là  số nguyên nên ta hiệu chỉnh lại bảng tần số ghép nhóm như sau:
			\begin{center}
				\begin{tabular}{|*6{c|}}
					\hline 
					Số bưu phẩm & $[19{,}5;  24{,}5)$ &$[24{,}5; 29{,}5)$ &$[29{,}5; 34{,}5)$ &$[34{,}5; 39{,}5)$ &$[39{,}5; 44{,}5)$ \\
					\hline 
					Giá trị đại diện & 22 & 27  & 32 & 37 & 42 \\
					 \hline
					Số ngày & 4 & 6 & 10 & 6 & 4 \\ 
					\hline 
				\end{tabular}
			\end{center}
			Số trung bình của mẫu số liệu ghép  nhóm là
			\[\overline{x}=\dfrac{4\cdot 22+6\cdot  27+10\cdot 32+6\cdot 37+4\cdot 42}{30}=32.\]
			 Nhóm chứa mốt của mẫu số liệu trên là nhóm $[29{,}5; 34{,}5)$.\\
			Do đó: $u_m=29{,}5$;  $n_{m-1}=6$; $n_m=10$; $n_{m+1}=6$; $u_{m+1}-u_m=34{,}5-29{,}5=5$.\\
			Mốt của mẫu số liệu ghép nhóm là $M_o=29{,}5+\dfrac{10-6}{(10-6)+(10-6)} \cdot  5=32$.
		\end{listEX}
	} 
\end{ex}

% nguồn: Ok_TuanLuan_Toan11_Lan2
\begin{ex}%[1D5H1-3]%[ma:01496cfa]
	Mẫu số liệu sau ghi lại cân nặng của 30 bạn học  sinh (đơn vị: kilogam)
	\begin{center}
		\begin{tabular}{|c|c|c|c|c|c|c|c|c|c|}
			 \hline
			$17$ & $40$ & $39$ & $40{,}5$ & $42$ & $51$ & $41{,}5$ & $39$ & $41$  & $30$ \\
			\hline 
			$40$ & $42$ & $40{,}5$ & $39{,}5$ & $41$ & $40{,}5$ & $37$ & $39{,}5$ & $40$  & $41$ \\
			 \hline
			$38{,}5$ & $39{,}5$ & $40$ & $41$ & $39$ & $40{,}5$ & $40$ &  $38{,}5$ & $39{,}5$ & $41{,}5$ \\
			\hline 
		\end{tabular}
	\end{center}
	\begin{enumerate}
		\item Lập bảng tần số ghép nhóm cho mẫu số liệu trên có  tám nhóm ứng với tám nửa khoảng
		$$\left[15;20\right),\left[20;25\right),\left[25;30\right),\left[30;35\right),\left[35;40\right),\left[40;45\right),\left[45;50\right),\left[50;55\right).$$ 
		\item  Xác định số trung bình cộng của mẫu số liệu ghép nhóm trên.
		 \item Mốt của mẫu số liệu ghép nhóm trên là bao nhiêu.
	\end{enumerate}
	\loigiai{ 
		\begin{enumerate}
			\item Bảng tần số ghép nhóm cho mẫu số liệu là 
			\begin{center}
				\begin{tabular}{|c|c|c|c|c|c|c|c|c|}
					 \hline
					Cân  nặng(m) & [$15$;$20)$ & [$20$;$25)$ & [$25$;$30)$ & [$30$;$35)$ & [$35$;$40)$ & [$40$;$45)$ & [$45$;$50)$ & [$50$;$55)$ \\
					 \hline
					Giá trị đại  diện & $17{,}5$ & $22{,}5$ & $27{,}5$ & $32{,}5$ & $37{,}5$ & $42{,}5$ & $47{,}5$ & $52{,}5$ \\
					 \hline
					Số học sinh  & $1$ & $0$ & $0$ & $1$ & $10$ & $17$ & $0$ & $1$ \\
					\hline 
				\end{tabular}
			\end{center}
			\item Số  trung bình cộng $\overline{x}=\dfrac{17{,}5+32{,}5+37{,}5\cdot 10+42{,}5\cdot 17+52{,}5}{30}=25{,}8$.
			\item Mốt  của mẫu số liệu $$M_0=40+\dfrac{7}{7+17}\cdot5\approx 41{,}5.$$
		\end{enumerate}
	} 
\end{ex}

% nguồn: Ok_TuanLuan_Toan11_Lan2
\begin{ex}%[1D5H1-3]%[ma:b8b13f1d]
	Phỏng vấn một số học sinh khối $11$ về thời gian (giờ) ngủ  của một buổi tối, thu được bảng số liệu ở bên
	\begin{center}
		\begin{tabular}{|c|c|c|c|}
			 \hline
			Thời gian & Số học sinh nam & Số học sinh nữ & Số học sinh khối $11$  \\
			\hline 
			 $[4;5)$ & $6$ & $4$ & $10$ \\
			\hline 
			$[5;6)$ & $10$ & $8$  & $18$ \\
			\hline 
			 $[6;7)$ & $13$ & $10$ & $23$ \\
			\hline 
			$[7;8)$ & $9$  & $11$ & $20$ \\
			\hline 
			$[8;9)$ & $7$ & $8$ & $15$ \\ 
			 \hline
		\end{tabular}
	\end{center}
	\begin{enumerate}
		\item So sánh thời gian ngủ trung bình của các bạn học sinh  nam và nữ.
		\item Hãy cho  biết $75\%$ học sinh khối $11$ ngủ ít nhất bao nhiêu giờ?
	\end{enumerate}
	\loigiai{ 
		\begin{enumerate}
			\item Trong mỗi khoảng thời gian,  giá trị đại diện là trung bình cộng của giá trị hai đầu mút nên ta có bảng sau
			\begin{center}
				\begin{tabular}{|c|c|c|}
					 \hline
					Thời gian & Số học sinh nam & Số  học sinh nữ \\
					\hline 
					$4{,}5$ & $6$ &  $4$ \\
					 \hline
					$5{,}5$ & $10$  & $8$ \\
					 \hline
					$6{,}5$ & $13$  & $10$ \\
					\hline 
					$7{,}5$ & $9$ & $8$ \\ 
					\hline 
					$8{,}5$  & $7$ & $8$ \\
					 \hline
				\end{tabular}
			\end{center}
			 Tổng số các bạn nam là $n_1=6+10+13+9+7=45$.\\
			Thời gian ngủ trung bình của các bạn học sinh nam là \[\overline{x}_1=\dfrac{6\cdot4,5+10\cdot5,5+13\cdot6,5+9\cdot7,5+7\cdot8,5}{45} \approx  6{,}52.\]
			Tổng số các bạn  nữ là $n_2=4+8+10+11+8=41$.\\
			Thời gian ngủ trung bình của các bạn học sinh  nữ là \[\overline{x}_2=\dfrac{4\cdot4,5+8\cdot5,5+10\cdot6,5+11\cdot7,5+8\cdot8,5}{41} \approx 6{,}77.\]
			Vì $6{,}52 < 6{,}77$ nên thời gian ngủ trung bình của các học  sinh nam ít hơn các học sinh nữ.
			\item  Ta có
			\begin{center}
				\begin{tabular}{|c|c|c|c|}
					 \hline
					Thời gian & Số học sinh nam & Số học sinh nữ & Số học sinh khối  11 \\
					\hline 
					$[4;5)$ & $6$ & $4$ & $10$ \\ 
					\hline 
					$[5;6)$ & $10$ & $8$ &  $18$ \\
					\hline 
					$[6;7)$  & $13$ & $10$ & $23$ \\
					\hline 
					$[7;8)$ &  $9$ & $11$ & $20$ \\
					 \hline
					$[8;9)$ & $7$ & $8$  & $15$ \\
					\hline 
				\end{tabular}
			\end{center}
			Tổng số học sinh khối $11$ được khảo sát  là $n=45+41=86$.\\
			Gọi $x_1$, $x_2$, $x_3$, $\ldots$, $x_{86}$ là thời gian ngủ của các học sinh khối  $11$ được khảo sát và giả sử dãy này đã sắp xếp theo thứ tự tăng dần. Khi đó trung vị của mẫu số liệu là $\dfrac{x_{43}+x_{44}}{2}$.\\
			Do đó, tứ phân vị thứ ba $Q_3=x_{65}$. Vì $x_{65}$ thuộc nhóm  $[7;8)$ nên nhóm này chứa $Q_3$. Do đó, $p=4$; $a_4=7$; $m_4=20$; $m_1+m_2+m_3=10+18+23=51$; $a_5-a_4=8-7=1$ và ta có
			\[Q_3=7+\dfrac{\dfrac{3\cdot86}{4}-51}{20}\cdot1=7{,}675 
			 \]
			Tứ phân vị thứ ba $Q_3$ chia mẫu số liệu thành $2$ phần, phần dưới chiếm $75\%$ số liệu  của mẫu và phần trên chiếm $25\%$ số liệu của mẫu.
			Vậy  $75\%$ học sinh khối $11$ ngủ ít nhất $7{,}675$ giờ.
		\end{enumerate}
	} 
\end{ex}

% nguồn: Ok_TuanLuan_Toan11_Lan2
\begin{ex}%[1D5H1-3]%[ma:52518575]
	Tìm  hiểu thời gian xem ti vi trong tuần trước (đơn vị: giờ) của một số học sinh thu được kết quả sau:
	\begin{center}
		\begin{tabular}{|c|c|c|c|c|c|}
			 \hline
			Thời gian (giờ) & $[0; 5)$ & $[5; 10)$ & $[10; 15)$ & $[15; 20)$ &  $[20; 25)$ \\
			 \hline
			Số học sinh & $8$ & $16$ & $4$ & $2$ & $2$  \\
			\hline 
		\end{tabular}
	\end{center}
	Tính thời gian xem ti vi trung bình trong  tuần trước của các bạn học sinh này.
	 \loigiai{
		Trong mỗi khoảng thời gian, giá trị đại  diện là trung bình cộng của giá trị hai đầu mút nên ta có bảng sau:
		\begin{center}
			\begin{tabular}{|c|c|c|c|c|c|}
				\hline 
				Thời gian (giờ) & $2{,}5$ & $7{,}5$ & $12{,}5$ &  $17{,}5$ & $22{,}5$ \\
				\hline 
				Số học sinh & $8$ &  $16$ & $4$ & $2$ & $2$ \\
				 \hline
			\end{tabular}
		\end{center}
		Tổng  số học sinh là $n=8+16+4+2+2=32$.\\
		Thời gian xem ti vi trung bình trong  tuần trước của các học sinh là \[\overline{x}=\dfrac{8\cdot 2{,}5+16\cdot 7{,}5+4\cdot 12{,}5+2\cdot 17{,}5+2\cdot 22{,}5}{32}=8{,}4375\text{ (giờ).}\]
	} 
\end{ex}

% nguồn: Ok_TuanLuan_Toan11_Lan2
\begin{ex}%[1D5H1-3]%[ma:d3021ac7]
	Khảo sát thời  gian tự học của các học sinh trong lớp theo mẫu bên.
	\begin{enumEX}{2}
		\item! Hãy lập bảng thống kê cho mẫu số liệu ghép nhóm thu được. 
		\item! Có thể tính chính xác thời gian tự học trung bình của các học sinh trong  lớp không?
		\item! Có cách nào tính gần đúng thời gian tự học trung bình của các học sinh  trong lớp dựa trên mẫu số liệu ghép nhóm này không?
	\end{enumEX}
	 \loigiai{
		\begin{enumerate}[a)]
			\item Giả sử lớp $11A$ có $30$ học sinh và sau  khi khảo sát, ta có được bảng thống kê như sau:
			\begin{center}
				\begin{tabular}{|c|c|c|c|c|}
					\hline 
					Thời gian (giờ) &  Dưới $1{,}5$ giờ & $[1{,}5; 3)$ & $[3; 4{,}5)$ & Từ $4{,}5$ giờ trở lên \\
					\hline 
					Số học sinh & $5$  & $15$ & $8$ & $2$ \\
					 \hline
				\end{tabular}
			\end{center}
			\item Ta không thể tính chính xác thời gian tự học trung bình của các học sinh trong lớp vì không có mẫu số liệu cụ thể về thời gian tự học của từng  học sinh.
			\item Có thể tính gần đúng thời gian tự học trung bình của các học sinh trong lớp dựa trên mẫu số liệu ghép nhóm bằng cách chọn thời gian đại diện cho mỗi nhóm, sau đó sử dụng tần số tương ứng để  tính số trung bình, cụ thể:
			\begin{itemize}
				\item Thời gian tự học dưới $1{,}5$ giờ, ta chọn giá  trị đại diện là $0{,}75$ giờ, tần số tương ứng là 5 .
				\item Thời gian tự học  từ $1{,}5$ đến dưới $3$ giờ, ta chọn giá trị đại diện là $\dfrac{1{,}5+3}{2}=2{,}25$, tần số tương ứng là $15$.
				\item Thời gian tự học từ $3$ đến dưới $4{,}5$  giờ, ta chọn giá trị đại diện là $\dfrac{3+4{,}5}{2}=3{,}75$, tần số tương ứng là $8$.
				 \item Thời gian tự học là từ $4{,}5$ giờ trở lên, ta chọn giá trị đại diện là $5{,}25$, tần số tương ứng là $2$ .
			\end{itemize}
			Số trung bình là $\overline{x}=\dfrac{5\cdot  0{,}75+15\cdot 2{,}25+8\cdot 3{,}75+2\cdot 5{,}25}{30}=2{,}6$.\\
			Vậy thời gian tự học  trung bình của học sinh lớp $11A$ xấp xỉ khoảng $2{,}6$ giờ.
		\end{enumerate}
	} 
\end{ex}

% nguồn: Ok_TuanLuan_Toan11_Lan2
\begin{ex}%[1D5H1-3]%[1D5H1-4]%[ma:799132f0]
	Thống kê lại thu nhập trong một tháng  của nhân viên hai công ty $A$ và $B$ (đơn vị: triệu đồng) được thể hiện trong biểu đồ dưới đây.
	\begin{center}
		\begin{tikzpicture}[line cap=round,line join=round,font=\footnotesize,>=stealth,scale=1]
				 \draw[->] (0,0)--(0,6) node[left] {(số nhân viên)};
				\draw[->] (0,0)--(11,0) node[below] {(triệu đồng)}; 
				\draw[black!30!white] (11,1) -- (0,1)  node[left] {$10$};
				\draw[black!30!white]  (11,2) -- (0,2) node[left] {$20$};
				 \draw[black!30!white] (11,3) -- (0,3) node[left] {$30$};
				\draw[black!30!white] (11,4) -- (0,4) node[left] {$40$}; 
				\draw[black!30!white] (11,5)  -- (0,5) node[left] {$50$};
				\draw[fill=red!40!white] (0,0)  rectangle (1,1.2);
				\draw[fill=red!40!white]  (2,0) rectangle (3,4.5);
				 \draw[fill=red!40!white] (4,0) rectangle (5,1.6);
				\draw[fill=red!40!white] (6,0) rectangle (7,1.4); 
				\draw[fill=red!40!white] (8,0) rectangle (9,0.7); 
				\draw[fill=red!5!white]  (1,0) rectangle (2,1.6);
				\draw[fill=red!5!white]  (3,0) rectangle (4,4.0);
				 \draw[fill=red!5!white] (5,0) rectangle (6,2.9);
				\draw[fill=red!5!white]  (7,0) rectangle (8,1.5);
				\draw[fill=red!5!white] (9,0) rectangle  (10,1.2);
				\path (1,0) node[below] {$[5; 9)$}; 
				 \path (3,0) node[below] {$[9; 13)$};
				\path (5,0)  node[below] {$[13; 17)$};
				\path (7,0) node[below] {$[17;  21)$};
				\path (9,0) node[below] {$[21;  25)$};
				 \path (0.5,1.2) node[above] {$12$};
				 \path (1.5,1.6) node[above] {$16$};
				\path  (2.5,4.5) node[above] {$45$};
				\path (3.5,4.0)  node[above] {$40$};
				\path (4.5,1.6) node[above]  {$16$};
				\path  (5.5,2.9) node[above] {$29$};
				\path (6.5,1.4) node[above] {$14$}; 
				\path (7.5,1.5) node[above] {$15$}; 
				\path (8.5,0.7) node[above] {$7$}; 
				\path  (9.5,1.2) node[above] {$12$};
				\path (5.5,6) node {\makecell{\textbf{Thu nhập trong một tháng của nhân viên}\\\textbf{hai  công ty \textit{A} và \textit{B}}}};
				\draw[fill=red!40!white] (12.3,3.2)  rectangle (11.7,3.8);
				 \path (12.5,2.5) node[right] {\textbf{Công ty \textit{B}}};
				\draw[fill=red!5!white] (12.3,2.2) rectangle  (11.7,2.8);
				 \path (12.5,3.5) node[right] {\textbf{Công ty \textit{A}}};
			\end{tikzpicture}
	\end{center}
	Hãy so sánh thu nhập trung bình của nhân viên hai công ty theo số  trung bình và mốt của mẫu số liệu ghép nhóm.
	 \loigiai{
		Ta lập  bảng thống kê thu nhập của nhân viên hai công ty như sau:
		\begin{center}
			\begin{tabular}{|*{4}{c|}}
				\hline 
				\makecell{Thu nhập  \\
				(đơn vị:  triệu đồng)} & Giá trị đại diện & Số nhân viên công ty A & Số nhân viên công ty B \\
				\hline 
				$[5; 9)$ & 7 & 12 & 16  \\
				\hline 
				$[9; 13)$ &  11 & 45 & 40 \\
				\hline 
				 $[13; 17)$ & 15 & 16 & 29 \\
				 \hline
				$[17; 21)$ & 19 & 14 & 15 \\ 
				\hline 
				$[21; 25)$ & 23 &  7 & 12 \\
				\hline 
			\end{tabular}
		\end{center}
		Số trung bình thu thập trong một tháng của  nhân viên công ty $A$ là
		\[\overline{x}_A=\dfrac{12\cdot7+45\cdot11+16\cdot15+14\cdot19+7\cdot23}{94}\approx 13{,}26\,\text{(triệu đồng)}.\] 
		Số trung bình thu thập trong một tháng của nhân viên công ty  $B$ là
		\[\overline{x}_B=\dfrac{16\cdot7+40\cdot11+29\cdot15+15\cdot19+12\cdot23}{112}\approx 13{,}82\,\text{(triệu  đồng)}.\]
		Vậy nếu so sánh theo số trung bình của mẫu số liệu ghép nhóm thì thu nhập của nhân  viên công ty $A$ thấp hơn thu nhập của nhân viên công ty $B$.\\
		Ta ước  lượng mốt của thu nhập của nhân viên công ty $A$.\\
		Nhóm  mốt của mẫu số liệu là nhóm $[9; 13)$.\\
		Do đó $u_m=9$;  $n_{m-1}=12$; $n_m=45$; $n_{m+1}=16$; $u_{m+1}-u_m=13-9=4$.\\
		Mốt của mẫu số liệu ghép  nhóm là
		\[M_o=9+\dfrac{45-12}{(45-12)+(45-16)} \cdot 4=\dfrac{345}{31} \approx 11{,}13.\] 
		Ta ước lượng mốt của thu nhập của nhân  viên công ty $B$.\\
		Nhóm mốt của  mẫu số liệu là nhóm $[9; 13)$.\\
		Do đó $u_m=9$;  $n_{m-1}=16$; $n_m=40$; $n_{m+1}=29$; $u_{m+1}-u_m=13-9=4$.\\
		Mốt của  mẫu số liệu ghép nhóm là
		\[M_o=9+\dfrac{40-16}{(40-16)+(40-29)} \cdot 4=\dfrac{411}{35} \approx 11{,}74.\] 
		Vậy nếu so sánh theo mốt của mẫu số liệu ghép nhóm, thu nhập của nhân viên công ty $A$ thấp hơn thu  nhập của nhân viên công ty $B$.
	 }
\end{ex}

% nguồn: Ok_TuanLuan_Toan11_Lan2
\begin{ex}%[1D5H1-3]%[ma:6cd01148]
	 Kết quả khảo sát cân nặng của $20$ quả cam Canh ở mỗi lô hàng 1 và lô hàng 2 được cho ở bảng sau:
	\begin{center}
		\begin{tabular}{|*6{c|}}
			\hline 
			Cân nặng (g) & $[100; 110)$ &$[110; 120)$ &$[120; 130)$ & $[130; 140)$ & $[140;  150)$ \\
			 \hline
			Số quả cam Canh ở lô  hàng 1 & 1 & 4 & 5 & 4 & 6 \\
			 \hline
			Số quả cam Canh ở lô hàng 2 & 2 & 3 & 6 & 4  & 5 \\
			 \hline
		\end{tabular}
	\end{center}
	\begin{listEX}
		\item Hãy  ước lượng cân nặng trung bình của mỗi quả cam Canh ở lô hàng 1 và lô hàng 2.
		\item Nếu so sánh theo số trung bình thì  cam Canh ở lô hàng nào nặng hơn?
	\end{listEX}
	\loigiai{ 
		Ta có bảng thống kê số lượng cam Canh theo giá trị  đại diện:
		\begin{center}
			\begin{tabular}{|*6{c|}}
				\hline 
				Cân nặng đại diện (g) & 105  & 115 & 125 & 135 & 145 \\
				\hline 
				Số quả cam Canh ở lô hàng 1 & 1 & 4 & 5 & 4 & 6 \\ 
				 \hline
				Số quả  cam Canh ở lô hàng 2 & 2 & 3 & 6 & 4 & 5 \\
				\hline 
			\end{tabular}
		\end{center}
		\begin{listEX}
			\item  Cân nặng trung bình của mỗi quả cam Canh ở lô hàng 1 là
			 \[\dfrac{1\cdot 105+4\cdot 115+5\cdot 125+4\cdot 135+6\cdot 145}{20}=130\,\text{(g)}.\]
			Cân nặng trung bình của mỗi quả cam Canh ở lô  hàng 2 là
			\[\dfrac{2\cdot 105+3\cdot  115+6\cdot 125+4\cdot 135+5\cdot 145}{20}=128{,}5\,\text{(g)}.\]
			\item  Nếu so sánh theo số trung bình thì số cam Canh ở lô hàng 1 nặng hơn số cam Canh ở lô hàng 2.
		\end{listEX}
	 }
\end{ex}

% nguồn: Ok_TuanLuan_Toan11_Lan2
\begin{ex}%[1D5H1-3]%[1D5H1-4]%[ma:b15ac403]
	Quãng đường (km)  từ nhà đến nơi làm việc của $40$ công nhân một nhà máy được ghi lại như sau:
	\begin{center}
		\begin{tabular}{cccccccccc}
			$5$ & $3$  & $10$ & $20$ & $25$ & $11$ & $13$ & $7$ & $12$ & $31$ \\ $19$ & $10$ & $12$ & $17$ & $18$ & $11$ & $32$ & $17$ & $16$ & $2$ \\ $7$ & $9$ & $7$ & $8$ & $3$ & $5$ & $12$ & $15$ & $18$ & $3$ \\ $12$ & $14$ & $2$ & $9$ & $6$ & $15$ & $15$ & $7$ & $6$ & $12$.
		\end{tabular}
	\end{center}
	\begin{listEX}
		\item Ghép nhóm dãy số liệu trên thành các khoảng có độ rộng  bằng nhau, khoảng đầu tiên là $[0; 5)$. Tìm giá trị đại diện cho mỗi nhóm.
		\item Tính số trung bình của mẫu số liệu  không ghép nhóm và mẫu số liệu ghép nhóm. Giá trị nào chính xác hơn?
		\item Xác định nhóm chứa mốt của mẫu số  liệu ghép nhóm thu được.
	\end{listEX}
	 \loigiai{
		\begin{listEX}
			\item Giá trị nhỏ nhất của mẫu số liệu là $2$, giá trị lớn nhất của mẫu số  liệu là $32$, do đó khoảng biến thiên là $32-2=30$.
			\item Các nhóm có độ rộng bằng nhau và độ rộng của mỗi nhóm là $5$. Để cho thuận tiện, ta chia thành $7$ nhóm là các nhóm $\left[0;5 \right)$, $\left[5;10 \right)$, $\left[10;15 \right)$, $\left[15;20 \right)$, $\left[20;25 \right)$, $\left[25;30 \right)$, $\left[30;35 \right)$. Đếm số giá trị thuộc mỗi  nhóm, ta có mẫu số liệu ghép nhóm như sau:
			\begin{center}
				\begin{tabular}{|*8{c|}}
					\hline 
					Quãng đường (km) &  $[0; 5)$ & $[5; 10)$ &$[10; 15)$ & $[15; 20)$ & $[20; 25)$&$[25; 30)$ & $[30; 35)$\\
					 \hline
					Số công nhân  & 5 & 11 & 11 & 9 & 1 & 1 & 2 \\
					 \hline
				\end{tabular}
			\end{center}
			Giá trị đại diện cho mỗi nhóm là trung bình của hai đầu mút của nhóm. Ta có bảng giá trị đại diện  như sau:
			\begin{center}
				\begin{tabular}{|*8{c|}}
					 \hline
					 \makecell{Quãng đường (km)\\
					(giá trị đại diện)} & $2{,}5$ & $7{,}5$ & $12{,}5$ & $17{,}5$ & $22{,}5$ & $27{,}5$ &  $32{,}5$\\
					\hline 
					Số công nhân & $5$  & $11$ & $11$ & $9$ & $1$ & $1$ & $2$ \\
					\hline 
				\end{tabular}
			\end{center}
			\item Số trung bình của mẫu  số liệu ghép nhóm là \[\overline{x_g}=\dfrac{5.2{,}5+11.7{,}5+11.12{,}5+9.17{,}5+1.22{,}5+1.27{,}5+2.32{,}5}{40}=12{,}625.\]
			Ta có $\begin{aligned}
			 5+3+10+20+25+11+13+7+12+31+19+10+12+17& \\
			 {}+18+11+32+17+16+2+7+9+7+8+3+5+12+15& \\
			 {}+18+3+12+14+2+9+6+15+15+7+6+12&=476.
		\end{aligned}$\\
		Số trung bình của mẫu  số liệu không ghép nhóm là $\overline{x}=\dfrac{476}{40}=11{,}9$.\\
		Giá trị trung bình của mẫu số liệu không ghép  nhóm chính xác hơn vì nó là giá trị của mẫu số liệu gốc.
		\item Tần số lớn nhất trong bảng tần số của mẫu số liệu ghép nhóm là $11$. Do đó, nhóm chứa mốt của mẫu số liệu ghép nhóm là các  nhóm $\left[5;10 \right)$ và $\left[10;15 \right)$.
	\end{listEX}
} 
\end{ex}

% nguồn: Ok_TuanLuan_Toan11_Lan2
\begin{ex}%[1D5H1-3]%[ma:f9720d35]
	Thống kê số lần đi học muộn trong học kì của các  bạn trong lớp, Nam thu được kết quả sau:
	\begin{center}
		\begin{tabular}{|*6{c|}}
			 \hline
			Số lần đi muộn & $0-2$ & $3-5$ & $6-8$ & $9-11$  & $12-14$ \\
			 \hline
			Số học sinh & 23 & 8 & 5 & 3 &  1 \\
			 \hline
		\end{tabular}
	\end{center}
	Trung bình mỗi học sinh trong lớp đi muộn bao nhiêu buổi trong học kì? 
	\loigiai{ 
		Ta có mẫu số  liệu ghép nhóm
		\begin{center}
			\begin{tabular}{|*6{c|}}
				\hline 
				Số lần đi muộn & $[-0{,}5;2{,}5)$ &  $[2{,}5;5{,}5)$ & $[5{,}5;8{,}5)$ & $[8{,}5;11{,}5)$ & $[11{,}5;14{,}5)$ \\
				\hline 
				Giá trị đại  diện & $1$ & $4$ & $7$ & $10$ & $13$ \\
				 \hline
				Số học sinh & 23 & 8 & 5  & 3 & 1 \\
				\hline 
			\end{tabular}
		\end{center}
		Trung bình trong học kì mỗi học sinh đi  muộn \[\overline{x}=\dfrac{1\cdot 23+4\cdot 8+7\cdot 5+10\cdot 3+13\cdot 1}{23+8+5+3+1}=3{,}325\,\text{(buổi)}.\]
	} 
\end{ex}

% nguồn: Ok_TuanLuan_Toan11_Lan2
\begin{ex}%[1D5H1-3]%[1D5H1-4]%[ma:03f9f5b2]
	Nhân ngày hội đọc sách, các học sinh của một trường trung học phổ thông mang sách cũ đến tặng thư viện trường và trao đổi với các bạn  học sinh khác. Bảng sau thống kê số sách cũ mà các bạn học sinh lớp 11B mang đến trường.
	\begin{center}
		\begin{tabular}{|*6{c|}}
			 \hline
			 Số sách &$[1; 3]$ & $[4; 6]$ & $[7; 9]$ & $[10; 12]$ &$[13; 15]$ \\
			 \hline
			Số học sinh & $5$ & $14$ & $10$  & $8$ & $3$ \\
			\hline 
		\end{tabular}
	\end{center}
	Hãy ước lượng số trung bình và mốt  của mẫu số liệu ghép nhóm trên.
	 \loigiai{
		Bảng số liệu ghép  nhóm hiệu chỉnh như sau:
		\begin{center}
			\begin{tabular}{|*6{c|}}
				\hline 
				Số sách & $[0{,}5; 3{,}5)$ & $[3{,}5; 6{,}5)$& $[6{,}5;  9{,}5)$ & $[9{,}5; 12{,}5)$& $[12{,}5; 15{,}5)$ \\
				 \hline
				Số học sinh & 5 & 14 & 10 & 8 &  3 \\
				\hline 
			\end{tabular}
		\end{center}
		Số trung bình  của mẫu số liệu ghép nhóm là $7{,}25$.
		Mốt của mẫu số liệu ghép nhóm là $\dfrac{145}{26} \approx 5{,}58$. 
	} 
\end{ex}

% nguồn: Ok_TuanLuan_Toan11_Lan2
\begin{ex}%[1D5H1-3]%[ma:7e12101d]
	Anh Văn ghi lại cự li $30$ lần ném lao của mình ở bảng sau (đơn vị:  mét)
	\begin{center}
		\begin{tabular}{|c|c|c|c|c|c|c|c|c|c|}
			\hline 
			 $72{,}1$ & $72{,}9$ & $70{,}2$ & $70{,}9$ & $72{,}2$ & $71{,}5$ & $72{,}5$ & $69{,}3$ & $72{,}3$ & $69{,}7$ \\
			 \hline
			 $72{,}3$ & $71{,}5$ & $71{,}2$ & $69{,}8$ & $72{,}3$ & $71{,}1$ & $69{,}5$ & $72{,}2$ & $71{,}9$ & $73{,}1$ \\
			\hline 
			$71{,}6$ & $71{,}3$ & $72{,}2$ & $71{,}8$ & $70{,}8$ & $72{,}2$ & $72{,}2$ &  $72{,}9$ & $72{,}7$ & $70{,}7$ \\
			\hline 
		\end{tabular}
	\end{center}
	\begin{enumerate}
		\item Tính cự li trung bình của mỗi  lần ném.
		\item Tổng hợp lại kết quả ném của anh Văn vào bảng tần số ghép nhóm theo  mẫu sau
		\begin{center}
			\begin{tabular}{|c|c|c|c|c|c|}
				 \hline
				Cự  li (m) & [$69{,}2$;$70)$ & [$70$;$70{,}8)$ & [$70{,}8$;$71{,}6)$ & [$71{,}6$;$72{,}4)$ & [$72{,}4$;$73{,}2)$ \\
				 \hline
				Số lần  &? &? &? &? &? \\
				\hline 
			\end{tabular}
		\end{center}
		\item Hãy ước lượng cự li trung bình mỗi lần ném từ bảng tần số ghép nhóm  trên.
		\item Khả năng anh Văn ném được khoảng bao nhiêu  mét là cao nhất?
	\end{enumerate}
	\loigiai{ 
		\begin{enumerate}
			\item Cự li trung bình của mỗi lần  ném là $71{,}6$ (m).
			\item Bảng tần số ghép  nhóm
			\begin{center}
				\begin{tabular}{|c|c|c|c|c|c|}
					 \hline
					Cự li (m) & [$69{,}2$;$70)$ & [$70$;$70{,}8)$ & [$70{,}8$;$71{,}6)$ &  [$71{,}6$;$72{,}4)$ & [$72{,}4$;$73{,}2)$ \\
					 \hline
					Số lần & $4$ & $2$ & $9$  & $10$ & $5$ \\
					\hline 
				\end{tabular}
			\end{center}
			\item  \,
			\begin{center}
				\begin{tabular}{|c|c|c|c|c|c|}
					 \hline
					Cự  li (m) & [$69{,}2$;$70)$ & [$70$;$70{,}8)$ & [$70{,}8$;$71{,}6)$ & [$71{,}6$;$72{,}4)$ & [$72{,}4$;$73{,}2)$ \\
					 \hline
					Giá trị đại diện  & $69{,}6$ & $70{,}4$ & $71{,}2$ & $72$ & $72{,}8$ \\
					 \hline
					Số lần & $4$ & $2$  & $9$ & $10$ & $5$ \\
					 \hline
				\end{tabular}
			\end{center}
			Cự li trung bình  mỗi lần ném xấp xỉ bằng $$\dfrac{69{,}6\cdot 4+70{,}4\cdot 2+71{,}2\cdot 9+72\cdot 10+72{,}8\cdot 5}{30}=71{,}5 \,(\text{m}).$$
			\item Nhóm chứa mốt của mẫu số liệu trên là  $\left[71{,}6;72{,}4\right)$.\\
			Do  đó $u_m=71{,}6$; $n_{m-1}=9$; $n_{m+1}=5$; $u_{m+1}-u_m=72{,}4-71{,}6=0{,}8$.\\
			Mốt của mẫu số liệu ghép nhóm  là
			$$M_0=71{,}6+\dfrac{10-9}{\left(10-9\right)+\left(10-5\right)} \cdot 0{,}8=71{,}7\,(\text{m}).$$ 
			Vậy khả năng anh Văn ném được $71{,}7$ m là cao nhất. 
		\end{enumerate}
	 }
\end{ex}

% nguồn: Ok_TuanLuan_Toan11_Lan2
\begin{ex}%[1D5H1-3]%[ma:899f5ec5]
	Các  bạn học sinh lớp 11A1 trả lời 40 câu hỏi trong một bải kiểm tra. Kết quả được thống kê ở bảng sau:
	\begin{center}
		\begin{tabular}{|*{6}{c|}}
			\hline 
			Số  câu trả lời đúng&$[16;21)$&$[21;26)$&$[26;31)$&$[31;36)$&$[36;41)$\\
			 \hline
			Số  học sinh&4&6&8&18&4\\
			\hline 
		\end{tabular}
	\end{center}
	\begin{listEX}
		\item Tính giá trị đại  diện $c_p$ với $1\le p\le 5$, của từng nhóm số liệu.
		\item Tính $n_1c_1+n_2c_2+n_3c_3+n_4c_4+n_5c_5$ 
		\item Tinh  $\overline{x}=\dfrac{n_1c_1+n_2c_2+n_3c_3+n_4c_4+n_5c_5}{40}$.
	\end{listEX}
	\loigiai{ 
		\begin{listEX}
			\item  $c_1=18{,}5$; $c_2=23{,}5$; $c_3=28{,}5$; $c_4=33{,}5$; $c_5=38{,}5$.
			 \item $n_1c_1+n_2c_2+n_3c_3+n_4c_4+n_5c_5=1200$.
			\item  $\overline{x}=\dfrac{n_1c_1+n_2c_2+n_3c_3+n_4c_4+n_5c_5}{40}=\dfrac{1\,200}{40}=30$.
		\end{listEX}
	 }
\end{ex}

% nguồn: Ok_TuanLuan_Toan11_Lan2
\begin{ex}%[1D5H1-3]%[1D5H2-3]%[ma:87f7c216]
	Kiểm tra điện lượng của một số viên pintiểu do một hãng sản xuất thu được kết  quả sau
	\begin{longtable}{|c|*{5}{c|}}
		 \hline
		Điện lượng (nghìn mAh)& [$0{,}9$;  $0{,}95$) & [$0{,}95$; $1$) & [$1$; $1{,}05$) & [$1{,}05$; $1{,}1$) & [$1{,}1$; $1{,}15$) \\
		\hline 
		Số viên pin  & $10$ & $20$ & $35$ & $15$ & $5$ \\
		 \hline
	\end{longtable}
	Hãy ước lượng số trung bình, mốt và tứ phân  vị của mẫu số liệu ghép nhóm trên.
	\loigiai{ 
		Số trung bình của mẫu số liệu ghép nhóm trên  là:
		\[\overline{x}=\dfrac{10\cdot 0{,}925+20\cdot 0{,}975+35\cdot 1{,}025+15\cdot 1{,}075+5\cdot 1{,}125}{10+20+35+15+5} \approx 1{,}016.\\ 
		\] 
		Nhóm chứa mốt của mẫu số liệu  trên là nhóm $\left[1{,}0; 1{,}05\right)$.\\
		Mốt của mẫu số ghép nhóm  trên là $M_o=1{,}0+\dfrac{35-20}{\left(35-20\right)+\left(35-15\right)}\cdot\left(1{,}05-1{,}0\right)\approx 1{,}021$.\\
		Gọi $x_1, x_2, \ldots, x_{85}$ là điện lượng của mỗi viên pintiểu xếp theo thứ tự không  giảm.\\
		Ta có $x_1, x_2, \ldots, x_{10} \in \left[0{,}9; 0{,}95\right)$; $x_{11}, x_{12},  \ldots, x_{30} \in \left[0{,}95; 1{,}0\right)$; $x_{31}, x_{32}, \ldots, x_{65} \in \left[1{,}0; 1{,}05\right)$;\\
		$x_{66},  x_{67}, \ldots, x_{80} \in \left[1{,}05; 1{,}1\right)$;
		$x_{81}, x_{82}, \ldots, x_{85} \in  \left[1{,}1; 1{,}15\right)$.\\
		Tứ phân vị thứ hai của  dãy số liệu $x_1, x_2, \ldots, x_{85}$ là $x_{43}$.\\
		 Vì $x_{43} \in \left[1{,}0; 1{,}05\right)$ nên tứ phân vị thứ hai của mẫu số liệu là
		\begin{align*}
			 Q_2=1{,}0+\dfrac{\dfrac{2\cdot85}{4}-\left(10+20\right)}{35} \left(1{,}05-1{,}0\right)\approx 1{,}018.
		\end{align*}
		Tứ phân vị thứ nhất của dãy số liệu $x_1, x_2,  \ldots, x_{85}$ là $x_{21}$ và $x_{22}$.\\
		Vì $x_{21}, x_{22} \in \left[0{,}95; 1{,}0\right)$ nên tứ phân vị thứ nhất  của mẫu số liệu là
		\begin{align*}
			Q_1=0{,}95+\dfrac{\dfrac{1\cdot85}{4}-10}{20} \left(1{,}0-0{,}95\right)\approx 0{,}978. 
		\end{align*}
		Tứ  phân vị thứ ba của dãy số liệu $x_1, x_2, \ldots, x_{85}$ là $x_{64}$ và $x_{65}$.\\
		Vì $x_{64}, x_{65} \in \left[1{,}0; 1{,}05\right)$ nên  tứ phân vị thứ ba của mẫu số liệu là
		\begin{align*}
			Q_3=1{,}0+\dfrac{\dfrac{3\cdot85}{4}-\left(10+20\right)}{35} \left(1{,}05-1{,}0\right)\approx  1{,}048.
		\end{align*}
	} 
\end{ex}

% nguồn: Ok_TuanLuan_Toan11_Lan2
\begin{ex}%[1D5H1-3]%[ma:24c5f11c]
	Số cuộc gọi điện thoại một người thực hiện mỗi ngày trong $30$ ngày được lựa chọn ngẫu nhiên được  thống kê trong bảng sau:
	\begin{center}
		\begin{tabular}{|c|c|c|c|c|c|}
			\hline Số cuộc gọi & $[3;5]$ & $[6;8]$  & $[9;11]$ & $[12;14]$ & $[15;17]$\\
			\hline 
			Số ngày & $5$ & $13$ & $7$ & $3$ &  $2$ \\
			\hline 
		\end{tabular}
	\end{center}
	\begin{enumerate}
		\item Tìm mốt của mẫu số liệu ghép nhóm trên. 
		\item Hãy dự đoán xem  khả năng người đó thực hiện bao nhiêu cuộc gọi mỗi ngày là cao nhất.
	\end{enumerate}
	 \loigiai{
		Hiệu  chỉnh bảng lại như sau:
		\begin{center}
			\begin{tabular}{|c|c|c|c|c|c|}
				\hline Số cuộc gọi  & $[2{,}5;5{,}5)$ & $[5{,}5;8{,}5)$ & $[8{,}5;11{,}5)$ & $[11{,}5;14{,}5)$ & $[14{,}5;17{,}5)$\\
				 \hline
				Số ngày & $5$ & $13$ & $7$ & $3$ & $2$  \\
				\hline 
			\end{tabular}
		\end{center}
		\begin{enumerate}
			\item Nhóm chứa mốt của mẫu số  liệu trên là nhóm $\left[5{,}5;8{,}5\right)$.\\
			Do đó $u_m=5{,}5$, $n_{m-1}=5$, $n_m=13$, $n_{m+1}=7$, $u_{m+1}-u_m=8{,}5-5{,}5=3$.\\ 
			Mốt của mẫu số  liệu ghép nhóm là
			 \[M_o=5{,}5+\dfrac{13-5}{\left(13-5\right)+\left(13-7\right)} \cdot 3=\dfrac{101}{14} \approx 7{,}2
			 \]
			\item Dựa vào kết quả trên ta có thể dự đoán rằng khả năng người đó  thực hiện $7$ cuộc gọi mỗi ngày là cao nhất.
		\end{enumerate}
	 }
\end{ex}

% nguồn: Ok_TuanLuan_Toan11_Lan2
\begin{ex}%[1D5H1-3]%[ma:33385cac]
	Tìm cân nặng  trung bình của học sinh lớp $11D$ cho trong bảng sau
	\begin{center}
		\begin{tabular}{|c|c|c|c|c|c|c|}
			\hline 
			$\textbf{Cân nặng}$ & $[40{,}5;45{,}5)$ & $[45{,}5;50{,}5)$ & $[50{,}5;55{,}5)$ & $[55{,}5;60{,}5)$ & $[60{,}5;65{,}5)$ &  $[65{,}5;70{,}5)$\\
			\hline 
			$\textbf{Số học sinh}$ & $10$ & $7$ & $16$ & $4$  & $2$ & $3$\\
			\hline 
		\end{tabular}
	\end{center}
	\loigiai{ 
		 Trong mỗi khoảng cân nặng, giá trị đại diện là trung bình cộng của giá trị hai đầu mút nên ta có bảng sau
		\begin{center}
			\begin{tabular}{|c|c|c|c|c|c|c|}
				\hline 
				$\textbf{Cân nặng (kg)}$ & $43$ & $48$ & $53$ & $58$ & $63$  & $68$\\
				\hline 
				$\textbf{Số học sinh}$ & $10$ & $7$ & $16$ & $4$  & $2$ & $3$\\
				\hline 
			\end{tabular}
		\end{center}
		Tổng số học  sinh là $n=42$.\\
		Cân nặng trung bình của học sinh lớp $11D$  là
		\begin{align*}
			\overline{x}=\dfrac{43\cdot 10+7\cdot 48+16\cdot 53+4\cdot 58+2\cdot  63+3\cdot 68}{42}\approx 51{,}81 (\text{kg}).
		\end{align*}
	 }
\end{ex}

% nguồn: Ok_TuanLuan_Toan11_Lan2
\begin{ex}%[1D5H1-3]%[ma:f17f3e0f]
	Tìm hiểu thời gian xem ti vi trong tuần trước (đơn vị: giờ) của một số học sinh thu  được kết quả sau
	\begin{center}
		\begin{tabular}{|c|c|c|c|c|c|}
			 \hline
			$\textbf{Thời gian (giờ)}$ & $[0;5)$ &  $[5;10)$ & $[10;15)$ & $[15;20)$ & $[20;25)$ \\
			\hline 
			$\textbf{Số học sinh}$ &  $8$ & $16$ & $4$ & $2$ & $2$ \\
			 \hline
		\end{tabular}
	\end{center}
	Tính thời gian xem ti vi trung bình trong tuần trước của các bạn  học sinh này.
	\loigiai{ 
		\begin{center}
			\begin{tabular}{|c|c|c|c|c|c|}
				 \hline
				$\textbf{Thời gian (giờ)}$ & $[0;5)$ & $[5;10)$ & $[10;15)$ & $[15;20)$ & $[20;25)$ \\ 
				\hline 
				$\textbf{Giá trị  đại diện}$ & $2{,}5$ & $7{,}5$ & $12{,}5$ & $17{,}5$ & $22{,}5$ \\
				 \hline
				$\textbf{Số học sinh}$ & $8$ & $16$ & $4$ & $2$ & $2$ \\ 
				\hline 
			\end{tabular}
		\end{center}
		Thời gian xem ti vi trung bình trong tuần trước của các bạn  học sinh này là
		\begin{align*}
			\overline{x}=\dfrac{8\cdot 2{,}5+16\cdot 7{,}5+4\cdot 12{,}5+2\cdot 17{,}5+2\cdot 22{,}5}{8+16+4+2+2}=8{,}4375 (\text{giờ}). 
		\end{align*}
	 }
\end{ex}

% nguồn: Ok_TuanLuan_Toan11_Lan2
\begin{ex}%[1D5H1-3]%[1D5H1-4]%[ma:f0d1014d]
	Thảo thống kê lại số bước chân bạn đi mỗi ngày trong $3$ tháng. Kết quả được biểu diễn  ở biểu đồ ở bên.
	\begin{center}
		\begin{tikzpicture}[line cap=round,line join=round,font=\footnotesize,>=stealth,scale=1.3]
				\draw[->] (0,0)--(0,4) node[left] {(số ngày)}; 
				\draw[->]  (0,0)--(9,0) node[below] {(số bước chân)};
				 \draw[black!30!white] (9,1) -- (0,1) node[left] {$10$};
				\draw[black!30!white] (9,2) -- (0,2)  node[left] {$20$};
				\draw[black!30!white] (9,3) -- (0,3) node[left] {$30$}; 
				\draw[fill=red!40!white] (0,0) rectangle  (1.5,1.2);
				\draw[fill=red!40!white]  (1.5,0) rectangle (3,2.4);
				 \draw[fill=red!40!white] (3,0) rectangle (4.5,2.9);
				\draw[fill=red!40!white] (4.5,0) rectangle (6,1.8); 
				 \draw[fill=red!40!white] (6,0) rectangle (7.5,0.9);
				\path (0.75,-1) node  {\rotatebox{45}{$[4\,000; 5\,000)$}};
				\path (0.75,1.2)  node[above] {$12$};
				\path (2.25,-1) node {\rotatebox{45}{$[5\,000; 6\,000)$}}; 
				 \path (2.25,2.4) node[above] {$24$};
				\path (3.75,-1)  node {\rotatebox{45}{$[6\,000; 7\,000)$}};
				 \path (3.75,2.9) node[above] {$29$};
				\path (5.25,-1) node {\rotatebox{45}{$[7\,000;  8\,000)$}};
				 \path (5.25,1.8) node[above] {$18$};
				\path  (6.75,-1) node {\rotatebox{45}{$[8\,000; 9\,000)$}};
				\path (6.75,0.9) node[above] {$9$}; 
			\end{tikzpicture}
	\end{center}
	\begin{listEX}
		 \item Hãy lập bảng tần số ghép nhóm, kèm theo giá trị đại diện biểu diễn dữ liệu thống kê trên.
		\item Hãy ước lượng số trung bình và mốt của  mẫu số liệu ghép nhóm trên.
	\end{listEX}
	\loigiai{ 
		\begin{listEX}
			\item Bảng tần  số ghép nhóm:
			\begin{center}
				\begin{tabular}{|c|c|c|}
					\hline 
					 Số bước chân &Giá trị đại diện& Số ngày\\
					\hline 
					 $[4\,000; 5\,000)$ &$4\,500$&$12$\\
					 \hline
					 $[5\,000; 6\,000)$ &$5\,500$&$24$\\
					\hline 
					$[6\,000; 7\,000)$ &$6\,500$&$29$\\ 
					\hline 
					$[7\,000; 8\,000)$ &$7\,500$&$18$\\ 
					\hline 
					$[8\,000; 9\,000)$  &$8\,500$&$9$\\
					\hline 
				\end{tabular}
			\end{center}
			\item Số trung bình của mẫu số liệu ghép  nhóm là
			\[\overline{x}=\dfrac{12\cdot 4\,500+24\cdot 5\,500+29\cdot 6\,500+18\cdot  7\,500+9\cdot 8\,500}{92} \approx 6\,369{,}57.\]
			Nhóm  chứa mốt của mẫu số liệu trên là nhóm $\left[6\,000;7\,000 \right)$.\\
			Do  đó $u_m=6\,000$; $n_{m-1}=24$; $n_m=29$; $n_{m+1}=18$; $u_{m+1}-u_m=7\,000-6\,000=1\,000$.\\
			Mốt của  mẫu số liệu ghép nhóm là
			\[M_o=6\,000+\dfrac{29-24}{(29-24)+(29-18)}  \cdot 1\,000=6\,312{,}5.\]
		\end{listEX}
	} 
\end{ex}

% nguồn: Ok_TuanLuan_Toan11_Lan2
\begin{ex}%[1D5H1-3]%[ma:8738c7ee]
	Kết quả khảo sát cân nặng của $25$ quả cam ở mỗi lô hàng $A$ và  $B$ được cho ở bảng sau:
	\begin{center}
		\begin{tabular}{|c|c|c|c|c|c|}
			\hline Cân nặng (g) & $[150;155)$ & $[155;160)$  & $[160;165)$ & $[165;170)$ & $[170;175)$\\
			 \hline
			Số cam ở lô hàng $A$ & $2$ & $6$ & $12$ & $4$ & $1$ \\ 
			\hline 
			Số cam ở lô hàng $B$ & $1$ & $3$ & $7$ & $10$  & $4$\\
			 \hline
		\end{tabular}
	\end{center}
	\begin{enumerate}[a)]
		\item Hãy ước lượng cân  nặng trung bình của mỗi quả cam ở lô hàng $A$ và lô hàng $B$.
		\item Nếu so sánh theo số trung bình thì  cam ở lô hàng nào nặng hơn?
	\end{enumerate}
	\loigiai{ 
		\begin{enumerate}[a)]
			\item Ta  có bảng thống kê số lượng cam theo giá trị đại diện:
			\begin{center}
				\begin{tabular}{|c|c|c|c|c|c|}
					\hline Cân nặng (g)  & $[150;155)$ & $[155;160)$ & $[160;165)$ & $[165;170)$ & $[170;175)$\\
					 \hline
					Cân nặng đại diện (g) & $152{,}5$ & $157{,}5$ & $162{,}5$ & $167{,}5$ & $172{,}5$  \\
					 \hline
					Số cam ở lô hàng $A$ & $2$ & $6$ & $12$ &  $4$ & $1$ \\
					 \hline
					 Số cam ở lô hàng $B$ & $1$ & $3$ & $7$ & $10$ & $4$\\
					\hline 
				\end{tabular}
			\end{center}
			Cân nặng trung bình của mỗi quả cam ở lô hàng  $A$ xấp xỉ bằng
			\[\dfrac{2\cdot  152{,}5+6\cdot157{,}5+12\cdot162{,}5+4\cdot 167{,}5+1\cdot 172{,}5}{25}=161{,}7 \,(\mathrm{g}).\;
			 \]
			Cân nặng trung bình  của mỗi quả cam ở lô hàng $B$ xấp xỉ bằng
			 \[\dfrac{1\cdot 152{,}5+3\cdot 157{,}5+7\cdot 162{,}5+10\cdot 167{,}5+4\cdot 172{,}5}{25}=165{,}1\,(\mathrm{g}).
			 \]
			\item Nếu so sánh theo số trung bình thì  cam ở lô hàng $B$ nặng hơn cam ở lô hàng $A$.
		\end{enumerate}
	 }
\end{ex}

% nguồn: Ok_TuanLuan_Toan11_Lan2
\begin{ex}%[1D5H1-3]%[ma:4a874aff]
	Nồng độ cồn trong hơi thở (đơn vị tính là miligam/$1$ lít khí thở) của $20$ lái xe  ô tô vi phạm được cho như sau:
	\begin{center}
		\begin{tabular}{*{10}{c}}
			 0{,}09&0{,}18&0{,}47&1{,}20&0{,}28&0{,}45&0{,}72&0{,}15&0{,}75&0{,}36 \\
			 0{,}21&0{,}15&0{,}23&0{,}30&0{,}41&0{,}13&0{,}05&0{,}38&0{,}42&0{,}79 . \\
		\end{tabular}
	\end{center}
	 Theo quy định, mức phạt nồng độ cồn đối với lái xe ô tô như sau:
	\begin{description}
		\item[Mức  1.] Nồng độ cồn trong hơi thở chưa vượt quá $0{,}25$ phạt từ $6$ đến $8$ triệu đồng;
		\item[Mức 2.] Nồng độ cồn trong hơi thở từ trên $0{,}25$ đến $0{,}4$ phạt từ  $16$ đến $18$ triệu đồng;
		\item[Mức 3.] Nồng độ cồn trong hơi thở vượt quá $0{,}4$ phạt từ $30$  đến $40$ triệu đồng.
	\end{description}
	\begin{listEX}
		\item Lập bảng thống kê biểu diễn số lượng lái xe vi phạm  theo mức tiền bị phạt.
		\item Trung bình mỗi lái xe bị phạt bao nhiêu tiền? Tổng số tiền phạt của $20$ lái xe khoảng  bao nhiêu?
	\end{listEX}
	\loigiai{ 
		\begin{listEX}
			\item Bảng  thống kê:
			\begin{center}
				\begin{tabular}{|*4{c|}}
					\hline 
					Số tiền phạt (triệu đồng) & $6-8$ & $16-18$  & $30-40$ \\
					\hline 
					Số người vi phạm & 8 &  4 & 8 \\
					 \hline
				\end{tabular}
			\end{center}
			\item  Cỡ mẫu $n=20$. Số tiền trung bình một người bị phạt là $\overline{x}=\dfrac{8\cdot 7+4\cdot 17+8\cdot 35}{20}=20{,}2$(triệu đồng).
			Tổng số tiền 20  lái xe bị phạt là $20{,}2 \cdot 20=404$ (triệu đồng).
		\end{listEX}
	} 
\end{ex}

% nguồn: Ok_TuanLuan_Toan11_Lan2
\begin{ex}%[1D5H1-3]%[1D5H1-4]%[ma:aadc4945]
	Bảng sau thống kê cân  nặng (đơn vị: kg) của một số con ngan đực 88 ngày tuổi ở một trang trại.
	\begin{center}
		\begin{tabular}{|*{10}{c|}}
			\hline 
			4{,}60 & 4{,}62 & 4{,}64 & 4{,}65 & 4{,}67 &  4{,}67 & 4{,}68 & 4{,}68 & 4{,}70 & 4{,}70 \\
			 \hline
			4{,}70 & 4{,}70 & 4{,}71 & 4{,}71 & 4{,}72 & 4{,}73 & 4{,}74 & 4{,}76 &  4{,}77 & 4{,}77 \\
			\hline 
			4{,}77 & 4{,}78 & 4{,}78 & 4{,}80 & 4{,}82 & 4{,}84  & 4{,}84 & 4{,}85 & 4{,}87 & 4{,}89 \\
			 \hline
			4{,}89 &  4{,}90 & 4{,}92 & 4{,}92 & 4{,}93 & 4{,}94 & 4{,}94 & 4{,}95 & 4{,}97 & 4{,}97 \\
			 \hline
			4{,}97 & 4{,}99 & 4{,}99 & 5{,}01 & 5{,}02 & 5{,}03 & 5{,}04 & 5{,}05 &  5{,}06 & 5{,}07 \\
			 \hline
		\end{tabular}
	\end{center}
	\begin{listEX}
		\item Hãy lập bảng  tần số ghép nhóm của mẫu số liệu trên với nhóm đầu tiên là $[4{,}6; 4{,}7)$.
		 \item Hãy ước lượng số trung bình và mốt của mẫu số liệu ghép nhóm trên.
	\end{listEX}
	\loigiai{ 
		\begin{listEX}
			\item Bảng tần  số ghép nhóm của mẫu số liệu đã cho như sau:
			\begin{center}
				\begin{tabular}{|*{6}{c|}}
					\hline 
					Cân  nặng (kg) & $[4{,}6; 4{,}7)$& $[4{,}7; 4{,}8)$& $[4{,}8; 4{,}9)$& $[4{,}9; 5{,}0)$& $[5{,}0; 5{,}1)$\\
					\hline 
					Số con ngan & 8 & 15 & 8  & 12 & 7 \\
					 \hline
				\end{tabular}
			\end{center}
			\item  Số trung bình của mẫu số liệu ghép nhóm là $4{,}84$. Mốt của mẫu số liệu ghép nhóm là $4{,}75$.
		\end{listEX}
	 }
\end{ex}

% nguồn: Ok_TuanLuan_Toan11_Lan2
\begin{ex}%[1D5H1-3]%[ma:a21ca2c2]
	Biểu  đồ bên dưới thống kê chiều cao (đơn vị: cm) của các vận động viên của hai đội bóng rổ Sao La và Kim Ngưu.
	\begin{center}
		 \resizebox{0.8\linewidth}{!}{
		\begin{tikzpicture}[>=stealth,declare function={ox=-2; kc=3; rc=1.0; kcy=0.5; }]
				\foreach[count=\i] \yt in  {2,4,6,8,10,12}{
					\draw[gray!35] (0,{\i*2*kcy})  node[left,text=black]{$\yt$}--+(17,0);
				} 
				\foreach[count=\i] \yt/\yh/\xt in {2/2/{[170; 175)},4/3/{[175;  180)},5/4/{[180; 185)},5/10/{[185; 190)},4/1/{[190; 195)}}{
					\fill[blue!75!black]  ({ox+\i*kc},0) rectangle +(rc,{\yt*kcy});
					\fill[orange] ({ox+\i*kc+rc+0.25},0)  rectangle +(rc,{\yh*kcy});
					\path  ({ox+\i*kc+rc},0) node[below]{\xt};
				} 
				\draw[->]  (0,0)--(17,0) node[anchor=west,shift={(-15pt,-6pt)}]{cm};
				\draw[->] (0,0)--(0,6.75) node[above]{Số vận động  viên};
				\path  (0,0) node[left]{$0$};
				\path (current bounding box.north) node[below,font=\bfseries\sffamily\large]{Chiều cao  của các vận động viên};
			\end{tikzpicture}
		 }
	\end{center}
	Hãy so sánh chiều cao của các vận động viên hai đội bóng theo số trung bình và trung  vị.
	\loigiai{ 
		Ta có bảng tần  số ghép nhóm như sau
		\begin{longtable}{|*{6}{c|}}
			 \hline
			Chiều cao (cm) & [$170$; $175$) & [$175$; $180$)  & [$180$; $185$) & [$185$; $190$) & [$190$; $195$) \\
			\hline 
			Chiều cao đại diện & $172{,}5$ &  $177{,}5$ & $182{,}5$ & $187{,}5$ & $192{,}5$ \\
			\hline 
			Đội Sao La & $2$ & $4$  & $5$ & $5$ & $4$ \\
			 \hline
			Đội Kim Ngưu & $2$ & $3$ & $4$ & $10$ &  $1$ \\
			\hline 
		\end{longtable}
		 \noindent
		Theo số  trung bình\\
		Chiều cao trung bình của  các vận động viên đội Sao La xấp xỉ bằng\\
		 $\overline{x_1}=\dfrac{2\cdot 172{,}5+4\cdot 177{,}5+5\cdot 182{,}5+5\cdot 187{,}5+4\cdot 192{,}5}{2+4+5+5+4} \approx 183{,}75$ (cm).\\
		Chiều cao trung bình của các  vận động viên đội Kim Ngưu xấp xỉ bằng\\
		$\overline{x_2}=\dfrac{2\cdot 172{,}5+3\cdot 177{,}5+4\cdot 182{,}5+10\cdot 187{,}5+1\cdot 192{,}5}{2+3+4+10+1}  \approx 183{,}75$ (cm).\\
		 Vậy theo chiều cao trung bình cả hai đội có chiều cao xấp xỉ như nhau.\\
		 Theo trung vị\\
		Gọi $x_1, x_2, \ldots,  x_{20}$ là chiều cao của 20 thành viên đội Sao La xếp theo thứ tự không giảm.\\
		Số trung vị của mẫu  số liệu trên là $\dfrac{1}{2} \left(x_{10}+x_{11} \right)$.\\
		Ta có $x_1, x_2 \in \left[170; 175\right)$; $x_3, x_4, x_5, x_6 \in \left[175; 180\right)$; $x_7, x_8, x_9, x_{10}, x_{11} \in \left[180;  185\right)$.\\
		Nhận thấy $\dfrac{1}{2} \left(x_{10}+x_{11} \right)$  thuộc $\left[180; 185\right)$.\\
		Số trung vị của  mẫu số liệu đội Sao La là: $M_o=180+\dfrac{\dfrac{20}{2}-\left(2+4\right)}{5} \left(185-180\right)=184$.\\
		Gọi $y_1, y_2, \ldots, y_{20}$ là chiều cao của 20 thành viên  đội Kim Ngưu xếp theo thứ tự không giảm.\\
		Số trung vị của mẫu số liệu trên  là $\dfrac{1}{2} \left(y_{10}+y_{11} \right)$.\\
		Ta có $y_1, y_2 \in \left[170;  175\right)$; $y_3, y_4, y_5 \in \left[175; 180\right)$; $y_6, y_7, y_8, x_9 \in \left[180; 185\right)$;
		 $y_{10}, y_{11}, \ldots, y_{19} \in \left[185; 190\right)$;
		$y_{20} \in  \left[190; 195\right)$.\\
		Nhận thấy $\dfrac{1}{2} \left(y_{10}+y_{11} \right)$  thuộc $\left[185; 190\right)$.\\
		 Số trung vị của mẫu số liệu đội Kim Ngưu là: $M_o ^{{'}}=185+\dfrac{\dfrac{20}{2}-\left(2+3+4\right)}{10} \left(190-185\right)=185{,}5$.\\
		Vậy theo trung vị, chiều  cao của đội Kim Ngưu cao hơn so với đội Sao La.
	} 
\end{ex}

% nguồn: Ok_TuanLuan_Toan11_Lan2
\begin{ex}%[1D5H1-3]%[ma:fede9c4b]
	Cân nặng của $28$ học sinh nam lớp $11$ được cho  như sau:
	\[\begin{tabular}{lllllll}
	{55{,}4} & {62{,}6} & {54{,}2} & {56{,}8} & {58{,}8} & {59{,}4} & {60{,}7}\\ 
	{58} & {59{,}5} & {63{,}6} & {61{,}8} & {52{,}3} & {63{,}4} &  {57{,}9}\\
	{49{,}7} & {45{,}1} &  {56{,}2} & {63{,}2} & {46{,}1} & {49{,}6} & {59{,}1}\\
	{55{,}3} & {55{,}8} & {45{,}5} & {46{,}8} & {54} & {49{,}2}  & {52{,}6}.
\end{tabular}\]
\begin{enumerate}[a)]
	\item Hãy chia mẫu dữ liệu trên thành $5$ nhóm, lập bảng tần số ghép nhóm và  xác định giá trị đại diện cho mỗi nhóm.
	\item Hãy ước lượng cân nặng trung bình của học sinh lớp $11$. 
\end{enumerate}
 \loigiai{
	\begin{enumerate}[a)]
		\item Khoảng biến thiên của  mẫu số liệu trên là $R=63{,}6-45{,}1=18{,}5$.\\
		Độ dài mỗi nhóm $L> \dfrac{R}{k}=\dfrac{18{,}5}{5}=3{,}7$.\\ 
		Ta chọn $L=4$ và chia dữ liệu thành các  nhóm $\left[45;49\right)$, $\left[49;53\right)$, $\left[53;57\right)$, $\left[57;61\right)$, $\left[61;65\right)$.\\
		Khi đó ta có bảng tần số  ghép nhóm sau:
		\begin{center}
			\begin{tabular}{|c|c|c|c|c|c|}
				\hline Cân nặng (g) & $[45;49)$ &  $[49;53)$ & $[53;57)$ & $[57;61)$ & $[61;65)$\\
				\hline 
				Giá trị đại diện & $47$ &  $51$ & $55$ & $59$ & $63$ \\
				\hline 
				Số học sinh & $4$  & $5$ & $7$ & $7$ & $5$\\
				\hline 
			\end{tabular}
		\end{center}
		\item  Cân nặng trung bình của học sinh trong lớp $11$ xấp xỉ là
		\[\dfrac{47\cdot 4+51\cdot 5+55\cdot 7+59\cdot 7+63\cdot 5}{28}=55{,}6\,  (\mathrm{kg}).
		\] 
	\end{enumerate}
} 
\end{ex}

% nguồn: Ok_TuanLuan_Toan11_Lan1
\begin{bt}%[1D5V1-3]%[ma:5918a377]
	Để đảm bảo bảng số liệu được phân bố  đồng đều người ta điều chỉnh các giá trị $x$, $y$ sao cho số trung bình cộng và số trung vị bằng nhau. Khi đó bảng số liệu được cho như sau
	\begin{center}
		\begin{tabular}{|c|c|c|c|c|c|}
			\hline Giá trị&$40$&$50$&$60+x$&$90+y$&\\ 
			\hline Tần  số&$20$&$30$&$60-x$&$90-y$&$200$\\
			 \hline
		\end{tabular}
	\end{center}
	 Biết rằng $x\le 0$. Tìm $x$?
	 \loigiai{
		Số trung  bình của bảng số liệu bằng $$\overline{x}=\dfrac{(60+x)(60-x)+(90-y)(90+y)+2300}{200}.$$
		Số trung vị là ${M_e}=\dfrac{30+60-x}{2}=\dfrac{90-x}{2}$.\\ 
		Ta có $20+30+60-x+90-y=200\Leftrightarrow x+y=0$.\\ 
		Theo đề bài  ta có phương trình
		$$\dfrac{(60+x)(60-x)+(90-y)(90+y)+2300}{200}=\dfrac{90-x}{2}.$$ 
		Giải trình ta được nghiệm $x=25-25\sqrt {5}$ hoặc $x=25+5\sqrt {21}$, vì $x\le 0$ nên nhận $x=25-25\sqrt {5}$, số nguyên gần với nó nhất  là $-31$.
	 }
\end{bt}

% nguồn: Ok_TuanLuan_Toan11_Lan2
\begin{ex}%[1D5V1-3]%[ma:7f9807b8]
	Để đảm bảo bảng số liệu được phân bố đồng đều người ta điều chỉnh các giá trị $x$, $y$ sao cho số  trung bình cộng và số trung vị bằng nhau. Khi đó bảng số liệu được cho như sau
	\begin{center}
		\begin{tabular}{|c|c|c|c|c|c|}
			\hline Giá trị&$40$&$50$&$60+x$&$90+y$&\\ 
			\hline  Tần số&$20$&$30$&$60-x$&$90-y$&$200$\\
			 \hline
		\end{tabular}
	\end{center}
	 Biết rằng $x\le 0$. Tìm $x$?
	\loigiai{ 
		Số trung bình của bảng số liệu bằng  $$\overline{x}=\dfrac{(60+x)(60-x)+(90-y)(90+y)+2300}{200}.$$
		Số trung vị  là ${M_e}=\dfrac{30+60-x}{2}=\dfrac{90-x}{2}$.\\
		Ta có $20+30+60-x+90-y=200\Leftrightarrow x+y=0$.\\ 
		Theo đề  bài ta có phương trình
		$$\dfrac{(60+x)(60-x)+(90-y)(90+y)+2300}{200}=\dfrac{90-x}{2}.$$ 
		Giải trình ta được nghiệm $x=25-25\sqrt {5}$ hoặc $x=25+5\sqrt {21}$, vì $x\le 0$ nên nhận $x=25-25\sqrt {5}$, số nguyên gần  với nó nhất là $-31$.
	 }
\end{ex}
